% concatenated from AAVII-combine-2026-0216arx.tex, lstyles.tex
% This is a general template file for the LaTeX package SVJour3f
% for Springer journals.          Springer Heidelberg 2010/09/16
%
% Copy it to a new file with a new name and use it as the basis
% for your article. Delete % signs as needed.
%
% This template includes a few options for different layouts and\red
% content for various journals. Please consult a previous issue of
% your journal as needed.
%
%%%%%%%%%%%%%%%%%%%%%%%%%%%%%%%%%%%%%%%%%%%%%%%%%%%%%%%%%%%%%%%%%%%
%
% First comes an example EPS file -- just ignore it and
% proceed on the \documentclass line
% your LaTeX will extract the file if required
%\begin{filecontents*}{example.eps}
%	%!PS-Adobe-3.0 EPSF-3.0
%	%%BoundingBox: 19 19 221 221
%	%%CreationDate: Mon Sep 29 1997
%	%%Creator: programmed by hand (JK)
%	%%EndComments
%	gsave
%	newpath
%	20 20 moveto
%	20 220 lineto
%	220 220 lineto
%	220 20 lineto
%	closepath
%	2 setlinewidth
%	gsave
%	.4 setgray fill
%	grestore
%	stroke
%	grestore
%\end{filecontents*}
%%
%\RequirePackage{fix-cm}
%
%\documentclass{svjour3}                     % onecolumn (standard format)
%%\documentclass[smallcondensed]{svjour3}     % onecolumn (ditto)
%\documentclass[smallextended]{svjour3}       % onecolumn (second format)
%\documentclass[twocolumn]{svjour3}          % twocolumn
%
%\smartqed  % flush right qed marks, e.g. at end of proof
%

%\documentclass[12pt]{amsart} %a4
%\usepackage{graphicx}
%%
%\usepackage{mathptmx}      % use Times fonts if available on your TeX system
%%
%% insert here the call for the packages your document requires
%\usepackage{latexsym}
%% etc.
%%
%% please place your own definitions here and don't use \def but
%% \newcommand{}{}
%%
%% Insert the name of "your journal" with
%% \journalname{myjournal}
%%
%
%\pdfoutput=1 

%\NeedsTeXFormat{LaTeX2e}% LaTeX 2.09 can't be used (nor non-LaTeX)
%[1994/12/01]% LaTeX date must December 1994 or later
%\documentclass{ijmart}
%    Some definitions useful in producing this sort of documentation:
%\chardef\bslash=`\\ % p. 424, TeXbook
%%    Normalized (nonbold, nonitalic) tt font, to avoid font
%%    substitution warning messages if tt is used inside section
%%    headings and other places where odd font combinations might
%%    result.
%\newcommand{\ntt}{\normalfont\ttfamily}
%%    command name
%\newcommand{\cn}[1]{{\protect\ntt\bslash#1}}
%%    LaTeX package name
%\newcommand{\pkg}[1]{{\protect\ntt#1}}
%%    File name
%\newcommand{\fn}[1]{{\protect\ntt#1}}
%%    environment name
%\newcommand{\env}[1]{{\protect\ntt#1}}
%\hfuzz1pc % Don't bother to report overfull boxes if overage is < 1pc

\documentclass{amsart}

\usepackage{amscd,amssymb,dsfont} %epsf
\usepackage{url}
%\usepackage{graphics}
\usepackage{rotating}
\usepackage[utf8]{inputenc}
\usepackage{amsmath}
%\usepackage{graphicx}

%\usepackage[backref,bookmarks=true,colorlinks=true,pdfstartview=FitV,linkcolor=blue, citecolor=red,urlcolor=green]{hyperref}  % hyperref package

\usepackage[colorlinks=true]{hyperref}
\usepackage{mathrsfs}
%\usepackage{natbib}
%\usepackage{hyperref]

%\usepackage{hyperref}
%\usepackage{academicons}
%\usepackage{xcolor}
%
%\newcommand{\orcid}[1]{\href{https://orcid.org/#1}{\textcolor[HTML]{A6CE39}{\aiOrcid}}}

%\usepackage{graphicx}
%\usepackage{amssymb,amsmath}
%\usepackage{amsthm}
%\usepackage{epstopdf}
%\DeclareGraphicsRule{.tif}{png}{.png}{`convert #1 `dirname #1`/`basename #1 .tif`.png}
%\usepackage{sfmath}
\usepackage{bbm}
\usepackage{framed}

% \usepackage[pagewise, modulo]{lineno}
% \linenumbers

\usepackage{xcolor}

\newcommand{\red}[1]{\textcolor{black}{#1}} 
\newcommand{\reD}[1]{\textcolor{black}{#1}} 
 \newcommand{\green}[1]{\textcolor{black}{#1}}
 \newcommand{\purple}[1]{\textcolor{black}{#1}}
\newcommand{\blue}[1]{\textcolor{blue}{#1}} 
\newcommand{\newch}[1]{\textcolor{black}{#1}} 
 \newcommand{\oldch}[1]{\textcolor{black}{#1}} 

\newcommand{\dN}{d_{\tilde N}}


 

\usepackage{stmaryrd}
\usepackage{url}
%\usepackage{mathabx}
%\usepackage{lpic}
%\usepackage{pict2e}
%\usepackage{tikz}
\usepackage{pinlabel}


\usepackage{subfigure}

\usepackage[export]{adjustbox}
\usepackage{tikz-cd}


%\usepackage[pagewise]{lineno}\linenumbers

%How to find this unicode character U+202F in your LaTeX documents?
\makeatletter
\def\UTFviii@defined#1{%
	\ifx#1\relax
	!!FIXME!!%
	\else
	\expandafter#1%
	\fi
}
\makeatother

\DeclareFontFamily{OT1}{pzc}{}
\DeclareFontShape{OT1}{pzc}{m}{it}{<-> s * [1.10] pzcmi7t}{}
\DeclareMathAlphabet{\mathpzc}{OT1}{pzc}{m}{it}

\usepackage{syntonly}
%\syntaxonly

%\spnewtheorem*{xproof}{}{\itshape}{\rmfamily}% the label is assigned later
%\newcommand\xprooftitle{}
%\renewenvironment{proof}[1][\proofname]
%{\renewcommand\xproofname{#1}\xproof}
%{\endxproof}

%\usepackage[usenames,dvipsnames,svgnames,table]{xcolor}
%\newcommand*{\commt}[1]{\textcolor{DarkRed}{#1}}
%\newcommand*{\yes}[1]{\textcolor{YellowGreen}{#1}}
%\newcommand*{\hmm}[1]{\textcolor{MidnightBlue}{#1}}

\newcommand{\deltxt}[1]{}

%\setlength{\textwidth}{15.5cm}
%\setlength{\textheight}{22cm}
%\setlength{\marginparwidth}{1.5cm}
%\setlength{\voffset}{-1.0cm}
%\setlength{\hoffset}{-1.4cm}
%\setlength{\baselineskip}{14pt}

%\pagestyle{plain}
%
%\newtheorem{theorem}{Theorem}[section]
%\newtheorem{lemma}[theorem]{Lemma}
%\newtheorem{corollary}[theorem]{Corollary}
%\newtheorem{proposition}[theorem]{Proposition}
%\newtheorem{claim}[theorem]{Claim}
%\newtheorem{conj}[theorem]{Conjecture}

%\newtheorem{thmA}{Theorem A}
%\renewcommand{\thethmA}{}
%\newtheorem{thmB}{Theorem B}
%\renewcommand{\thethmB}{}
%\newtheorem{thmC}{Theorem C}
%\renewcommand{\thethmC}{}
%\newtheorem{thmD}{Theorem D}
%\renewcommand{\thethmD}{}
%\newtheorem{corA}{Corollary A}
%\renewcommand{\thecorA}{}
%\newtheorem{corB}{Corollary B}
%\renewcommand{\thecorB}{}
%\newtheorem{corC}{Corollary C}
%\renewcommand{\thecorC}{}
%\newtheorem{corD}{Corollary D}
%\renewcommand{\thecorD}{}

%\theoremstyle{remark}
%\newtheorem{remark}[theorem]{Remark}
%\newtheorem{hypothesis}[theorem]{Hypothesis}
%\newtheorem{example}[theorem]{Example}
%\theoremstyle{definition}
%\newtheorem{definition}[theorem]{Definition}

% TOP MATTER
% Theorem-like environments are defined with spnewtheorem
%
% Usage:
%
%     spnewtheorem{env_nam}{caption}[within]{cap_font}{body_font}
% or  spnewtheorem{env_nam}[numbered_like]{caption}{cap_font}{body_font}
% or  spnewtheorem*{env_nam}{caption}{cap_font}{body_font}

%\spnewtheorem{theorem}{Theorem}[section]{\bf}{\it}
%\spnewtheorem{thm}{Theorem}{\bfseries}{\itshape}
%\spnewtheorem{corollary}[theorem]{Corollary}{\bfseries}{\itshape}
%\spnewtheorem{lemma}[theorem]{Lemma}{\bfseries}{\itshape}
%\spnewtheorem{ntn}{Notations}[section]{\bfseries}{\rm}
%\spnewtheorem{proposition}{Proposition}[section]{\bf}{\it}
%\spnewtheorem{pro}[theorem]{Proposition}{\bfseries}{\itshape}
%\spnewtheorem{definition}{Definition}{\bfseries}{\rm}% maybe upshape?
%\spnewtheorem{hypothesis}{Hypothesis}{\bfseries}{\rm}
%\spnewtheorem{remark}{Remark}{\bfseries}{\rm}
%\spnewtheorem{ob}{Observation}[section]{\bfseries}{\rm}
%%\newproof{proof}{Proof}

%\begin{frontmatter}

\newtheorem{theorem}{Theorem}[section]
\newtheorem{corollary}[theorem]{Corollary}
\newtheorem{lemma}[theorem]{Lemma}
\newtheorem{proposition}[theorem]{Proposition}
%\newtheorem{axiom}{Axiom}

\theoremstyle{definition}
\newtheorem{definition}{Definition}[section]
\newtheorem{remark}{Remark}[section]
\newtheorem*{notation}{Notation}

\theoremstyle{remark}
\newtheorem{step}{Step}

%\numberwithin{equation}{section}

\newcommand{\thmref}[1]{Theorem~\ref{#1}}
\newcommand{\secref}[1]{\S\ref{#1}}
\newcommand{\lemref}[1]{Lemma~\ref{#1}}


%       Math definitions

\newcommand{\A}{\mathcal{A}}
\newcommand{\B}{\mathcal{B}}
\newcommand{\st}{\sigma}
\newcommand{\XcY}{{(X,Y)}}
\newcommand{\SX}{{S_X}}
\newcommand{\SY}{{S_Y}}
\newcommand{\SXY}{{S_{X,Y}}}
\newcommand{\SXgYy}{{S_{X|Y}(y)}}
\newcommand{\Cw}[1]{{\hat C_#1(X|Y)}}
\newcommand{\G}{{G(X|Y)}}
\newcommand{\PY}{{P_{\mathcal{Y}}}}
\newcommand{\X}{\mathcal{X}}
\newcommand{\wt}{\widetilde}
\newcommand{\wh}{\widehat}

\DeclareMathOperator{\per}{per}
\DeclareMathOperator{\cov}{cov}
\DeclareMathOperator{\non}{non}
\DeclareMathOperator{\cf}{cf}
\DeclareMathOperator{\add}{add}
\DeclareMathOperator{\Cham}{Cham}
\DeclareMathOperator{\IM}{Im}
\DeclareMathOperator{\esssup}{ess\,sup}
\DeclareMathOperator{\meas}{meas}
\DeclareMathOperator{\seg}{seg}

%    \interval is used to provide better spacing after a [ that
%    is used as a closing delimiter.
\newcommand{\interval}[1]{\mathinner{#1}}

%    Notation for an expression evaluated at a particular condition. The
%    optional argument can be used to override automatic sizing of the
%    right vert bar, e.g. \eval[\biggr]{...}_{...}
\newcommand{\eval}[2][\right]{\relax
	\ifx#1\right\relax \left.\fi#2#1\rvert}

%    Enclose the argument in vert-bar delimiters:
\newcommand{\envert}[1]{\left\lvert#1\right\rvert}
\let\abs=\envert

%    Enclose the argument in double-vert-bar delimiters:
\newcommand{\enVert}[1]{\left\lVert#1\right\rVert}
\let\norm=\enVert


\newcommand\bR{{\mathbb{R}}}
\newcommand\bC{{\mathbb C}}
\newcommand\bZ{{\mathbb Z}}
\newcommand\bQ{{\mathbb Q}}
\newcommand\bN{{\mathbb N}}
\newcommand\bT{{\mathbb T}}
\newcommand\bV{{\mathbb V}}
\newcommand\bE{{\mathbb E}}
\newcommand\bH{{\mathbb{H}}}
\newcommand\bA{{\mathbb{A}}}
\newcommand\bU{{\mathbb{U}}}
\newcommand\Hom{{\rm Hom}}
\newcommand\dev{{\bf dev}}
\newcommand\SI{{\mathbb{S}}}
\newcommand\ima{{\rm Im}}
\newcommand\inte{{\rm int}}
\newcommand\Bd{{\rm bd}}
\newcommand\clo{{\rm Cl}}
\newcommand\bdd{{\mathbf{d}}}
\newcommand\ideal[1]{\tilde #1_\infty}
\newcommand\ra{\rightarrow}
\newcommand\longra{\longrightarrow}
\newcommand\che{\check}
\newcommand\emp{\emptyset}
\newcommand\eps{\epsilon}
\newcommand\disk{{\cal D}}
\newcommand\disc{{\rm disc}\,}
\newcommand\vth{\vartheta}
\newcommand\vpi{\varphi}
\newcommand\Aff{{\mathbf{Aff}}}
\newcommand\ovl{\overline}
\newcommand\Aut{{\mathbf{Aut}}}
\newcommand\Idd{{\rm I}}
\newcommand\Rar{\Rightarrow}
\newcommand\uar{\uparrow}
\newcommand\Out{{\mathbf{Out}}}
\newcommand\Inn{{\mathbf{Inn}}}
\newcommand\calR{{\mathcal R}}
\newcommand\calS{{\mathcal S}}
\newcommand\CP{{\mathbb{CP}}}
\newcommand\tr{{\mathrm{tr}}}
\newcommand\ssl{{\mathsf{l}}}
\newcommand\bv{{\mathbf{v}}}
\newcommand\bu{{\mathbf{u}}}
\newcommand\CN{{\mathcal{N}}}
\newcommand\Isom{{\mathbf{Isom}}}

%Text
\newcommand\awlg{may assume without loss of generality }
\newcommand\cO{{\rm ,}\ }
\newcommand\hll{\hfill\break}
\newcommand\pO{{\rm .}}
\newcommand\tgt{\noindent}


%\newcommand\twotwo{$(\hbox{{\rm ; }} 2,2)$-disk }
%\newcommand\twotwos{$(\hbox{{\rm ; }} 2,2)$-disks }
\newcommand\smc{\hbox{\rm ; }}
\newcommand\tri{\triangle}
\newcommand\rpt{\mathbb{RP}^2}
\newcommand\rptd{\mathbb{RP}^{2\dagger}}
\newcommand\rpn{\mathbb{RP}^n}
\newcommand\rpts{\rpt{}^*}
\newcommand\rpo{\mathbb{RP}^1}
\newcommand\rpno{\mathbb{RP}^{n-1}}
\newcommand\RP{\mathbb{RP}}
%\newcommand\Idd{\mathrm{I}}
\newcommand\SOThr{{\mathsf{SO}}(3)}
\newcommand\SOn{{\mathsf{SO}}(n)}
\newcommand\Oon{{\mathsf{O}}(1, n)}
\newcommand\Pon{{\mathsf{PO}}(1, 2)}
\newcommand\SUTw{{\mathsf{SU}}(2)}
\newcommand\rep{\mathrm{rep}}
\newcommand\SLThr{{\mathsf{SL}}(3, \bR)}
\newcommand\PGLThr{{\mathsf{PGL}}(3, \bR)}
\newcommand\SL{{\mathsf{SL}}}
\newcommand\PSL{{\mathsf{PSL}}}
\newcommand\SO{{\mathsf{SO}}}
\newcommand\Ort{{\mathsf{O}}}
\newcommand\PGL{{\mathsf{PGL}}}
\newcommand\SLnp{{\mathsf{SL}}_\pm(n+1, \bR)}
\newcommand\SLn{{\mathsf{SL}}_\pm(n, \bR)}
\newcommand\SLf{{\mathsf{SL}}_\pm(4, \bR)}
\newcommand\SLt{{\mathsf{SL}}_\pm(3, \bR)}
\newcommand\GL{{\mathsf{GL}}}
\newcommand\GLnp{{\mathsf{GL}}(n+1, \bR)}
\newcommand\PGLnp{{\mathsf{PGL}}(n+1, \bR)}
\newcommand\PGLn{{\mathsf{PGL}}(n, \bR)}
\newcommand\Def{\mathrm{Def}}
\newcommand\CDef{\mathrm{CDef}}
\newcommand\TCDef{{\mathrm{TCDef}}}
\newcommand\SDef{\mathrm{SDef}}
\newcommand\hol{\mathrm{hol}}
\newcommand\orb{\mathcal{O}} 
\newcommand\torb{\tilde{\mathcal{O}}}
\newcommand\bGamma{{\boldsymbol \Gamma}}
\newcommand\leng{{\mathrm{length}}}
\newcommand\Ints{\mathrm{Int}}
\newcommand\sgn{{\mathrm{sgn}}}
\newcommand\SG{{\mathsf{SG}}}

\newcommand\cG{{\mathcal{G}}}
\newcommand\cg{{\mathsf{g}}}
\newcommand\logt{{\log_{\frac{1}{t}}}}
\newcommand\Logt{{{\mathrm{Log}}_{\frac{1}{t}}}}
\newcommand\trop{{\mathrm{trop}}}
\newcommand\Dis{{\mathcal{D}}}
\newcommand\ssb{{\mathsf{b}}}
\newcommand\mx{{\mathrm{max}}}
\newcommand\cwl{{\mathrm{cwl}}}
\newcommand\Ss{{\mathbb{S}}}
\newcommand\supp{{\mathrm{supp}}}






% vectors? i prefer in bold
%\newcommand{\vu}{\mathsf{u}}
\newcommand{\vv}{\mathbf{v}}
\newcommand{\vw}{\mathbf{w}}
\newcommand{\va}{\mathbf{a}}
\newcommand{\vb}{\mathbf{b}}
\newcommand{\vc}{\mathbf{c}}
\newcommand{\vn}{\mathbf{n}}
\newcommand{\vp}{\mathbf{p}}
\newcommand{\vx}{\mathbf{x}}
\newcommand{\vu}{\mathbf{u}}
\newcommand{\vy}{\mathbf{y}}
\newcommand{\be}{\mathbf{e}}
\newcommand{\bp}{\mathbf{p}}
%\newcommand{\vecv}{\mathsf{v}}   %%%
\newcommand{\CH}{\mathcal{CH}}
\newcommand{\rec}{\mathrm{rec}}
\newcommand{\tbV}{\widetilde{\mathbf{V}}}
\newcommand{\bg}{\mathsf{g}}
\newcommand{\Uu}{\mathbf{U}} 
\newcommand{\Uc}{\mathbf{UC}}

\newcommand{\llrrbracket}[1]{% \llrrbracket{..}
	\left[\mkern-2mu\left[#1\right]\mkern-2mu\right]}
\newcommand{\llrrparen}[1]{% \llrrparen{..}
	\left(\mkern-4mu\left(#1\right)\mkern-4mu\right)}
\newcommand{\llrrV}[1]{% \llrrparen{..}
	\left|\mkern-2mu\left|#1\right|\mkern-2mu\right|}

\newcommand{\Paren}[1]{% \Paren{..}
	\left(#1\right)}

\newcommand*{\Scale}[2][4]{\scalebox{#1}{$#2$}}%

%\newcommand*{\commt}[1]{\textcolor{red}{#1}}
%\newcommand*{\yes}[1]{\textcolor{green}{#1}}
%\newcommand*{\hmm}[1]{\textcolor{blue}{#1}}

%\newcommand{\deltxt}[1]{}

\newcommand\vol{{\mathsf{vol}}}
\newcommand\algvol{{\mathsf{algvol}}}




\setcounter{tocdepth}{3} % maybe set it to 2 for final version

\begin{document}
	
	\title[Partially hyperbolic flows]{
		Partially hyperbolic flows on flat vector bundles with 
an application to complete affine manifolds.} 

	\author[S. Choi]{Suhyoung Choi}
	%%\thanks{orcid number: 0000-0002-9201-8210}}  
	\thanks{This research was supported by the Basic Science Research Program through the National Research Foundation
		of Korea(NRF) funded by the Ministry of Education(2019R1A2C108454412), 
		Author orcid number: 0000-0002-9201-8210}
	\address{Department of Mathematical Sciences \\ KAIST \\Daejeon 34141, South Korea}
	
	\email{schoi@math.kaist.ac.kr}
	
	%\title{Closed complete special affine manifolds 
	%		cannot have Anosov holonomy }   
	%\subtitle{Anosov-affine structures} 
	%
	%\author{Suhyoung Choi\thanks{This research was supported by the Basic Science Research Program through the National Research Foundation
	%		of Korea(NRF) funded by the Ministry of Education(2019R1A2C108454412)}
	%	\thanks{orcid number: 0000-0002-9201-8210}}  
	%\orcid{0000-0002-9201-8210}
	
	%\institute{Department of Mathematical Science \\ KAIST \\Daejeon 34141, South Korea\\
	%\email{schoi@math.kaist.ac.kr}}
	
%	\date{Received on MONTH, YEAR}
%	\issueinfo{VOL}{NUM}{MONTH}{YEAR}
%	\doiinfo{10.1007/DOI-NUMBER}
	\date{\today}
	
	%\author[J. Danciger]{Jeffrey Danciger} 
	%
	%\address{Department of Mathematics,\\ The University of Texas at Austin, \\1 University Station C1200, \\ Austin TX 78712}
	%
	%\email{Jeffrey.Danciger@math.utexas.edu}
	%
	%\author[M. Kapovich]{Michael Kapovich}
	%
	%\address{Department of Mathematics MSB\\
	%	1 Shields Ave.,\\
	%	University of California, Davis\\
	%	Davis, CA 95616-8633, USA }
	%
	%\email{kapovich@math.ucdavis.edu} 
	%
	%\author[F. Stecker]{Florian Stecker} 
	%
	%\address{Department of Mathematics,\\ The University of Texas at Austin, %\\1 University Station C1200, \\ Austin TX 78712}
	%
	%\email{stecker@utexas.edu} 
	%
	
	%\date{\today}
	
	%\subjclass[2010]{Primary 57M50; Secondary 53A20, 53C15}
	
	%\keywords{affine manifold, Anosov representation, word-hyperbolic group, Gromov hyperbolic space }
	
	%\thanks{  }
	
	%\dedicatory{Dedicated to the memory of Todd Drumm}
	%\date{\today}
	%
	\subjclass[2010]{Primary 57M50; Secondary 22E40}
	%
	\keywords{affine manifold, Anosov representation, word-hyperbolic group, Gromov hyperbolic space, coarse geometry }
	%%\maketitle
	
	%\maketitle
	\begin{abstract}
	%	The conjecture of Auslander is that a closed complete affine manifold must have a virtually solvable fundamental group. 
	%	We present some supporting evidence. 
	%	We will show that a closed manifold with a complete affine structure %with a hyperbolic fundamental group 
	%	cannot have partially hyperbolic  linear holonomy
	%	nor can it have a $P$-Anosov linear holonomy for any parabolic subgroup 
	%	$P$ of $\GL(n, \bR)$. 
		%, provided the Zariski closure of the linear
		%holonomy is reductive. 

Let $N$ be a manifold of dimension $m$ with a flat vector bundle given by 
a representation $\rho:\pi_1(N) \ra \GL(n, \bR)$
where $\pi_1(N)$ is finitely generated. 
%affine transformation group acting on the complete affine space 
%$\mathds{A}^n$,. %and $K(\Gamma, 1)$ is realized as a finite CW-complex. 
The holonomy group 
$\rho$ is a \emph{$k$-partially hyperbolic holonomy representation} if the flat bundle pulled back over the unit tangent bundle of a sufficiently large compact submanifold of $N$ splits into expanding, neutral, and contracting subbundles along the geodesic flow, where the expanding and contracting subbundles are $k$-dimensional with $k < n/2$.
%Let $N$ be a complete affine manifold $\mathds{A}^n/\Gamma$ of dimension $n$ where  %$\Gamma$ is an affine transformation group acting on the complete affine space 
%$\mathds{A}^n$ and $K(\Gamma, 1)$ is realized as a finite CW-complex. 
% $N$ has a {\em $k$-partially hyperbolic holonomy group} if the tangent bundle pulled over the unit tangent bundle over a sufficiently large 
%	compact part 
%	splits into expanding, neutral, and contracting subbundles along the geodesic flow
%where the expanding and contracting subbundles have the dimension $k$, $k< n/2$. 
Suppose that each element of $\rho(\pi_1(N))$ has an eigenvalue of norm $1$,
or, alternatively,  \red{$\rho$ has some singular values of subexponential growth
in terms of word length. }
We show that $\rho$ is a $P$-Anosov representation for a parabolic subgroup 
$P$ of $\GL(n, \bR)$ 
if and only if $\rho$ is a partially hyperbolic representation. 
We will primarily employ techniques from representation theory.
\purple{As an application, we will show that the equivalence holds} when $N$ is a complete affine $n$-manifold
and $\rho$ is a linear part of the holonomy representation. 
This had never been done over the full general linear group. 
%
%		In part 1, we will show that the complete affine $n$-manifold has a $P$-Anosov linear holonomy group for a parabolic subgroup $P$ of $\GL(n, \bR)$
%	if and only if  it has a partially hyperbolic linear holonomy group. 
%\red{This had never been done over the full general linear group  before this paper.
%The part 1 will be where the techniques are mostly representation theoretical. }
%We show that if the holonomy group is partially hyperbolic of index $k$, then 
%$\mathrm{cd}(\Gamma) \leq n-k$. Moreover, 
% if a finitely-presented affine group $\Gamma$ acts on $\mathds{A}^n$ properly discontinuously and freely with the $k$-Anosov linear group, 
%	then $\mathrm{cd}(\Gamma) \leq n-k$.  
%	Also, there exists a compact collection of $n-k$-dimensional affine subspaces where $\Gamma$ acts on. 
	%The principal tools are Gromov's theory of hyperbolic metric spaces, coarse geometry, the spectrum theory of linear groups by Benoist, and the Anosov representation theory.  
		
		%\keywords{affine manifold, Anosov representation, word-hyperbolic group, Gromov hyperbolic space }
		%\subclass{57M50\and 53A20 \and 53C15}
	%	The conjecture of Auslander is that a closed complete affine manifold must have a virtually solvable fundamental group. 
	%	We present some supporting evidence. 
	%	We will show that a closed manifold with a complete affine structure %with a hyperbolic fundamental group 
	%	cannot have partially hyperbolic  linear holonomy
	%	nor can it have a $P$-Anosov linear holonomy for any parabolic subgroup 
	%	$P$ of $\GL(n, \bR)$. 
		%, provided the Zariski closure of the linear
		%holonomy is reductive. 
%Let $N$ be a complete affine manifold $\mathds{A}^n/\Gamma$ of dimension $n$ where  $\Gamma$ is an affine transformation group and $K(\Gamma, 1)$ is realized as a finite CW-complex. 
% $N$ has a {\em partially hyperbolic holonomy group} if the tangent bundle pulled over the unit tangent bundle over a sufficiently large 
%	compact part 
%	splits into expanding, neutral, and contracting subbundles along the geodesic flow.
	%	We will show that the complete affine $n$-manifold has a $P$-Anosov linear holonomy group for a parabolic subgroup $P$ of $\GL(n, \bR)$
	%if and only if  it has a partially hyperbolic linear holonomy group. 
%In part 2, we demonstrate that if the holonomy group is partially hyperbolic of index $k$, where $k < n/2$, then $\mathrm{cd}(\Gamma) \leq n-k$. Moreover, if a finitely-presented affine group $\Gamma$ acts properly discontinuously and freely on $\mathds{A}^n$ with a $k$-Anosov linear subgroup for $k \leq n/2$, then $\mathrm{cd}(\Gamma) \leq n-k$.
%In part 2, we show that if the holonomy group is partially hyperbolic of index $k$, $k < n/2$, then 
%$\mathrm{cd}(\Gamma) \leq n-k$. Moreover, 
% if a finitely-presented affine group $\Gamma$ acts on $\mathds{A}^n$ properly %discontinuously and freely with the $k$-Anosov linear group for $k \leq n/2$, 
%	then $\mathrm{cd}(\Gamma) \leq n-k$.  
%	This will show that if a finitely-presented affine group $\Gamma$ acts on $\mathds{A}^n$ properly discontinuously and freely with the $k$-Anosov linear group, 
%	then $\mathrm{cd}(\Gamma) \leq n-k$.  
%Also, there exists a compact collection of $n-k$-dimensional affine subspaces where $\Gamma$ acts. The techniques employed here mostly stem from the coarse geometry theory.
%Canary and Tsouvalis previously proved the same result using the powerful method of Bestvina and Mess for word hyperbolic groups; however, our approach differs in that  our method projects the holonomy cover to a stable affine subspace, and we plan to generalize to relative Anosov groups.

%	Also, there exists a compact collection of $n-k$-dimensional affine subspaces where $\Gamma$ acts on.  The techniques here are mostly from coarse geometry.
%\red{Canary and Tsouvalis proved the same result before using the powerful method of Bestvina and Mess for word hyperbolic groups; however, our techniques are  different
%in that our method is to project the holonomy cover to a stable affine subspace, and we plan to generalize to relative Anosov groups.}
	%The principal tools are Gromov's theory of hyperbolic metric spaces, coarse geometry, %the spectrum theory of linear groups by Benoist, and the Anosov representation theory.  
		
		%\keywords{affine manifold, Anosov representation, word-hyperbolic group, Gromov hyperbolic space }
		%\subclass{57M50\and 53A20 \and 53C15}
	\end{abstract}
	\maketitle
	\tableofcontents
%	\maketitle 


%\part{Hyperbolic bundles and  Anosov representations} \label{part2}  

\section{Introduction}

%We continue from our previous paper \cite{AAI} with the same notations. 
%The purpose is to drop the condition of $M$ having a negative curvature. 

\subsection{Notation}
%Let $\Aff(\mathds{A}^n)$ denote the group of affine transformations of 
%an affine space $\mathds{A}^n$
%whose elements are of the form: 
%\[  x \mapsto A x + \bv   \]
%for a vector $\bv \in \bR^{n}$ and $A \in \GL(n, \bR)$. 
%Let $\mathcal{L}: \Aff(\mathds{A}^n) \ra \GL(n, \bR)$ denote
%the map sending the affine transformations to their linear parts.

%We basically follow the definition of Guichard and Wienhard \cite{GW12}. 
%
%$M$ has a Riemannian metric and this induces metric on $T\mathds{A}^{ n}$. 
%Hence, $\Uu \mathds{A}^n$ can be identified with unit sphere bundle in $T\mathds{A}^n$. 
%Let $M$ be a manifold and $\hat M$ the $G$-cover of $M$
%and $\rho: G \ra \SL_{\pm}(n, \bR)$ denote a representation, and let $\mathcal{S}$ 
%an $\SL_{\pm}(n, \bR)$-space. We denote by 
%$\mathcal{S}_{\rho} = (\hat M \times \mathcal{S})/\Gamma$ where 
%$\Gamma$ acts diagonally. 


%Let $\Uu M$ be the unit tangent bundle over a compact negatively 
%curved Riemannian manifold $M$
%with a geodesic flow $\phi_{t}: \Uu M \ra \Uu M, t \in \bR$.  
%Denote by $\hat M$ the universal cover of $M$ with the covering map
%$p_M$. 
%There is a covering $\Uu p_M:\Uu \hat M \ra \Uu M$ from the unit tangent bundle 
%$\Uu \hat M$ of $\hat M$.  
%Let $\pi_{\Uu \hat M} : \Uu M \ra M$ denote the fibration 
%and $\pi_{\Uu \hat M} : \Uu \hat M \ra \hat M$ the induced fibration. 
%\begin{itemize} 
%\item We define 
%$\bR^n_{\rho\circ \Pi_{M\ast}}:= (\Uu \hat M \times \bR^n) /\Gamma_M$. 
%\item For an affine representation $\rho':\Gamma_M \ra \Aff(\mathds{A}^n)$, 
%we define $\mathds{A}^n_{\rho'\circ \Pi_{M\ast}}:= (\Uu \hat M \times %\mathds{A}^n)/\Gamma_M)$. 
%\end{itemize}  
%From now on, we will simply denote by $\bR^n_{\rho}$ and 
%$\mathds{A}_{\rho'}$. 
%There exists a flow $\Phi_t, t\in \bR,$ on $\bR^n_{\rho\circ \Pi_{M\ast}}$ acting 
%as the geodesic flow on the base space $\Uu M$ and 
%acting trivially on $\bR^n$ in $\Uu \hat M \times \bR^n$
%when lifted. 
%%There is also a flow $\hat \Phi_t$, $t \in \bR$, on 
%  $\mathds{A}^n_{\rho'\circ \Pi_\ast}$. 

%\marginpar{Should define completeness, developing maps and so on}
 
 %Let $\Gamma$ be a finitely-generated affine group that acts properly and freely on $\mathds{A}^n$. 
 %We assume that the Eilenberg-MacLane space $K(\Gamma, 1)$ for dimension one is realized as a finite complex, which is a fairly mild assumption in our field. 
 %Let $\mathds{A}^n/\Gamma$ be a manifold equipped with a complete Riemannian metric structure. 

Let $N$ be an $m$-manifold with a finitely generated fundamental group $\pi_1(N)$. 
Each flat vector bundle over $N$ can be described as 
$\bR^n_\rho(N):=\tilde N \times \bR^n/\sim$ where $(x, \vec{v}) \sim (\gamma(x), \rho(\gamma)(\vec{v}))$ for 
$\gamma \in \pi_1(N)$ where $\rho:\pi_1(N) \ra \GL(n, \bR)$ is a representation. 
A differential-geometry description of a flat bundle is a vector bundle with a flat linear connection. 
%We are allowed to have nonzero torsions even when $m =n$. 

Let $M$ be a compact submanifold (with or without boundary) in $N$.
%where $\Gamma$ acts properly discontinuously and freely on $\mathds{A}^n$. 
We assume that the inclusion map $M \hookrightarrow N$ induces a surjection on fundamental groups. The inverse image $\hat{M}$ of $M$ in $\tilde N$ is connected, and the deck transformation group $\Gamma_M$ of $\hat{M}$ is isomorphic to 
\green{$\pi_1(N)$. }
In these cases, we refer to $M$, the map $M \hookrightarrow N$, $\hat{M}$, the projection $p_{\hat{M}}: \hat{M} \ra M$, and $\Gamma_M$ as a \emph{fundamental group surjective {\rm (}FS\/{\rm }) submanifold}, an \emph{FS map}, an \emph{FS cover}, an \emph{FS covering map}, and an \emph{FS deck transformation group}, respectively.

% Let $M$ be a compact submanifold (with or without boundary) in
% $\mathds{A}^n/\Gamma$ where $\Gamma$ acts properly 
% discontinuously and freely on $\mathds{A}^n$. 
% We suppose that the inclusion map 
% $M \hookrightarrow \mathds{A}^n/\Gamma$ 
% induces a surjection of  the fundamental groups.
% The inverse image $\hat M$ of $M$ in $\mathds{A}^n$ is connected
% and the deck transformation group $\Gamma_M$ of $\hat M$ is isomorphic to $\Gamma$. 
% In these cases, we shall say 
%\[M, M \hookrightarrow \mathds{A}^n/\Gamma, \hat M, p_{\hat M}: \hat M \ra M, \hbox{ and } \Gamma_M\]  
% a {\em fundamental group surjective {\em (FS)} submanifold}, a {\em FS map}, 
%a {\em FS cover}, a {\em FS covering map} and a {\em FS deck transformation group} respectively. 

A \emph{compact core} of a manifold is a compact submanifold homotopy equivalent to the ambient manifold. 
When a manifold is homotopy equivalent to a finite complex $K$, for $n=3$, Scott and Tucker \cite{ST89} proved the existence of a compact core, while for $n - \dim K \geq 3$, Stallings \cite{Stallings65} established this fact. Unfortunately, compact cores do not always exist (see Venema \cite{Venema80}). 
Although compact cores are preferable, we need to resort to FS submanifolds. This is fine for our purposes since we are working in coarse geometry essentially. 


%A {\em compact core} of a manifold is a compact submanifold homotopy equivalent to the ambient manifold. 
%When a manifold is homotopy equivalent to a finite complex $K$, 
%for $n=3$, Scott and Tucker \cite{ST89} proved the existence of the compact core, and 
%while for $n- \dim K  \geq 3$, Stallings \cite{Stallings65} established this fact. 
%Unfortunately, sometimes they don't exist (See Venema \cite{Venema80}).
%Although it is better for our arguments when we have compact cores, 
%we need to use FS submanifolds. 
%Let $d_E$ denote the chosen standard Euclidean metric on $\mathds{A}^n$, fixed for this paper. Now, $\mathds{A}^n$ has an induced complete $\Gamma$-equivariant Riemannian metric from $N := \mathds{A}^n/\Gamma$, denoted by $\dN$, where $\tilde{N} = \mathds{A}^n$.

%Let $d_E$ denote a chosen standard Euclidean metric of $\mathds{A}^n$
%fixed for this paper. 
%Now,  $\mathds{A}^n$ has an induced complete  $\Gamma$-equivariant Riemannian 
%metric from $N:= \mathds{A}^n/\Gamma$ to be denoted by $\dN$,
%where $\tilde N = \mathds{A}^n$. 
 % giving us a Riemannian metric 
%  $d_M$ on $\mathds{A}^n/\Gamma = M$. 
%  Hence, $d_M = d_{\tilde N}$ in this article. 
We will assume that $\partial M$ is convex and smooth in $N$. Let $d_M$ denote the path metric on \purple{$M$} induced from a Riemannian metric on $N$, and let $d_{\hat{M}}$ denote the path metric on $\hat{M}$ induced from it.
%We will assume that $\partial M$ is convex in this paper. 
% Let $d_M$ denote the path metric induced  from a Riemannian metric on  $\mathds{A}^n/\Gamma$, 
%and we let $d_{\hat M}$ denote the path metric on $\hat M$ induced from it. 

%	Theorem \ref{thm:main} and Theorem  \ref{thm:main2} imply the result.  
  
  %which is $\delta$-hyperbolic.  
%  Since $\mathds{A}^n$ is identical with $\hat M$ under the developing map, 

  
%  We have fiber-wise norm $\llrrV{\cdot}_{\mathds{A}^n_{\rho'}}$ 
%  on $\mathds{A}^n_{\rho'}$ 
%  and a norm $\llrrV{\cdot}_{\bR^n_{\rho}}$ 
%   on $\bR^n_{\rho}$ using partition of unity.
%  This gives us translation invariant Euclidean metric-norms in the fiber-wise directions. 
%  Since $M$ has a Riemannian metric $d_M$, 
%  we obtain Riemannian metrics $d_{\tilde N_{\rho'}}$ on $\mathds{A}^n_{\rho'}$
%  and $d_{\bR^n_{\rho'}}$ on $\bR^n_{\rho}$ by choosing arbitrary sections. 
%  We will fix these in this paper. 

%Here, $M$ is a compact core of $\bR^n/\Gamma$. 
%We don't really know
%$\bR^n/\Gamma$  always has a compact core but this seems to be a reasonable assumption.
%(See \cite{Venema80}.) 
%If $n \geq cd(\Gamma) + 3$, the compact core exists by Stalling \cite{Stallings65}. )

\begin{itemize} 
\item Let $\Uu M$ denote the space of direction vectors on $M$, i.e., 
the space of elements of the form 
$(x, \vec{v})$ for $x\in M$ and a unit vector $\vec{v} \in T_x M \setminus \{O\}$ 
with the projection $\pi_{\Uu M}: \Uu M \ra M$. 
\item Let $\Uu \hat{M}$ denote the space of direction vectors on $\hat{M}$
covering $\Uu M$ with the deck transformation group $\Gamma_M$.  
Let $\pi_{\Uu \hat{M}} : \Uu \hat{M} \ra \hat{M}$ denote the projection.
 \item  Let $d_{\Uu M}$ denote the path metric on $\Uu M$ obtained from
  the natural Riemannian metric
on the unit tangent bundle $\Uu M$ of $M$ 
obtained from the Riemannian metric of $M$
as a sphere bundle.   We will use 
$d_{\Uu \hat{M}}$ to denote the induced path metric on $\Uu \hat{M}$. 
\end{itemize} 

%\marginpar{Jan 27: prob not needed} 
%Given metric space $(X, d_X)$ and $(Y, d_Y)$ and $L> 1, C> 0$, 
%a map $f: X \ra Y$ is an {\em $(L, C)$-quasi-isomeric embedding} if 
%$f$ satisfies 
%\[L^{-1} d_X(x, x') - C \leq d_Y(f(x), f(x')) \leq L d_X(x, x')  + C \hbox{ for all }x, x'\in X.\] 

%\begin{lemma} \label{lem:piUUM} 
%	$\pi_{\Uu \hat M}$ is a quasi-isometry.
%	\end{lemma} 
%\begin{proof}
%	Each fiber is uniformly bounded under $d_{\Uu \hat M}$.  
%	\end{proof} 
  
  A {\em complete isometric geodesic} in $\hat{M}$ is a geodesic that is an isometry of 
  $\bR$ into $\hat{M}$ equipped with \green{$d_{\hat{M}}$}. 
 Some but not all Riemannian geodesics are isometries. 
Actually, in \green{geodesic metric spaces}, these are identical to the notion of geodesics.
Since we are working on Riemannian metric spaces, we use the adjectives.  
  A {\em complete isometric geodesic} in $M$ is a geodesic that lifts to 
  a complete isometric geodesic in $\hat{M}$. 
%  Because the unfortunate terminology difference in the metric geometry and in Riemannian geometry, we use this terminology. 
\green{An} {\em isometric ray} in $\hat M$ is an isometric geodesic defined on  %\marginpar{ ray: meaning correctly used everywhere?}
$[0, \infty)$. 
%In this paper, the geodesic is always an isometric geodesic. 

% $a \in \bR$. 
  
%  We say that a sequence of rays (resp. complete isometric geodesics $l_i$ {\em converges} to a ray (resp. complete geodesic) $l$ 
%if $l_i(t)\ra l(t)$ for each $t \in [0, \infty)$ (resp. $t\in \bR$).
%
%By the Arzel\`a-Ascoli theorem and the properness of $\hat M$, 
%any sequence of rays (resp. complete isometric geodesics) 
%converges to a ray (resp. complete isometric geodesic) if 
%the sequence of the pairs of their $0$-points and directions at the %$0$-points converges in $\Uu \hat M$. 
%(See Section 11.11 of \cite{DK2018}.)
  
%  Suppose that $\hat M$ is Gromov hyperbolic. 
%We consider the subspace of $\Uu M$ where complete isometric geodesics pass. We denote this subspace by $\Uc M$, and call it the \emph{complete isometric geodesic unit-tangent bundle}. The inverse image in $\Uu \hat M$ is a closed set denoted by $\Uc \hat M$. Since a limit of a sequence of isometric geodesics is an isometric geodesic, $\Uc M$ is compact and $\Uc \hat

We consider the subspace of $\Uu M$ where complete isometric geodesics pass. We denote this subspace by $\Uc M$, and call it a \emph{complete isometric geodesic \green{unit tangent} bundle}. The inverse image in $\Uu \hat{M}$ is a closed set denoted by $\Uc \hat{M}$. Since a limit of a sequence of isometric geodesics is an isometric geodesic, $\Uc M$ is compact and $\Uc \hat{M}$ is closed and locally compact. However, $\pi_{\Uu M}(\Uc M)$ may be a proper subset of $M$ for the
projection $\pi_{\Uu M}: \Uu M \ra M$, and
$\pi_{\Uu \hat{M}}(\Uc \hat{M})$ may be a proper subset of $\hat{M}$.



%  We consider the subspace of $\Uu M$ where complete isometric geodesics pass. 
%  We denote this subspace  by $\Uc M$, and call it the {\em complete isometric geodesic
%  unit-tangent bundle}. 
%  The inverse image in $\Uu \hat M$ is a closed set 
%  denoted by $\Uc \hat M$. 
%  Since a limit of a sequence of isometric geodesics is 
%  an isometric geodesic, % by Theorem 11.45 of \cite{DK2018},  
%  $\Uc M$ is compact and $\Uc \hat M$ is closed and locally compact. 
%  However, $\pi_{\Uu M}(\Uc M)$ may be a proper subset of $M$
%  for the projection $\pi_{\Uu M}: \Uu M \ra M$, 
%  and $\pi_{\Uu \hat M}(\Uc \hat M)$ may be one of $\hat M$. 
 % Note that we will define these only when $\hat M$ is Gromov hyperbolic. 
  
 

%\begin{definition}[Strong semiproximality]\label{defn:semiproximal} 
%
%A representation $\rho: \Gamma_M \ra \SL_{\pm}(n, \bR)$ is {\em partially hyperbolic}  if 
%the following hold:  
%\begin{enumerate} 
%%\item[(i)] the flat bundle $\bR^n_{\rho}$ 
%%admits a section $\sigma: \Uu M \ra \bR^n_{\rho}$ which is flat along the geodesic lines
%%under the flat section. 
%\item[(i)] There exist nontrivial $C^0$-subbundles $\bV_+, \bV_0$, and $\bV_-$ in $\bR^n_{\rho}$ 
%with fibers $V_+(x), V_0(x), V_-(x) \subset V(x)$ for the fiber $V(x)$ over each point $x$ of $\Uc M$
%invariant under the flow $\Phi_t$. 
%%and $\iota(\bV_+) = \bV_-$ for 
%%an involution $\iota: \bR^n_{\rho} \ra \bR^n_{\rho}$ induced by 
%%an involution induced by the map reversing the orientation in $\Uu M$. 
%($V_+, V_0$ and $V_- $ are independent subspaces and their sum equals $V$.)
%\item[(ii)] For any fiber-wise metric on $\bR^n_{\rho}$ over $\Uc M$, 
%the lifted action of $\Phi_{t}$ on $\bV_{+}$ (resp. $\bV_{-}$) is dilating (resp. contracting). 
%\end{enumerate}
%\end{definition} 



%Let $\mathds{A}^n$ be a complete affine space. 
%An affine manifold is a manifold equipped with an atlas of charts to $\mathds{A}^n$ with 
%affine transition maps. 
%There is a homomorphism $\rho':\Gamma_M \ra \Gamma \subset \Aff(\mathds{A}^n)$ %called a {\em holonomy homomorphism}. 
%There is a inmmersion. $\dev: \hat M \ra \mathds{A}^n$ 
%so that 
%$\dev \circ \gamma = \rho'\circ \dev$ for each deck transformation 
%$\gamma \in \Gamma_M$. 
%An affine manifold is {\em special} if $\mathcal{L}(\Gamma)\subset \SL_{\pm}(n, \bR)$. 
%
%A {\em complete affine manifold} is a manifold $M$ 
%with a diffeomorphism $M \ra \mathds{A}^n/\Gamma$. 
%$\dev$ is a diffeomorphism if and only if the affine manifold $M$ is complete. 
%(See Thurston \cite{Thurston97}.)



Let $\rho: \pi_1(N) = \Gamma_M \ra \GL(n, \bR)$ be a homomorphism. 


There is an action of $\Gamma_M$ 
on $\Uc \hat{M} \times \bR^n$ given by 
\[ \gamma\cdot(x, \vec{v})  = (\gamma(x), \rho(\gamma)(\vec{v})) 
\text{ for } x \in \Uc\hat{M}, \vec{v} \in \bR^n.\]
Let $\bR^n_{\rho}$ denote the bundle over $\Uc M$
 given by taking the quotient $\Uc \hat{M} \times \bR^n$ by $\Gamma_M$. 
%by $\Gamma_M$ acting by $\rho$ on the $\bR^n$-factor.

Recall that there is a geodesic flow $\phi$ in $\Uu M$ that restricts to one in $\Uc M$, also denoted by $\phi$. This lifts to a flow, still denoted by $\phi$, on $\Uc \hat M$. Hence, there is a flow $\Phi$ on $\Uc \hat M \times \bR^n$, which descends to a flow, also denoted by $\Phi$, on $\bR^n_{\rho}$.


%Recall that there is a geodesic flow $\phi$ on $\Uu M$
%restricting to one on $\Uc M$, again denoted by $\phi$.  
%This lifts to a flow denoted by $\phi$ of $\Uc \hat M$.
%Hence, there is a flow $\Phi$ on $\Uc \hat M \times \bR^n$
%again descending to a flow to be denoted by $\Phi$ on $\bR^n_{\rho}$.  

\subsection{Main result} 
In this paper, we establish connections between the partial hyperbolicity in the context of dynamical systems and the \purple{Anosov} properties. 
%The significance and motivation of the main result are detailed in Part 2.

\reD{Unlike the abstract dominated splitting discussed in many previous articles such as \cite{BPS} and \cite{KP22}, here we require actual splitting of smooth bundles and smooth flows, along with the presence of expanding and contracting properties in addition to the domination property.}

%In this paper, we connect the partial hyperbolicity in the sense of the dynamical system and 
%the P-Anosov properties.  The significance and  the motivation of the main result are written in  Part 2. 

%\red{Unlike the abstract dominated splitting of many previous articles such as \cite{BPS} and \cite{KP22}, 
%we need actual splitting of smooth bundles and smooth flows and the expanding and contracting properties 
%in addition to the domination property.}

We need the following objects
since we need to work on manifolds and not just on groups. 
\begin{definition}[Partial hyperbolicity]\label{defn:phyp} 
	Suppose that $M$ is a compact Riemannian manifold with convex boundary.
	Let $\hat{M}$ be a regular cover of $M$ with the deck transformation group $G_M$.
%\purple{Assume that $G_M$ is word hyperbolic}. 
	A representation $\rho: G_M \ra \GL(n, \bR)$ is {\em partially hyperbolic in a bundle sense 
over $M$} if 
	the following hold for a Riemannian metric on $M$: 
	\begin{enumerate} 
		%\item[(i)] the flat bundle $\bR^n_{\rho}$ 
		%admits a section $\sigma: \Uu M \ra \bR^n_{\rho}$ which is flat along the geodesic lines
		%under the flat section. 
		\item[(i)] 
		There exist nontrivial $C^0$-subbundles $\bV_+, \bV_0$, and $\bV_-$ in $\bR^n_{\rho}$ 
		with fibers \[\bV_+(x), \bV_0(x), \bV_-(x) \subset \purple{\bR^n_{\rho}(x)}\] for the fiber $\purple{\bR^n_{\rho}}(x)$ of $\bR^n_{\rho}$ over each point $x$ of $\Uc M$
		invariant under the flow $\Phi_t$ of $\bR^n_{\rho}$ over 
		the geodesic flow $\phi_t$ of $\Uc M$. 
		%and $\iota(\bV_+) = \bV_-$ for 
		%an involution $\iota: \bR^n_{\rho} \ra \bR^n_{\rho}$ induced by 
		%an involution induced by the map reversing the orientation in $\Uu M$. 
		\item[(ii)] $\bV_+(x), \bV_0(x)$ and $\bV_-(x) $ for 
		each $x \in \Uc M$ are independent subspaces, and their sum equals $\bR^n_{\rho}(x)$.
		\item[(iii)] For any fiberwise metric on $\bR^n_{\rho}$ over $\Uc M$, 
		the lifted action of $\Phi_{t}$ on $\bV_{+}$ (resp. $\bV_{-}$) is dilating (resp. contracting). 
		That is, for coefficients $A> 0, a> 0$, $A' > 0, a'>0$\/: 
		\begin{enumerate} 
			\item $\llrrV{\Phi_{-t}(\vv)}_{\bR^n_\rho, \phi_{-t}(m)} \leq 
A \exp(-a t)\llrrV{\vv}_{\bR^n_\rho, m}$ for $\vv \in \bV_+(m)$ as $t \ra \infty$. 
			\item $\llrrV{\Phi_{t}(\vv)}_{\bR^n_\rho, {\phi_{t}(m)}} \leq 
A \exp(-a t)\llrrV{\vv}_{\bR^n_\rho, m}$ for \green{$\vv \in \bV_-(m)$} as $t \ra \infty$. 
			%\item $1/A \llrrV{\vv}_{m} \leq \llrrV{\vv}_{\phi_{t}^{\ast}(m)} \leq A \llrrV{\vv}_{m}$
			%for $\vv\in V_{0, m}$. 
			\item (Dominance properties) 
			\begin{multline} \label{eqn:dominance} 
				\frac{\llrrV{\Phi_{t}(\vw)}_{\bR^n_\rho,\phi_{t}(m)}}{\llrrV{\Phi_{t}(\vv)}_{\bR^n_\rho, \phi_{t}(m)}} 
				\leq A' \exp(-a't)\frac{\llrrV{\vw}_{\bR^n_\rho, m}}{\llrrV{\vv}_{\bR^n_\rho, m}} \\
				\text{ for } \vv \in \bV_+(m), \vw \in \bV_0(m)
				\text{ as } t \ra \infty, \\
				 \text{ or for }
				 \vv \in \bV_0(m), \vw \in \bV_-(m)
				 \text{ as } t \ra \infty.
			\end{multline} 
		\end{enumerate} 
	\end{enumerate}
	We additionally require that $\dim \bV_+ = \dim \bV_- \geq 1$.  
	Here $\dim \bV_+$ is the {\em partial hyperbolicity index} of $\rho$
	\purple{for the associated bundle decomposition $\bV_+, \bV_0, \bV_-$ 
	%where $\dim \bV_+$ is maximal
	and $\dim \bV_+ = \dim \bV_- < n/2$, so that $\dim \bV_0 \geq 1$. }
%	$\dim \bV_+$ is the index of the 
 Furthermore, $\bV_0$ is said to be the {\em neutral subbundle} of $\bR^n_{\rho}$. 
	(See Definition 1.5 of \cite{CP} or \cite{BLM}.)
We use the words ``in a bundle sense'' only as a convenience of this paper.)
\end{definition} 
%\green{We call the above properties without (iii)(a) and (iii)(b), the Anosov property of 
%$\rho$. The partial hyperbolicity is stronger than the Anosov property by Definition \ref{defn:phyp}. 
%By Theorem \ref{thm:identify}, there exists a quasi-isometry from $\Uc \hat M$ to the Gromov flow space $\partial^{(2)}_\infty M \times \mathbb{R}$ preserving the geodesic length parameter. Hence, the discussion reduces to one on the Gromov flow space, which is independent of the choice of FS submanifolds. 
%Our Anosov property implies that of Guichard-Wienhard \cite{GW12} and hence
%$\pi_1(N)$ is always hyperbolic. 
%} 

%The definition does depend on the Riemannian metric
%but not on $\llrrV{\cdot}_{\bR^n_{\rho}}$. 
%However, if
%$\Gamma_M$ is word hyperbolic, then we can replace $\Uc\hat M$ by 
%a Gromov flow space which depends only on $\Gamma_M$. 
%(See Theorem \ref{thm:identify} and 
%Gromov \cite{Gromov87} or Mineyev \cite{Mineyev}.)

%\purple{We plan to drop the condition of hyperbolicity of $G_M$ using \cite{KLP17}; however, we do not do so due to its length.} 

%\purple{The group $G_M$ is necessarily word-hyperbolic: 
%Geodesic flows contains all loops corresponding 
%to elements of $G_M$, 
%and $G_M$ is quasi-isometric to $\hat M$
%by the Svarc-Milnor theorem, and the word lengths
%are uniformly bounded above and below by the geodesic lengths.
%(See \cite{Milnor} and \cite{Svarc}). 
%Hence, $\rho$ is a $\dim V_+$-dominated representation.
%By Theorem 3.2 of Bochi-Potrie-Sambarino \cite{BPS}, this follows.}

The definition depends on the Riemannian metric, but not on $\llrrV{\cdot}_{\bR^n_{\rho}}$. If \purple{$G_M$} is word-hyperbolic, then we can replace $\Uc \hat M$ with a Gromov flow space which depends only on \green{$G_M$} up to quasi-isometry. (See Theorem \ref{thm:identify} and Gromov \cite{Gromov87} or Mineyev \cite{Mineyev}.)
\purple{Also, we leave as a question if Definition \ref{defn:phyp} forces 
$G_M$ to be word-hyperbolic.}

%\purple{Also, we probably do not need to require $G_M$ to be
%word hyperbolic; however, this question better be answered by people with more 
%expertise.} 

We say that $N$ has a \emph{partially hyperbolic linear holonomy group in 
a bundle sense} if
\begin{itemize}
\item there exists an FS submanifold $M$ in $N$ with convex boundary for a Riemannian metric on $N$, and
\item the linear part $\rho: \pi_1(N) = \Gamma_{M} \ra \GL(n, \bR)$ is a partially hyperbolic representation in a bundle sense over $M$. 
\end{itemize}

%We say that a complete affine manifold $N$ has a \emph{partially hyperbolic linear holonomy group} if
%\begin{itemize}
%\item there exists an FS submanifold $M$ in $N$ with convex boundary for a Riemannian %metric on $N$, and
%\item the linear part $\rho: \pi_1(N) \cong \Gamma_{\hat{M}} \to \GL(n, \mathbb{R})$ of an affine holonomy homomorphism $\rho'$ is partially hyperbolic in a bundle sense.
%\end{itemize}


%We say that a complete affine manifold $N$ has a {\em partially hyperbolic linear holonomy %group} if 
%\begin{itemize} 
%\item there exists an FS submanifold $M$ in $N$ with convex boundary for a Riemannian %metric of $N$, and
%\item  the linear part $\rho: \pi_1(N) = \Gamma_{\hat M} \ra \GL(n, \bR)$ of an affine %holonomy 
%homomorphism $\rho'$ is a partially hyperbolic  representation
%in a bundle sense. %This definition is independend of the choice of FS submanifold $M$ in %$N$. 
%\end{itemize} 

%By two paragraphs above, 
%the condition is independent of the choice of $FS$-manifolds.

%Given metric space $(X, d_X)$ and $(Y, d_Y)$ and $L> 1, C> 0$, 
%a map $f: X \ra Y$ is an {\em $(L, C)$-quasi-isomeric embedding} if 
%$f$ satisfies 
%\[L^{-1} d_X(x, x') - C \leq d_Y(f(x), f(x')) \leq L d_X(x, x')  + C \hbox{ for all }x, x'\in X.\] 
%%A quasi-isometry is an quasi-isometric embedding whose inverse 



%Note that the choice of FS submanifold $M$ does not change the partial hyperbolicity assumption. 
%Since partial hyperbolicity is equivalent to the Anosov property by Theorem \ref{thm:main2}, 
% $\pi_1(N)$ always is hyperbolic.
%By Theorem \ref{thm:identify} , there is a quasi-isometry from $\Uc \hat M$ to the Gromov flow space
%$\partial^{(2)}_\infty M \times \bR$ preserving the geodesic length parameter.
%Hence, the discussion reduces to one on the Gromov flow space, which is independent of 
%the choice of FS submanifolds. 

%\marginpar{need to show... I think}

%%  
%\begin{theorem}[Choi-Kapovich]\label{thm:main} 
%Let $N$ be a complete affine manifold  for $n \geq 3$ with the finitely presented %fundamental group. 
%Suppose that $N$ has a partially hyperbolic linear holonomy group and 
%$K(\pi_1(N), 1)$ is realized by a finite complex.  
%Then the cohomological dimension $\mathrm{cd}(\pi_1(N))$ is $\leq n-k$ for the partial %hyperbolicity index $k$ of $\rho$. 
%	% for any 
%	%opposite pair of parabolic groups $(P^{+}, P^{-})$.  
%\end{theorem} 
%A generalized stable affine subspace is an affine subspace parallel to lifted $\bV_0 \oplus \bV_-$ on $\mathds{A}^n$.  
%The main idea for proof is that 
%we will modify the developing map into a quasi-isometric embedding  into 
%a generalized stable affine subspace. 
%Hence, each boundary point of the group is associated with the affine subspaces. 


%	By Proposition 2.2 of \cite{AF1},following from Theorem 3.2 of \cite{BPS}, 
%	$\Gamma$ is hyperbolic. 
%Henceforth, we will assume that $\Gamma$ is hyperbolic.
%Also, $\hat M$ is $\delta$-hyperbolic since $\Gamma$ and $\hat M$ are roughly isometric. 


%\commt{Qustion: Strong semi-proximality also imply hyperbolic holonomy group? If %so, we
%	can drop a hypothesis above.}

%We generalize Corollaries 1.3 and 1.4 of \cite{AF1}: that is, we drop the condition of 
%negative curved Riemannian manifold for $M$. 

%The following includes the case when $\rho$ is strongly irreducible, i.e., $\rho$ restricted to every finite-index subgroup is irreducible.  
%(See Section 3.10 of \cite{BS2018p}.)

We denote by $a_1(A), \dots, a_n(A)$ the \green{nonincreasing} set of singular values of 
an $n\times n$-matrix $A$.
We denote by $\lambda_1(A), \dots, \lambda_n(A)$ the nonincreasing set of norms of 
eigenvalues of an $n\times n$-matrix $A$. (Here, the multiplicity is allowed.) 

Let $\log \lambda_1, \dots, \log \lambda_n$ denote the logarithmic eigenvalue functions defined on
the maximal \green{abelian} Lie algebra of $\mathfrak{gl}(n, \bR)$. 
For roots 
\[\theta=\left\{\log \lambda_{i_{1}}-\log \lambda_{i_{1}+1}, \dots, \log \lambda_{i_{m}}- \log \lambda_{i_{m}+1}\right\}, 
1 \leq i_{1}<\cdots<i_{m} \leq n-1,\] 
of $\GL(n, \bR)$, 
the parabolic
subgroup $P_{\theta}$ (resp. $\left.P_{\theta}^{-}\right)$ is 
the subgroup of block upper 
(resp. lower) triangular matrices in $\GL(n,\bR)$ with square diagonal blocks 
of sizes $i_{1}, i_{2}-i_{1}, \dots, i_{m}-i_{m-1}, n-i_{m}$ respectively. 
\red{We denote by $\theta^*$ the image of $\theta$ under the map $i \mapsto n-i$ for $i=1, \dots, n-1$.}


\purple{For a finitely generated group $G$, we denote by $w(g)$ the word length of $g$ in terms of a fixed finite generating set $S$ in $G$.}

\begin{theorem}\label{thm:main2} 
Let $N$ be an $m$-manifold with a finitely generated \purple{word-hyperbolic} fundamental group. 
%where $K(\pi_1(N), 1)$ is realized by a finite complex. 
	%Suppose that $M$ admits a negatively curved Riemannian metric. 
	Let $\rho: \pi_1(N)\ra \GL(n, \bR)$ be a representation.
Assume one of the following: 
\begin{itemize} 
\item[(i)] \purple{$|\lambda_{i_g}(\rho(g))| =1 $} for all \green{$g \in \pi_1(N)$} where 
the index $i_g \in \{1, \dots, n\}$ may depend on $g$. 
\item[(ii)] \purple{For any sequence $g_i \in \pi_1(N)$, if $w(g_i) \ra \infty$, then $\frac{|\log a_{j_{g_i}}(\rho(g_i))|}{w(g_i)} \ra 0$ for some choices of 
the indices $j_{g_i}\in \{1, \dots, n\}$ depending on $g_i$.} 
\end{itemize} 
%	Suppose that the image of $\rho$ has a reductive
%	Zariski closure.  
%	\marginpar{ I still need $\rho$ to be semi-simple}
%	Suppose that there is a number $k$, $1\leq k \leq n/2$ 
%	and $\rho$ satisfies the neutrality condition for $[k+1, n-k]$.
	%so 
	%that there is a function $j:g \in \Gamma_M \ra [k+1, n-k]$ 
	%so that the function $g\in \Gamma_M \ra a_{j(g)}(\rho(g))$ is
	%subexponential. 

Then $\rho$ is $P$-Anosov for 
	any parabolic subgroup $P$ of $\GL(n, \bR)$ with 
	$P = P_\theta$ %for $\theta$ 
\red{for $\theta = \theta^*$}
containing $\log \lambda_k-\log \lambda_{k+1}, k \leq n/2$ 
	if and only if $\rho$ is partially hyperbolic of the index $k$ in a bundle sense. 
	In these cases $k < n/2$. 
%	then $cd(\Gamma) \leq n-k$ if $P = P_\theta$ for $\theta$ containing $\lambda_k-\lambda_{k+1}, k \leq n/2$. 
	%where $I$ is a subset of 
	%$\{1, \dots, n\}$ not containing $k$ and $n-k$
	%for $k$, $1\leq k \leq n/2$. 
\end{theorem}

%The theory of Anosov representations as begun by Guichard and Wienhard \cite{GW12} are well established with currently plentitudes of surveys.
%For a survey of the growing field of partially hyperbolic dynamics, refer to Crovasier and Potrie \cite{CP}. 


\purple{We say that $\rho$ has {\em subexponential growth for some singular values} if (ii) holds.}  \purple{By Lemma 2.11 of \cite{BS2021} applied to the standard representation of 
$\GL(n, \bR)$, (ii) implies (i); however, we 
wrote out (ii) for the purpose of the usefulness of the result.
Also, the convergence in (ii) is uniform with respect to the word length since 
we can find a seqeunce not converging to $0$ otherwise.}

\red{It is easy to see that the conclusion is true for representations to 
$\SL_\pm(n, \bR)$ and for $k$ which is the least element of $\theta$ without our 
conditions.} \purple{If $\pi_1(N)$ is not simple, we can multiply any representation  
by a representation to $\bR^+$ so that a partial hyperbolicity cannot hold.}
%\purple{Also the irreducibilty of $\rho$ is necessary since we can take a direct sum with 
%a trivial representation and a representation as stated immediately above.}
We can ask if an Anosov representation is always partially hyperbolic up to multiplying by a positive scalar representation to $\bR$. We cannot answer this question since
controlling the singular values of representation elements is very hard in general, and
the only related results that we know of are due to Abel-Margulis-Soifer \cite{AMS95}, which is not applicable here
as far as we know. \red{See Kassel-Potrie \cite{KP25} for some discussions. 
However, only cases that are left are for representations 
whose singular values behave exponentially.}





\reD{The significance of this result is that we work with the full general linear group 
$\GL(n, \bR)$ instead of orthogonal groups. Also, Anosov representations correspond to only dominated representations and not actual partially hyperbolic flows as here. Promoting the dominated representation to Anosov ones is trivial for some Lie groups but not for general linear groups. }


\begin{corollary}\label{cor:affine}
Suppose that $N$ is a complete affine $n$-manifold with a finitely generated 
\purple{word-hyperbolic} fundamental group. 
Let $\rho:\pi_1(N) \ra \GL(n, \bR)$ be a linear part of the affine holonomy group. 
Then the conclusions of Theorem \ref{thm:main2} hold. 
\end{corollary} 

In the field of affine manifolds, 
the hyperbolic flow was already discovered by Goldman and Labourie \cite{GL12} for $3$-dimensional Margulis space-times, and Ghosh \cite{Ghosh17}, Ghosh-Treib \cite{GhoshTreib21}, and Danciger-Zhang \cite{DZ2019} prove these properties for affine actions on higher-dimensional affine spaces. However, the linear parts are always assumed to be in the orthogonal groups or Hitchin, where the properties are more or less immediate.



%If each holonomy element has a unit norm eigenvalue, then the whole argument of 
%Theorem \ref{thm:main2} passes. However, only significant class of examples where this condition holds seem to be the 
%holonomy representations of the complete affine manifolds.


%\red{
%The hyperbolic flow was already discovered by Goldman and Labourie \cite{GL12} for $3$-dimensional Margulis space-times, 
%and Ghosh \cite{Ghosh17}, Ghosh-Treib \cite{GhoshTreib21}  and Danciger-Zhang \cite{DZ2019} prove these properties for affine actions on higher dimensional affine spaces. However,  the linear parts are always assumed to be in the orthogonal groups where the properties are more or less immediate.}

%\red{
%The significance of this result is that we work with the full general linear group $\GL(n, \bR)$ instead of orthogonal groups. Also, Anosov representations correspond to only dominated represenations and not actual partially hyperbolic flows as here. Promoting the dominated representation to Anosov ones are trivial for the some Lie groups but not for general linear groups. }

%We aim to prove in Part 2 the following theorem worked out with Kapovich. 
%\begin{theorem}[Theorem \ref{thm:main-p2}]\label{thm:mainofpart2} 
%Let $N$ be a complete affine manifold  for $n \geq 3$ with the finitely presented fundamental group. 
%Suppose that $N$ has a partially hyperbolic linear holonomy group with index $k, k < n/2$,  and $K(\pi_1(N), 1)$ is realized by a finite complex.  
%
%Then the cohomological dimension $\mathrm{cd}(\pi_1(N))$ is $\leq n-k$ for the partial hyperbolicity index $k$ of $\rho$. 
	% for any 
	%opposite pair of parabolic groups $(P^{+}, P^{-})$.  
%\end{theorem} 
%\red{This result coincides with Canary-Tsouvalis \cite{CT2020} on linear Anosov representations for $\SL(n, \mathbb{R})$ up to minor additional arguments, while our paper employs a significantly different method of homotopying the manifold into an affine subspace. However, we aim to generalize most of these methods and results to complete affine manifolds with relatively Anosov holonomy representations, a field that is still developing. See Kapovich-Leeb \cite{kapovich2023relativizing}, Zhu \cite{Zhu21,Zhu23}, and Zhu-Zimmer \cite{zhu2022p}. The approach of Canary-Tsouvalis \cite{CT2020}, which builds upon the work of Bestvina-Mess \cite{BM91}, essentially does not extend straightforwardly to the context of relativizing.}

%\red{This result does coincide with Canary-Tsouvalis \cite{CT2020} on linear Anosov representations for $\SL(n, \bR)$ up to minor additional arguments
%while our paper uses much different method of homotopying the manifold into an affine subspace. 
%However, we are trying to generalize the most of methods and results to the complete affine manifolds with relatively Anosov holonomy representations. But the later field is still developing. See Kapovich-Leeb \cite{kapovich2023relativizing} and Zhu \cite{Zhu21},\cite{Zhu23}, and Zhu-Zimmer \cite{zhu2022p}. The work of Canary-Tsouvalis \cite{CT2020} which uses the work of Bestvina-Mess \cite{BM91} essentially does not work when we do relativizing.}


%\red{The reason for dividing into two parts is that the main techniques are distinct, and it seems distracting to switch gears in the middle of the paper. We believe this approach is more reader-friendly. The first part focuses on partially hyperbolic flows, primarily working with singular values and flows. The second part delves into coarse geometry to a greater extent.}

%\red{Additionally, the techniques involved--- Anosov representations, partially hyperbolic flows, and coarse geometry --- may not be familiar to traditional geometric topologists as they are rapidly developing fields. However, the results used are grounded in well-established papers.}

We wish to approach the conjecture that a closed complete affine manifold cannot have a word-hyperbolic fundamental group. One approach considered is using Oseledets-type theorems (See Chapter 8 of \cite{KH95}). However, generalizing our results to measurable settings is currently hindered by numerous technical issues. Discussions on these matters took place with researchers such as Nicholas Tholozan at BIRS, Banff, during the winter of 2020.

%At the moment, we have a result that such a closed complete affine manifold cannot have 


%red{The reason for dividing into two parts  is that the main techniques are distinct, and it %seems to be distracting to switch gear in the middle of the paper. We believe that this is kinder to the readers. The first part is on partially hyperbolic flows mostly working with singular values and flows. The second part uses much more coarse geometry.}

%\red{Also, the techniques involved are Anosov representations, partially hyperbolic flows, and coarse geometry. These are not necessarily familiar topics to usual geometric topologists since they are currently or recently developing rapidly. However, the
%used results are from well-established papers.} 

%This paper was an attempt to solve the conjecture that a closed complete affine manifold cannot have word hyperbolic fundamental groups.
%One avenue is to use the Osedelec type theorems (See Chapter 8 of \cite{KH95}.) 
%For now, we cannot yet generalize our results to measurable settings
%due to many technical issues.  We talked about these issues with some people such as Nicholas Tholozan at BIRS, Banff  duing the winter 2020. 

 
%Also,we 
%really produce the flow instead of just defining using the rough geodesics. 

\subsection{Some historical comments on complete affine manifolds} 
Originally, Milnor posed the question whether a free affine group of rank $\geq 2$ could act freely and properly discontinuously on affine space. Margulis provided an affirmative answer by constructing examples in the affine $3$-space, where the linear parts belong to $\SO(2, 1)$. Drumm and Goldman subsequently explained this result using fundamental domains in a series of papers (see \cite{DDGS} for a survey). Sullivan \cite{Sullivan76} observed that the linear part of each affine group element must have $1$ as an eigenvalue and proved the Chern conjecture for complete affine manifolds \cite{KS75}. Klingler \cite{Klingler2017} later extended this to closed special affine manifolds.

Further developments include Ghosh \cite{Ghosh17,Ghosh18,ghosh2018avatars,Ghosh23}, Ghosh-Treib \cite{GhoshTreib21}, and Smilga \cite{Smilga14,Smilga16,Smilga18,Smilga22,Smilga22ii}, who generalized these results and produced higher-dimensional examples of proper and free affine actions. Abels, Margulis, and Soifer also contributed to related topics such as the Auslander conjecture (see Abels \cite{Abels01} for a survey). Danciger and Zhang \cite{DZ2019} investigated affine deformations of Hitchin representations, showing that they do not exist. Danciger, Kassel, and Gu\'eritaud \cite{DGK2020} provided \purple{the} first examples of fundamental groups of closed manifolds acting properly and freely on affine spaces.

Although these works are related to the themes of this paper, their specific aims differ. This summary only scratches the surface of this rapidly evolving field.

%\red{
%The Chern conjecture states that a closed affine manifold must have zero Euler characteristic. 
%Originally Milnor asked if an affine free group with rank $\geq 2$ could act the affine space freely and properly discontinuously. 
%Sullivan \cite{Sullivan76} observed that the linear part of each element of the affine group has to have $1$ as an eigenvalue and proved the Chern conjecture when the affine manifold is complete \cite{KS75}, 
%Recently Bruner Klingler \cite{Klingler2017} proved it for all closed special affine manifolds.
%Margulis answered this question producing an example for 
%the affine $3$-space. The examples have linear parts in $\SO(2, 1)$. 
%Drumm and Goldman produced series of papers explaining the result using the fundamental domains. 
%(See \cite{DDGS} for the survey on this topic.) Ghosh \cite{Ghosh17}, \cite{Ghosh18}, \cite{ghosh2018avatars} , \cite{Ghosh23}, Ghosh-Treib \cite{GhoshTreib21}, and Smilga \cite{Smilga14}, \cite{Smilga16}, \cite{Smilga18}, \cite{Smilga22}, and \cite{Smilga22ii}
%generalized the results and produced many more higher dimensional examples of proper and free actions. Abels, Margulis, and Soifer also worked on the Auslander conjecture which is related to this work. (See Abels \cite{Abels01} for a survey.) 
%Danciger and Zhang \cite{DZ2019} proved that there are no proper affine deformations of Hitchin representations. 
%Danciger, Kassel, and Gu\'ertaud \cite{DGK2020} produced examples of closed manifold fundamental groups acting on affine spaces properly and freely. 
%Many of these works are related to this paper but their aims are somewhat different.
%Of course, we do not think that this short summary does any justice to this growing field.
%}





%Since we can always find FS submanifolds for $\mathds{A}^n/\Gamma$, 
%Theorem \ref{thm:main} and Theorem  \ref{thm:main2} imply the result: 
%we have 
%\begin{corollary}
%	Let a finitely presented group $G$ acts on $\mathds{A}^n$  faithfully, properly discontinuously,  and freely. 
%	Suppose that $K(G, 1)$ is realized by a finite complex. 
%	Suppose that the linear part of $G$ is P-Anosov. Then 
%	$\mathrm{cd}(G) \leq n-k$ where $P = P_\theta$ for $\theta$ containing %$\lambda_k-\lambda_{k+1}, k < n/2$.
%	
%\end{corollary}
%%\begin{proof} 
%%	 Theorem 4.1 of Stallings \cite{Stallings65} shows that given a homotopy equiavence of a $2$-complex $K$ into 
%%	a manifold $N$, we can find a simple homotopy equivalent $2$-complex $K'$ %embedding into $N$ so that 
%%	we can obtain a compact core from taking a neighborhood of the image of $K'$.  
%%	Theorem \ref{thm:main} and Theorem  \ref{thm:main2} imply the result.  
%%	\end{proof} 
%
%Of course this also essentially follows from Theorem 1.3 of Canary-Tsouvalas \cite{CT2020} using deep Corollary 1.4 of 
%Bestvina-Mess \cite{BM91}. Although we have more assumptions, 
%our methods are substantially different and use more direct geometrical arguments. 
%Our main point here is that we avoid heavy machinery. 



%When $n= 4$, we don't know the existence of the compact core. 

%We conjecture that we can drop the reductiveness here. 
%For closed flat complete pseudo-Riemannian or special symplectic affine manifolds, we do not need the semisimplicity. 
%
%\begin{corollary}[Choi-Kapovich] \label{cor:main2} 
%	Let $M$ be  a closed complete special affine $n$-manifold with 
%	$n \geq 3$ and 
%	a fundamental group $\Gamma_M$ with linear holonomy in 
%	$\SO(k, n-k)$ for each integer $k$, $0\leq k \leq n$, 
%	or $\mathsf{SP}(m, \bR)$ for $n= 2m$. 
%	Then the linear part of 
%	the holonomy homomorphism $\rho$ is not $P$-Anosov for 
%	any parabolic group $P$ of $\SL_\pm(n, \bR)$.
%	%where $I$ is a subset of 
%	%$\{1, \dots, n\}$ not containing $k$ and $n-k$. 
%\end{corollary}


%e proved the following which supports the Auslander conjecture. 
%
%\begin{corollary}
%	A closed complete affine manifold $M^n$ cannot have a $P$-Anosov linear holonomy group
%	for a parabolic subgroup $P$ of $\GL(n, \bR)$. 
%	\end{corollary}  
%\begin{proof} 
%	If otherwise,  $\mathrm{cd}(\Gamma_M) = n \leq n -k $ for any $k$, $1\leq k \leq n/2$ %and $k$ in $\theta$ for $P = P_\theta$. 
%	\end{proof} 
%%We conjecture that a closed affine manifold whether it is complete or not does not
%%have an Anosov linear holonomy. 

%\begin{corollary} \label{cor:cpt} 
%	Suppose that $\rho: \pi_1(N) \ra \GL(n, \bR)$ be a $k$-Anosov representation
%	that is a linear part of a properly discontinuous and free affine action on %$\mathds{A}^n$. 
%	There exists a compact collection of affine subspaces of dimension $n-k$ in the %affine Grassmannian space%
%	$\mathcal{AG}_{n-k}(\bR^n)$  invariant under the affine action. 
%\end{corollary} 


%%The above results are clearly true for $n=1, 2$
%%by Benzecri \cite{Benzecri} since closed affine $2$-manifolds must have
%%Euler characteristic $0$. 
%The Auslander conjecture is proved for closed 
%complete affine manifolds of dimension $\leq 3$ by 
%Fried-Goldman \cite{FG83}, for ones with linear holonomy groups in
%the Lorentz group by Goldman-Kamishima \cite{GK84}, 
%and for ones of dimension $\leq 6$ by 
%Abels-Margulis-Soifer \cite{AMS02}, \cite{AMS05}, \cite{AMS11}, and \cite{AMS97}. In particular, they showed that 
%the linear holonomy group is not Zariski dense in $\SO(k, n-k)$ 
%for $(n-k) -k \geq 2$ in \cite{AMS11}.
%%Also, there is really interesting work by Bucher-Connell-Lafont 
%%\cite{BCL} showing that 
%%the simplicial norm of a closed complete affine manifold 
%%is zero if there is a translation inside the affine holonomy group. 

%Examples of complete affine manifolds with free fundamental groups were obtained by 
%Drumm \cite{Drumm92} and Margulis \cite{Margulis84}. 
%The existence of properly discontinuous affine actions on 
%$\mathds{A}^n$ for large classes of groups including all cubulated hyperbolic groups 
%were discovered by Danciger, Kassel, and Gu\'eritaud \cite{DGK2020}
%where $n$ is somewhat large compare to 
%$\mathrm{cd}(G)$ of the properly acting affine group $G$.   

%Also, there are various recent results on properly discontinuous actions of affine groups by Danciger-Zhang \cite{DZ2019} and 
%Canary-Tsouvalas \cite{CT2020} and Tsouvalas \cite{Tsouvalasp} 
%giving us various obstructions to the existences of affine representations acting properly on $\mathds{A}^n$ with Anosov linear parts. 
%%We showed $n > vcd(G)$ in this article for group $G$ acting properly on $\bR^n$. 
%%Hopefully, there are better lower bounds on $n$ by
%%$vcd(G)$ by generalizing these theories. Our method in this article could work in a more general setting;
%%however, this may take much more work.  
%%We believe that our method also gives obstructions in a more general setting
%%than here perhaps giving us the lower bound on $n$ in terms of the 
%%virtual cohomological dimensions. 



%It is well-known that 
%some discrete representations of the surface group giving us
%infinite ends of hyperbolic $3$-manifolds as quotients of 
%the hyperbolic spaces are
%non-Anosov representations. 
%The discreteness is not enough to imply an Anosov condition. 

%We cannot generalize the theorem to non-reductive Zariski closure cases. In these cases, the linear holonomy of the affine manifolds 
%normalizes nontrivial unipotent radicals. 
%Perhaps, this can be the direction for further research. 

%A well-known conjecture weaker than the Auslander conjecture is that 
%a complete closed affine manifold cannot have a word hyperbolic 
%fundamental group. 
%We might try to generalize our arguments using 
%Osedelec's theorem. 
%We believe that our approach may be a step in the right direction. 

%\marginpar{Jan 27, I need to revise these remarks}



\subsection{Outline}


In Section \ref{sec:preliminary-p1}, we will introduce  affine structures and 
some basic facts on $\delta$-hyperbolic metric spaces. 
Suppose that $\Gamma_M$ is word-hyperbolic.
Then $\hat M$ has a Gromov hyperbolic Riemannian metric
by the Svarc-Milnor theorem. 
(See \cite{Milnor} and \cite{Svarc}. 
\purple{\v{S}varc is a transliteration of Albert Schwarz's name.
He published this theorem in Russian prior to Milnor. Currently, he is 
a professor at the University of California, Davis.)}
%First, we show that each affine subspace intersected with $\hat M$ is uniformly contractible. % when $\hat M$ is a complete affine space. 
%and is a Gromov hyperbolic covering space of a closed affine manifold $M$. 
We discuss the ideal boundary of $\hat M$, identifiable with the Gromov boundary, and the complete isometric geodesics in $\hat M$. 
We relate the space of complete isometric geodesics to
the Gromov flow space. 
%Finally, we show that the set of complete isometric geodesics in $\hat M$ ending 
%at a common point of the ideal boundary $\partial_\infty \hat M$ is $C$-dense in $\hat M$ for some $C> 0$. 


%We prove Theorem \ref{thm:main} in Sections \ref{sec:neutral}
%and \ref{sec:geoconv}: 


%%In Section \ref{sec:neutral}, we will define an affine bundle associated with an FS %%submanifold $M$ of a closed complete special affine manifold. 
%%We suppose that we have a partially hyperbolic linear representation. 
%%In Theorem \ref{thm:neutral}, 
%%we will modify the developing section of $\Uc \hat M$ so that
%%each complete isometric geodesic in $M$ develops inside an affine space in the neutral directions. The modification follows from the idea of 
%%Goldman-Labourie-Margulis \cite{GLM09}. 
%%We define $\mathcal{R}_p$ for $p\in \partial_\infty M$ to be the subspace of points on complete isometric geodesics on $\Uc \hat M$ ending at an ideal point $p$.
%%%A generalized stable subspace is an affine subspace parallel to lifted $\bV_0 \oplus \bV_-$ on $\mathds{A}^n$.  
%%Proposition \ref{prop:Finalpart} shows that $\mathcal{R}_p$ for each $p\in \partial_\infty M$ always develops into a generalized stable subspace. 
%%This follows since along the unstable directions, geodesics depart away from one another. 

%%\marginpar{I need to change: I proved the second part here and do the cohomology later} 
%In Section \ref{sec:geoconv}, we will prove Proposition \ref{prop:quasiiso}   %%\ref{prop:at_least_two}
%that $\hat M$ quasi-isometrically embed into generalized stable subspaces
%since $\mathcal{R}_p$ embeds quasi-isometrically into one of the subspace, and
%$\hat M$ and $\mathcal{R}_p$ are quasi-isometric. 
%%that there are at least two distinct generalized stable affine subspaces where geodesics
%%ending at a point of $\partial_\infty \hat M$ 
%%go under the modified developing map
%%using the coarse geometry idea of Section \ref{sec:dimension_pres}. 
%%The main tools are from Section \ref{sec:dimension_pres}. 
%%However, Proposition \ref{prop:Finalpart} will show that every geodesic goes to a single generalized stable affine subspace 
%%under the modified developing map, using the affine geometry and the partially %hyperbolic flow property. 
%%This contradiction proves Theorem \ref{thm:main}. 
%%Hueristically speaking what we are showing is that each generalized stable leave of the unit tangent bundle of 
%%$\hat M$ must go into generalized stable affine plane. We will formalize this statement later. 
%Then we prove Theorem \ref{thm:main2}: We use the quasi-isometric embedding of $\hat M$ into 
%generalized stable affine subspace to show that 
%the maximal dimension of the compactly supported cohomology of $\hat M$ is less than 
%the dimension of the generalized stable subspace $n-k$.
%Since $\mathds{A}^n/\Gamma$ homotopy equivalent to $K(\Gamma, 1)$ 
%has an exhaustion by a sequence of  FS submanifolds $M_i$, 
%we will obtain the upper bound  $n-k$ of the cohomological dimension of $\Gamma$.
%(Here, we may need to change the Riemannian metric of $N$ so that $M_i$ has convex boundary.)

%In Section \ref{sec:dimension_pres},  we will prove that $\hat M$ cannot be in a bounded neighborhood of 
%an affine subspace of strictly lower dimension using the coarse geometry.
%The uniform contractibility of affine subspaces in $\hat M$ 
%will be used.   
%\marginpar{I think I can generalize this to vcd argument} 

In Section \ref{sec:phyp}, 
%we will show that if the linear part of the holonomy representation is partially hyperbolic, then it is $P$-Anosov. 
%Here, 
we provide a review of the definition of the $P$-Anosov property of 
Guichard-Wienhard \cite{GW12}. For convenience for reductive groups, we will use 
the approaches of G\'ueritaud-Guichard-Kassel-Wienhard \cite{GGKW17ii} and Bochi-Potrie-Sambarino \cite{BPS}
using singular values of the linear holonomy group. (The theory of 
Kapovich, Leeb, and Porti \cite{KL18}, \cite{KLP18iii} is equivalent to these, but it is not so
tailored for our purposes here.)
We give various associated definitions. 
%In Section \ref{sub:background}, we will review  parabolic subgroup $P$.


Section \ref{sec:PAnosov} has the main arguments. 
In Section \ref{sub:kconvexity}, we will relate the partial hyperbolic property in the singular-value sense to that in the bundle sense. 
In Section \ref{sub:converse}, we show that the linear part of 
the holonomy representation of a complete affine manifold 
is $P$-Anosov if and only if it is partially hyperbolic. 
First, we do this when the linear holonomy group
has a reductive Zariski closure. 
\purple{We use the fact that $1$ has to be an eigenvalue of each element 
or the subexponential growth of some singular values of
the linear holonomy groups}. We relate it
to singular values using the clever ideas of Danciger and Stecker,
the spectrum theory of Breuillard-Sert \cite{BS2021}, Benoist \cite{Benoist97}, \cite{Benoist98}, and the work of Potrie-Kassel \cite{KP22}.
We could generalize to the cases of nonreductive linear holonomy groups using 
the stability of the Anosov and partial hyperbolic properties. 
This will prove Theorem  \ref{thm:main2}. % and \ref{cor:main2}. 
%Finally, we prove Corollary \ref{cor:cpt}.  

\purple{In Section \ref{sec:affine}, we apply these results to complete affine manifolds.  
proving Corollary \ref{cor:affine}.} 


\subsection{Acknowledgments}
We thank Michael Kapovich for various help with geometric group theory and coarse geometry. This article began with some discussions with Michael Kapovich during the conference honoring the 60th birthday of William Goldman at the University of Maryland, College Park, in 2016. 
%Important suggestions for proofs of our main Theorems \ref{thm:main2} and \ref{thm:mainofpart2} were made by Michael Kapovich.
We thank Jeffrey Danciger and Florian Stecker for help with the $P$-Anosov properties, which yielded the main idea for the proof of Theorem \ref{thm:main2}. 
%These authors declined to join me as the coauthors for various personal reasons. 
%We also thank Herbert Abels, Richard Canary, Virginie Charette, Todd Drumm,
%William Goldman, Fran\c{c}ois Gu\'eritaud, Fanny Kassel, 
%Andr\'es Sambarino Nicholas Tholozan, and Konstantinos Tsouvalas for various discussions helpful for this paper. 
%We also thank BIRS, Banff, Canada, and KIAS, Seoul, where some of this work was done. 
%\red{We remark that this series  of articles are produced by joint work with Danciger, Kapovich, and Stecker.
%	They did not wish to join the author in writing these papers unfortunately. 
%	We also thank some referees for pointing out how the article was unclear about the originality of specific results in 
%	this paper.)}
%
%
%We thank Michael Kapovich with various help with geometric group theory and coarse geometry. This article began with some discussions with Michael Kapovich during the conference honoring the 60th birthday of William Goldman at the University of Maryland, College Park, in 2016. 
%We thank Jeffrey Danciger and Florian Stecker for help with $P$-Anosov properties which yielded the main corollary of the paper. 
We also thank Herbert Abels, Richard Canary, Virginie Charette, Todd Drumm,
William Goldman, Fran\c{c}ois Gu\'eritaud, Fanny Kassel, 
Andr\'es Sambarino, Nicholas Tholozan, and Konstantinos Tsouvalas for various discussions helpful to this paper. 
We also thank BIRS, Banff, Canada, and KIAS, Seoul, where some of this work was done.
\reD{This paper is a joint work with Kapovich in the beginning and Danciger and Stecker later. Unfortunately, they did not wish to join as coauthors for various reasons.} 
\purple{Also, we used 
OpenAI to assist us with writing and Google Gemini to find related results and papers.}

\section{ Preliminary} \label{sec:preliminary-p1} 

%\subsection{Grassmannians} \label{sub:Haudorff} 
%We assume $n \geq 3$ in this article.  
%Let $\mathcal{G}_k(\bR^n)$ denote the space of $k$-dimensional subspaces of 
%$\bR^n$. 
%We consider the space $\mathcal{AG}_k(\bR^n)$ of affine $k$-dimensional subspaces of $\bR^n$.
%The space has a complete Riemannian metric which we denote by $d_{\mathcal{AG}_k(\bR^n)}$, which is a proper metric. 
%%We will use the Haudorff distances $\bdd^H$ on 
%%the spaces of subspaces of $\mathds{A}^n$ with a fixed Euclidean metric $d_E$. 
%%embedded as an open subspace of a projective $n$-space $\RP^n$ 
%%with Fubini-Study metric $\bdd$ on $\RP^n$. 
%We also use these on subspaces of $\bR^n$ considered as 
%$\mathds{A}^n$. 

\reD{We recall some standard well-known facts in this section to set the notations and so on. 
Also, we need a slight modification of the theories for our purposes. Most of the background on geometric groups and coarse geometry is from the comprehensive textbook by 
Dru\c{t}u and Kapovich \cite{DK2018}.}

The other purpose is to establish the Anosov properties for manifolds using coarse isometries since most of the literature works \green{with groups}.


\subsection{Convergences of geodesics} \label{sub:geoconv}  
Let $M$ be a manifold of dimension $n$ with a regular cover $\hat{M}$ and the deck transformation group $\Gamma_M$. We assume that $M$ has a Riemannian metric with convex boundary.

We say that a sequence of isometric rays (resp. complete isometric geodesics) $l_i$ \emph{converges} to an isometric ray (resp. complete geodesic) $l$ (denoted as $l_i \ra l$) if $l_i(t) \ra l(t)$ for each $t \in [0, \infty)$ (resp. $t \in \bR$).

By the Arzel\`a-Ascoli theorem and the properness of $\hat{M}$, any sequence of isometric rays (resp. complete isometric geodesics) converges to an isometric ray (resp. complete isometric geodesic) if the sequence of pairs of their $0$-points and directions at the $0$-points converges in $\Uu \hat{M}$. (See Section 11.11 of \cite{DK2018}.)

%Let $M$ be a manifold of dimension $n$ with a regular cover $\hat M$ and the deck transformation group $\Gamma_M$. 
%We suppose that $M$ has a Riemannian metric with convex boundary. 
%  We say that a sequence of rays (resp. complete isometric geodesics) $l_i$ {\em converges} to a ray (resp. complete geodesic) $l$ (that is, $l_i \ra l$)
%if $l_i(t)\ra l(t)$ for each $t \in [0, \infty)$ (resp. $t\in \bR$).

%By the Arzel\`a-Ascoli theorem and the properness of $\hat M$, 
%any sequence of rays (resp. complete isometric geodesics) 
%converges to a ray (resp. complete isometric geodesic) if 
%the sequence of the pairs of their $0$-points and directions at the $0$-points converges in $\Uu \hat M$. 
%(See Section 11.11 of \cite{DK2018}.)




%\subsection{Metrics and affine subspaces} 
% Let us recall some definitions: 
% Let $(X, d_X)$ and $(Y, d_Y)$ be metrics spaces. A map $f: X \ra Y$ is called 
%{\em  $(L, C')$-coarse Lipschitz} for $L, C > 0$ if 
% \[d_Y(f(x), f(x') \leq L d_X(x, x') + C' \hbox{ for all } x, x' \in X.\] 
% A map $f: X \ra Y$ is an {\em $(L, C)$-quasi-isomeric embedding} if 
% $f$ satisfies 
% \[L^{-1} d_X(x, x') - C \leq d_Y(f(x), f(x')) \leq L d_X(x, x')  + C \hbox{ for all }x, x'\in X.\] 
% A map $\bar f: Y\ra X$ is a {\em $C$-coarse inverse} of $f$ 
% for $C > 0$ if 
%\[d_X(\bar f \circ f, \Idd_X) \leq C, \hbox{ and }
%d_Y(f\circ \bar f, \Idd_Y)\leq C.\]  
%A map $f:X \ra Y$ between metric spaces is called 
% a {\em quasi-isometry} if it is a coarse Lipschitz and admits a 
% coarse Lipschitz coarse inverse map.

 
%A metric space $(X, d)$ is {\em proper} if every metric ball has a compact closure. 
%In other words, $d_p(\cdot) = d(p, \cdot)$ is a proper function $X \ra \bR$. Clearly,
%a complete Riemannian manifold is a proper metric space.
%We denote by $B^d(x, R)$ the ball of radius $< R$ in $X$ with center $x, x \in X$. 

%From Definition 8.27 of \cite{DK2018}, we recall: 
%A map $f:X \ra Y$ between two proper metric spaces $(X, d_X)$ and $(Y, d_Y)$ 
%is {\em uniformly proper} if $f$ is coarsely Lipschitz and there is a function $\psi: \bR_+ \ra \bR_+$ such that 
%\[d_X\hbox{-diam}(f^{-1} (B^{d_Y}(y, R))) < \psi(R) \hbox{ for each }y \in Y, R\in \bR_+.\]
%An equivalent condition is that there is a proper continuous function $\eta: \bR_+ \ra \bR_+$ 
%so that 
%\[ d_Y(f(x), f(y))\geq \eta(d_X(x, y))) \hbox{ for all } x, y \in X. \]
%Here, functions satisfying the properties of 
%$\psi$ and $\eta$ respectively are called an {\em upper} 
%and {\em lower 
%distortion functions}. 

%Recall that 
%$M$ be a closed manifold covered by $\hat M$  
%to be identified with $\mathds{A}^n$ with a $\Gamma_M$-invariant metric $ \dN$, 
%which is a path metric from the Riemannian metric pulled from $M$. 
%%That is, $d_M =  \dN$ under the identification. 
%%Let $d_{\hat M}$ denote a path metric on $\hat M$ induced from 
%%a Riemannian metric of $M$. 
%%\marginpar{notation? $d_M$? right or $d_{\mathds{A}}^n$}


%A subspace $Y$ in a  metric space $(X, d)$ is {\em uniformly contractible} in a subspace $Y'$, $Y \subset Y'$,  if 
%for every  $r> 0$, there exists a real number $R(r) > 0$ depending 
%only on $r$ so that $B^d_r(x) \cap Y$ is contractible in $B^d_{R(r)}(x) \cap Y'$
%for any $x \in Y$.
%(We generalize Block and Weinberger \cite{BW93} and Gromov \cite{Gromov93}.) 
%%for the definition of the uniform $k$-connectedness of a manifold, $k=1, 2, \dots$.) 
%%The $0$-level uniform contractibility is equivalent to 


%%% Oct 3 5:10pm 

%Let $d_E$ denote any Euclidean metric on $\mathds{A}^n$ fixed for this paper. 
%The restricted path metric on any affine subspace $L$ from $\mathds{A}^n= \hat M$ is denoted $d_L$. 
%We will later show that $L$ is uniformly connected in 
%the sense of Gromov \cite{Gromov93} in Theorem \ref{thm:k-conn}. 
%The following is the $0$-th level uniform connectedness. 

%For an affine subspace $L$ of $\mathds{A}^n$, we denote by 
%$d_L$ the restricted metric of $ \dN$. Note that this is not the path metric induced from 
%the restricted Riemannian metric to $L$. This is just the plain restriction of the distances. 

%We obviously have: 
%\begin{proposition}\label{prop:unifprop}
%	Suppose that $M$ is an FS submanifold of a complete affine manifold $N$ covered by $\mathds{A}^n$ with 
%	invariant path metric $d_N$ induced from a Riemannian metric. 
%	Let $\hat M$ be the cover of $M$ in $\mathds{A}^n$. 
%	Let $L$ be an affine subspace of $\mathds{A}^n$
%	with induced metric $d_L$ from $d_N$ on $L \cap \hat M$. 
%	Then $L\cap \hat M$ with metric $d_L$ is uniformly properly embedded in $\hat M$ with metric $d_M$. 
%\end{proposition}
%%\begin{proof}
	
%	The uniform contractibility for $L$ proved in 
%	Theorem \ref{thm:k-conn}. 
%	This property indeed gives us the upper distortion function
%	$\psi(r)$ set as $R(r)$ obtained there. 
	%$\psi(r) = R(r)$. %obtained there. 
%The proof of Proposition \ref{prop:unifprop} 
%is implied by the proof of Theorem \ref{thm:k-conn} 
%since the pthe uniform connectedness for $\pi_0(L)$
%and the proof applies to this case. 
%\end{proof} 

%%We adopt a result of Gromov \cite{Gromov93}: 
%\begin{theorem}[Choi-Kapovich] \label{thm:k-conn} 
%		Suppose that $M$ is an FS submanifold of a complete affine manifold $N$ covered by $\mathds{A}^n$ with 
%	invariant path metric $d_N$ induced from a Riemannian metric.
%	Let $L$ be an affine subspace of $\mathds{A}^n$ of $\dim \leq n$. 
%	Let $\hat M \subset \mathds{A}^n$ be the cover of $M$ under the covering map %$\mathds{A}^n \ra N$. 
%	Then $L \cap \hat M$ is uniformly contractible in $L$. 
%	% with respect to the path metric of 
%%	the restriction on $L $ of
%%	the Riemannian metric associated with $d_N$. 
%\end{theorem} 
%\begin{proof} 
%	%	Let $d_M$ denote the path metric on $\hat M$ induced from
%	%	the Riemannian metric of $M$. 
%	Let $F$ be a compact fundamental domain of $\hat M \subset \mathds{A}^n$, 
%	containing the origin $O$. 
%	Let $L'$ be any affine subspace of dimension $\dim L \leq n$.
%	Let $d_{L'}$ denote the path metric on $L'$ induced from  $ \dN$. 
%	Let $r$ be any positive real number. 
%	The $d_{L'}$-ball $B^{d_{L'}}_r(x)$ in $L'$ of radius $r> 0$ 
%	for $x \in F$ is a subset of 
%	$B^{d_{\hat M}}_r(x)$  for a $d_{\hat M}$-ball of radius $r$ 
%	with center $x \in F$ 
%	%	for $R(r, F)$ depending only on $r$ and $F$ 
%	since the endpoints of a $d_{L'}$-path of length $< r$ has 
%	$d_{\hat M}$-distances $< r$ from $x$. 
%	Since the $\bigcup_{x\in F} B^{d_{\hat M}}_r(x)$ is bounded in $d_N$, 
%	there is a constant $R(r, F)$ depending only on $r$ and $F$  
%	so that $B^{d_{\hat M}}_r(x) \subset B_{R(r, F)}(O)$
%	for the Euclidean ball $B_{R(r, F)}(O)$ of radius 
%	$R(r, F)$ with center $O$. 
	
%%	Let $R > 0$. 
%%	The values	$d_{L'}(x, \clo(B_R(O)) \cap L')$, $x\in F$, 
%%	form a continuous function of  the set 
%%Consider the set
%%	\[\{(x, L')| x \in F, x\in L', L'\hbox{ is an affine subspace of dimension }\dim L\}\]
%%	with the metric given by $d_{\hat M}$ and $\bdd^H$. 
%%	%times the space of affine subspaces of $\dim L$ 
%%	%passing $F$ with the $\bdd^H$-metric topology. 
%%	Since this space is a compact metric space,
%	We take $C(R, F)$ for each $R> 0$ to be the supremum of 
%	\[\{ \dN(x, B_R(O))|  x \in F\}.\] 
%%	for 
%%	every affine space $L'$ containing $x\in F$ of dimension $\leq n$. 
%	Then $B_R(O)\cap L' \subset B^{d_{L'}}_{C(R, F)}(x)\subset L'$
%	for $x\in F$ and any affine subspace $L'$ with $\dim L'=\dim L$ containing $x \in F$. 
	
%	Now, $B_R(O) \cap L'$ is convex
%	and is a subset of $B^{d_{L'}}_{C(R, F)}(x)$. 	
%	Since \[B^{d_{L'}}_r(x) \subset B_{R(r, F)}(O) \cap L', x\in F,\]  
%	$B^{d_{L'}}_r(x)$ is contractible to a point 
%	inside $B^{d_{L'}}_{C(R(r, F), F)}(x) \subset L'$.
	
%	Since we can put any  $B^{d_{L}}_r(x)$ for $x\in L \cap \hat M$ to 
%	a $d_{\gamma(L)}$-ball with the center in $F$ 
%	by a deck transformation $\gamma$ of $\hat M$, 
%	we obtained the radius $C(R(r, F), F)$ for each $r> 0$ so that 
%	the uniform contractibility holds.  
%	%\qed
%	%	
%	%	
%	%	
%	%	For any $x$ in $L$, we choose a deck transformation $\gamma$ so that 
%	%	$\gamma(x) \in F$. Then for each $B_{d_{\hat M}, r}(x)$ is contained 
%	%	in a Euclidean ball $B_R(0)$ for 
%	%	
%	%	
%\end{proof} 


%\begin{proof} 
%	It will be sufficient to show \eqref{eqn:dL} since 
%	for each $D$ we can find $C_D$ as the value of 
%	the upper distorsion function. 
%For each $D> 0$, 
%we have 
%\begin{multline} \label{eqn:dL} 
%	d_{M}(z_1, z_2)\leq 
%	d_{L}(z_1, z_2) \leq C_D,  
%	\hbox{ for }
%	z_1, z_2 \in 	L 
%	\hbox{ with } d_{M}(z_1, z_2) \leq D
%\end{multline} 
%where $C_D$ depends on $D$: 
%The left inequality follows since paths in $L$ can be shortened. 
%%Suppose first that $z_1$ is in a fundamental domain $F$. 
%Let $F$ be a compact fundamental domain of $\hat M$. 
%%The space of affinesubspaces of dimension $\dim L$ passing $F$
%%is compact. 
%%Take a constant $1/2 < r < 1$. 
%For any point $z \in L'\cap F$ and an affine subspace $L'$ of dimension $\dim L$, 
%we take a convex hull $C_{L', r, z}$ of $B^{d_{\hat M}}_{D+1}(z)\cap L'$ 
%%so that $d_{\hat M}(z, \partial_{L'} C_{L', r, z}) \geq R$ for the boundary $\partial_{L'} C_{L', r, z}$ 
%%of $C_{L', r, z}$ in $L'$. 
%Let $H(F)$ denote the set of hyperplane $L'$ of dimension $\dim L$ passing a 
%point of $F$. 
%We define 
%\[c_z := \max_{L'\in H(F)} C_{z, K'}
%\hbox{ where } c_{z, L'}:= d_E(z, \partial C_{L', r, z}).\]
%Since $C_{L', r, z}$ is a subset of the compact convex hull  of 
%$B^{d_{\hat M}}_{D+1}$
%Since $\partial C_{L', r, z} \subset \clo(B^{d_E}_{C_z}(z))$,  
%the space
% \[\{\partial_{L'} C'_{L', r, z} |L' \in H(F)\}\]
% is compact under geometric convergences with the Hausdorff metric from
% $d_{\hat M}$.  
% Then for every $L' \in H(F)$, 
%$d_{L'}(z, \partial C'_{L', r, z}) \leq D'$ for a constant $D'$ depending on %$r$ and $F$ since 
%we may take affine geodesics from $z$ to $\partial C'_{L', r, z}$ 
%and estimate the length. 
%Hence, any shortest path between $z_1$ and $z$ on $L'$ must be inside $C'$ 
%and is $\leq D'$. 
%\end{proof} 


%This follows since the curvature of $L$ is bounded above uniformly 
%with respect to the Riemannian metric of $d_{M}$ always
%and $L$ is properly embedded. (Proof: normalize 
%by elements of $\Gamma_M$, and use geometric convergences)

\subsection{Gromov hyperbolic spaces} \label{sub:Gromov} 

%We will work under the assumption that $\Gamma_M$ is $P$-Anosov. 
%Then by Kapovich-Lee-Porti \cite{KLP17} or Bochi-Potrie-Sambarino \cite{BPS}, $\Gamma_M$ is Gromov hyperbolic. 

 Let us recall some definitions: 
A metric space is {\em geodesic} if every pair of points are connected by a geodesic, i.e., 
a path that is an isometry from an interval to the metric space.  
We will only work with geodesic metric spaces.
 Let $(X, d_X)$ and $(Y, d_Y)$ be metric spaces. A map $f: X \ra Y$ is called 
{\em  $(L, C')$-coarse Lipschitz} for $L, C > 0$ if 
 \[\green{d_Y(f(x), f(x'))} \leq L d_X(x, x') + C' \hbox{ for all } x, x' \in X.\] 
 A map $f: X \ra Y$ is an \green{{\em $(L, C)$-quasi-isometric embedding}} if 
 $f$ satisfies 
 \[L^{-1} d_X(x, x') - C \leq d_Y(f(x), f(x')) \leq L d_X(x, x')  + C \hbox{ for all }x, x'\in X.\] 
 A map $\bar f: Y\ra X$ is a {\em $C$-coarse inverse} of $f$ 
 for $C > 0$ if 
\[d_X(\bar f \circ f, \Idd_X) \leq C, \hbox{ and }
d_Y(f\circ \bar f, \Idd_Y)\leq C.\]  
A map $f:X \ra Y$ between metric spaces is called 
 a {\em quasi-isometry} if it is a coarse Lipschitz map \purple{that} admits a 
 coarse Lipschitz coarse inverse map.
(Note that these are not necessarily continuous, as we follow 
Dru\c{t}u-Kapovich \cite{DK2018}.) 

\begin{lemma} \label{lem:piUUM} 
Let $M$ be a compact Riemannian manifold with possibly nonempty boundary. 
	$\pi_{\Uu \hat M}$ is a quasi-isometry.
	\end{lemma} 
\begin{proof}
	Each fiber is uniformly bounded under $d_{\Uu \hat M}$.  
	\end{proof} 

 
A metric space $(X, d)$ is {\em proper} if every metric ball has a compact closure. 
In other words, $d_p(\cdot) = d(p, \cdot)$ is a proper function $X \ra \bR$. Clearly,
a complete Riemannian manifold is a proper metric space.
We denote by $B^d(x, R)$ the ball of radius $< R$ in $X$ with the center $x \in X$. 

Let $X$ be a geodesic metric space with metric $d$. 
A {\em geodesic triangle $T$} is a concatenation of three geodesics $\tau_1, \tau_2, \tau_3$ where the indices are modulo $3$. 
The {\em thinness radius} of a geodesic triangle $T$ is the number 
\[ \delta(T) := \max_{j=1,2,3} \left( \sup_{p \in \tau_j} d(p, \tau_{j+1}\cup \tau_{j+2}) \right). \]


A geodesic metric space $X$ is called {\em $\delta$-hyperbolic} in the sense of Rips if 
every geodesic triangle $T$ in $X$ is $\delta$-thin. 

By Corollary 11.29 of \cite{DK2018} as proved by Section 6.3C of 
Gromov \cite{Gromov87}, 
the Gromov hyperbolicity is equivalent to the Rips hyperbolicity
for geodesic metric spaces.
We will use these concepts interchangeably. 


%We will assume that $\Gamma_M$ is Gromov hyperbolic to define these objects. 
We will assume that $\Gamma_M$ is word-hyperbolic for our discussions below. 
This is equivalent to assuming that
$\hat M$ is Gromov hyperbolic by the Svarc-Milnor lemma.



%	A {\em ray} in a geodesic metric space $X$ is a closure in $X$ of a component of an isometric geodesic in $X$ defined from $[0, \infty)$. 
%Since $\hat M$ is quasi-isometric to the Cayley graph $X$ of $\Gamma_M$, 
%$\hat M$ is $\delta$-hyperbolic by Corollary 11.43 of \cite{DK2018}.
Two isometric rays $\rho_1$ and $\rho_2$ of $X$ are {\em equivalent} if 
$t\mapsto d(\rho_1(t), \rho_2(t))$ is bounded.  
This condition is equivalent to the one that their Hausdorff distance under $d$ is bounded. 
Assume that $X$ is $\delta$-hyperbolic. $X$ has a well-defined ideal boundary 
$\partial_\infty X$ as in Definition 3.78 of \cite{DK2018}; 
that is, the space of the equivalence classes of isometric rays. 
%(see Eberlein \cite{Eberlein}). 
% Let $\partial_\infty \hat M$ denote the Gromov boundary of $\hat M$. 
 (See \cite{CDP90} and \cite{GdH90}.) 
 We denote by $\partial^{(2)}_\infty X = \partial_\infty X \times \partial_\infty X \setminus \Delta$ 
 where $\Delta$ is the diagonal set. 

The constant $C$ below is called a {\em quasi-geodesic} constant. 
 \begin{lemma} \label{lem:twogeo} 
	Let $M$ be a compact manifold with 
	the induced path metric $d_{\hat M}$ on $\hat M$
	from a Riemannian metric of $M$.
	Let $\Gamma_M$ be the deck transformation group of $\hat M \ra M$.  
	Suppose that $\Gamma_M$ is word-hyperbolic. 

Then there exists a constant $C >0$ so that
for every pair of
	complete isometric geodesics $l_1$ and $l_2$ in $\hat M$
with the identical set of endpoints in $\partial_\infty \hat M$, the Hausdorff distance between 
$l_1$ and $l_2$ under $d$ is bounded above by $C$. 
\end{lemma} 
\begin{proof}
%Since $M$ is compact, $\Gamma_M$ and $\hat M$ are quasi-isometric by 
%Svarc-Milnor lemma again. 
	Since $l_1$ and $l_2$ are both $(1, 0)$-quasi-geodesics, 
	this follows from Proposition 3.1 of Chapter 3 of \cite{CDP90}. 
	%\qed
\end{proof} 

 
 


 	We put on $\hat M \cup \partial_\infty \hat M$ the 
 shadow topology, which is a first-countable Hausdorff topology
 by Lemma 11.76 of \cite{DK2018}. It is also compact according to Section 11.11 of \cite{DK2018}. 
 Since it is first countable, 
 we do not need to consider nets on the space but only sequences
 to understand the continuity of the real-valued functions. 

We denote by $\partial_G X$ the Gromov boundary of 
a $\delta$-hyperbolic geodesic metric space $X$. (See Section 11.12 of \cite{DK2018}.)
 By Theorem 11.104 of \cite{DK2018}, 
 there is a homeomorphism $h: \partial_\infty X
 \ra \partial_G X$ given 
 by sending the equivalence class $[\rho]$ of an isometric ray $\rho$
 to the equivalence class of $\{\rho(n)\}_{n \in \bZ_+}$. 
 We will identify these two spaces using this map. 
 %The space is homeomorphic to the space of 
 %complete geodesics in $\hat M$. 
 %\marginpar{Some problem.. correct for this version}

Let $\hat{\Gamma}_M$ denote the Cayley graph of $\Gamma_M$. 
$\hat{\Gamma}_M$ and $\hat M$ are quasi-isometric by 
the Svarc-Milnor lemma again. 
The Gromov boundary $\partial_\infty \hat M$ of $\hat M$ and 
%the Gromov boundary of $X$ 
can be identified 
%which is again 
with the boundary $\partial_\infty \Gamma$ of
$\hat{\Gamma}_M$ with the word metric
by Theorem 11.108 of \cite{DK2018}. 
% and the theorem of 
%Svarc-Milnor \cite{Milnor} and \cite{Svarc}. 

For any Gromov hyperbolic geodesic metric space $X$, 
the set of $GX$ of complete isometric geodesics has a metric: 
for $g, h \in GX$, we 
define the metric $d_{GX}$ given by 
\begin{equation}\label{eqn:dGX} 
 d_{GX}(g, h) := \int^\infty_{-\infty} d_X(g(t), h(t)) 2^{-|t|} dt.    
\end{equation}
Notice that \purple{$d_X(g(t),h(t)) \leq 2|t| + C$ for a constant $C$}
since we can follow the geodesic $g$ backward and 
then forward in $h$ to obtain a distance $2|t|$. Thus, the above integral is always 
well defined. 
(See Gromov \cite{Gromov87}.)
%
% defined from 
%$\Uc \hat M$ by 
%
%considering the uniform convergences in the intersections with compact subsets in $\hat M$. 
Let $GX$ have this metric topology. 
%$GX$ is clearly locally closed since the limit of 
%a sequence of isometries from $\bR$ is an isometry from $\bR$. 
(See Section 11.11 of \cite{DK2018}.)


\begin{lemma} \label{lem:identify} 
	Suppose that $\Gamma_M$ is Gromov hyperbolic. 
	Then there is a quasi-isometric homeomorphism $\mathcal{F}:\Uc \hat M \ra G\hat M$
	by taking a unit vector $\vec{u}$ at $x$ to a complete isometric geodesic 
	$\bR \ra \hat M$ passing $x$ tangent to $\vec{u}$. 
	\end{lemma} 
\begin{proof}
There is a map $G\hat{M} \ra \Uc \hat{M}$ given by sending \purple{each} complete isometric geodesic $g: \bR \ra \hat{M}$ to $(g(0), \vec{v}_0)$ where $\vec{v}_0$ is the unit tangent vector at $g(0)$. This map is clearly bijective. In Section III of \cite{Matheus90}, the map $G\hat{M} \ra \hat{M}$ defined by $g \mapsto g(0)$ is shown to be a quasi-isometry. Therefore, post-composing this map with a lift that sends each $g(0)$ to a vector in the fiber of $\Uc \hat{M}$ over $g(0)$ results in a quasi-isometry. For a sequence $\{g_i\}$ of geodesics in $\hat{M}$, if $d_{GX}(g, g_i) \to 0$, then $(g_i(0), g_i'(0)) \to (g(0), g'(0))$ since otherwise, we would obtain a positive lower bound for the integral \eqref{eqn:dGX}.

The inverse map $\Uc \hat{M} \ra G\hat{M}$ is also continuous due to the continuity of the exponential map and considerations involving \eqref{eqn:dGX}, where $d_X(g(t), h(t))$ grows sublinearly. The continuity of the integral values under $g$ and $h$ can be shown by cutting off for $|t| > N$ where $N$ is independent of $g$ or $h$ by \eqref{eqn:dGX} and an estimate. 


%	There is a map $G\hat M \ra \Uc \hat M$ given by 
%	sending the complete isometric geodesic $g: \bR \ra \hat M$ to 
%	$(g(0), \vec{v}_0)$ where $\vec{v}_0$ is the unit tangent vector 
%	at $g(0)$. This map is clearly one-to-one and onto. 
%	In Section III in \cite{Matheus90}, the map $G\hat M \ra \hat M$ given by 
%	$g \ra g(0)$ is shown to be quasi-isometry. 
%	Hence, this map post-composed by a lift sending
%	each $g(0)$ to a vector in the fiber of 
%	$\Uc \hat M$ over $g(0)$ is a quasi-isometry. 
%	For a sequence $\{g_i\}$ of geodesics in $\hat M$, 
%	 if $d_{GX}(g, g_i)\ra 0$, then 
%	$(g_i(0), g_i'(0)) \ra (g(0), g'(0))$ 
%	since otherwise we easily obtain a positive lower bound of the integral \eqref{eqn:dGX}. 
%%	Also, by a sequence argment, the continuity of the unit tangent 
%%	vector at $g(0)$ follows by the continuity of the exponential map
%%	defined on $\Uu \hat M$. 
	
%	The inverse map $\Uc \hat M \ra G\hat M$ 
%	is also continuous by the continuity of
%	the exponential map and considering
%	equations of the form \eqref{eqn:dGX}
%	where $d_X(g(t), h(t))$ grows sublinearly. 
%	We can show the continuity of the integral values
%	under $g$ and $h$ by
%	cutting off for $|t| > N$. 
	\end{proof} 



%%% July 26 2020 12:10AM 
%
%Recall that an {\em  lower semi-continuous} function $f: \hat M \ra \bR_+$ is 
%a function $f(x_0) \leq \liminf_{x \ra x_0} f(x)$. 
%A lower semi-continuous function always achieve an infimum. 
%(See \cite{Vinogradova} for details.)
%A function $f$ is {\em $C$-roughly continuous} for $C> 0$ if 
%\[|\liminf_{x\ra x_0} f(x) - f(x_0)| \hbox{ and }
%|\limsup_{x\ra x_0} f(x) - f(x_0)| < C.\] 


%\begin{proposition} \label{prop:GtildeM} 
%	$\Uc \hat M$ is quasi-isometric with $\Gamma_M$. 
%\end{proposition}
%\begin{proof} 
%	There is a map $\Uc \hat M \ra \hat M$ given by 
%	a restriction of the projection. 
%	This has a $C$-dense image for some $C> 0$
%	since each point of the fundamental domain $F$ has an image point. 
%	Then we follow by a quasi-isometry 
%	$\hat M \ra \Gamma_M(x_0)$ for a point $x_0\in F
%	\cap \pi_{M}(\Uc \hat M)$
%	sending $\gamma(F^0)$ to $\gamma(x_0)$
%	and extending arbitrarily for $\gamma(\delta F)$ to 
%	adjacent orbit points. 
%%	By Proposition \ref{prop:Ucdense}, 
%	We can map $\Gamma_M$ into $\Uc \hat M$ 
%	by taking a point of $\Gamma_M(x_0)$ and lifting to $\Uc\hat M$ 
%	as a quasi-isometric inverse.
%	\end{proof} 

\subsection{Flow space $G\hat M$ quasi-isometric to $\partial_\infty^{(2)}\hat M \times \bR$}


We assume that $\Gamma_M$ is word-hyperbolic. 
For a word-hyperbolic group $\Gamma$, we let $\hat \Gamma$ denote its Cayley graph.
For simplicity, we denote by $\partial_\infty \Gamma$ the boundary 
$\partial_\infty \hat \Gamma$. 
We denote \[\partial_\infty^{(2)}\Gamma
:=\partial_\infty\Gamma\times\partial_\infty \Gamma \setminus \{(t,t)\mid t\in\partial_\infty \Gamma\}.\]
We denote by $F\Gamma_M$ the Gromov 
geodesic flow space of $\hat{\Gamma}_M$ obtained by the following proposition. 
(Note that we identified $\partial_\infty \hat M$ with $\partial_\infty \Gamma$ in
Section \ref{sub:Gromov}.) 


\begin{theorem}[Theorem 8.3.C of Gromov \cite{Gromov83}, Theorem 60 of Mineyev \cite{Mineyev}] 
	\label{thm:Mineyev} Let $\Gamma$ be a finitely generated word-hyperbolic group. 

Then there exists a proper hyperbolic metric space 
	$F\Gamma$ with the following properties\/{\em :}
	\begin{itemize} 
		\item[(i)] $\Gamma\times\bR\times\bZ/2\bZ$ acts on $F\Gamma$.
		\item[(ii)] The $\Gamma\times \bZ/2\bZ$-action is isometric.
		\item[(iii)] Every orbit $\Gamma\ra F\Gamma$ is a quasi-isometry. %In particular, %$\partial_\infty\Gamma=\partial_\infty\hat \Gamma$.
		\item[(iv)] The $\bR$-action is free, and every orbit $\bR \ra F\Gamma$ is a quasi-isometric embedding. 
		The induced map $F\Gamma/\bR \ra \partial_\infty^{(2)}\Gamma$ is a homeomorphism.
	\end{itemize}
\end{theorem}
We shall say that $F \Gamma$ is the {\em flow space} of $\Gamma$. 
In fact, $F \Gamma$ is unique up to a $\Gamma \times \bZ/2\bZ$-equivariant quasi-isometry sending $\bR$-orbits to $\bR$-orbits. 
We shall denote by $\phi_t$ the $\bR$-action on $F \Gamma$ and by 
$(\tau_+, \tau_-):F\Gamma \ra F\Gamma/\bR \cong \partial_\infty^{(2)}\Gamma$
the maps associating to an element of $F\Gamma$ 
the endpoint of its $\bR$-orbit.
Gromov identifies $F\Gamma$ with 
$\partial_\infty^{(2)} \Gamma \times \bR$ with a certain metric 
called the Gromov metric. 
%
%\begin{corollary} \label{cor:identify}
%	There is a $\Gamma_M \times \bZ/2\bZ$-equivariant 
%	quasi-isometry from $\Uc \hat M$ to  
%	 $\partial_\infty^{(2)} \hat M \times \bR$ with the
%	 Gromov metric by a map which sends each complete 
%	 isometric geodesic to $(x, y)\times \bR$, $(x, y) \in 
%	 \partial_\infty^{(2)} \hat M$  by an isometry.
%	\end{corollary} 
%\begin{proof} 
%	$\Uc \hat M$ satisfies the (i),(ii), and (iii) of 
%	Theorem \ref{thm:Mineyev}. 
%	The composition $f$ of maps in Lemma \ref{lem:identify} and Proposition \ref{prop:Uc} is obviously $\Gamma_M\times \bZ/2\bZ$-equivariant. 
%By Theorem \ref{thm:identify}, 
%$\partial_\infty^{(2)} \hat M \times \bR$ is a quotient 
%metric space of $\Uc \hat M$ which is also quasi-isometric by $f$. 
%With this metric $\partial_\infty^{(2)} \hat M \times \bR$ 
%satisfies (i)-(iv) of Theorem \ref{thm:Mineyev} also.
%$\Uc \hat M$ is quasi-isometric with 
%$\hat M$ by $f$, 
%and the Gromov space is quasi-isometric with 
%$\hat M$ by a map $\psi$ in  Section 8.5 of \cite{Mineyev}. 
%These maps give us a quasi-isometry between 
%$\Uc \hat M$ and the Gromov space 
%since we can use the same projection choices. 
%\end{proof} 

%Let $\Pi: \Uu M \ra M$ denote the fibration 
%and $\hat \Pi: \Uu \hat M \ra \hat M$ the induced fibration. 
%\begin{itemize} 
%	\item Take a representation $\rho: \Gamma_M \ra \SL_{\pm}(n, \bR)$, and 
%	we define 
%	$\chi_{\rho} := (\Uu \hat M \times \chi)/\Gamma_M$
%	with the bundle projection map $\pi_\rho: \chi_\rho \ra \Uu M$. 
%\end{itemize}  

%By Theorem \ref{thm:Mineyev}, 
%By Proposition 4.8 of \cite{Champetier94}, 
%the flow space $\hat \Gamma_M$ 
%is homeomorphic with $\partial_\infty^{(2)}\hat M \times \bR$. 
%We will give the induced metric on 
%$\partial_\infty^{(2)}\hat M \times \bR$. 



%We give the space 
%$\partial_\infty^{(2)}\hat M \times \bR$ the metric as in the paragraph after Theorem \ref{thm:Mineyev}. 


\begin{proposition} \label{prop:Uc} 
	Let $M$ be a compact manifold with a covering map $\hat M \ra M$
	with a deck transformation group $\Gamma_M$. 
	Suppose that $\Gamma_M$ is word-hyperbolic.

	Then there is 
	a quasi-isometric surjective map 
	$\mathcal{E}: G\hat M \ra \partial^{(2)}_\infty \Gamma_M \times \bR$ with compact fibers. 
	%Here $C = 4\delta$ if $\hat M$ is $\delta$-hyperbolic, and 
	The map obtained from composing with 
	a projection to the first factor is a map given by taking the endpoints of complete isometric geodesics. 
\end{proposition} 
\begin{proof} 
	Choose arbitrarily a base point $x_0$ in $\hat M$. 
	For each complete isometric geodesic $g: \bR \ra \hat M$, 
	let $g_{x_0}$ denote a projection point on $g(\bR)$ that is of 
	the minimal distance from $x_0$
	and let $\partial_+ g, \partial_-g \in \partial_\infty \Gamma_M$ denote
	the forward and backward endpoints,  
	and let $t(g) \in \bR$ denote $\pm d(g(0), g_{x_0})$ where we use $+$ 
	if $g(0)$ is ahead of $g_{x_0}$.  
	
	We define 
	$\mathcal{E}(g) = (\partial_+ g, \partial_- g, t(g))$. 
	The surjectivity follows from Proposition 2.1 of Chapter 2 of 
	\cite{CDP90} since $t$ is an isometry on each complete isometric  
	geodesic. 	The compactness of the fiber is implied by Lemma \ref{lem:twogeo}. 
	
	%	 The $C$-continuity of $\mathcal{E}$ composed with 
	%	 a projection to the $\bR$-factor in 
	%	 $\partial^{(2)}_\infty \hat M \times \bR$ 
	%	 follows from the fact that a closest point from 
	%	 $x_0$ to $g(\bR)$ is determined up to a distance $\leq 4\delta$
	%	 by Lemma 11.52 of \cite{DK2018}. 
	
	%	By Lemma \ref{lem:identify}, we use $\Uc M$ instead of $G\hat M$. 
	%	Given a sequence of elements $\vec{u}_i$ of $\Uc M$ converging to an 
	%	element $\vec{u}$ of $\Uc M$, the sequence 
	%	of corresponding geodesics $l_i$
	%	converges to a geodesic $l$ by Theorem 11.45 of \cite{DK2018}
	%	in every compact subset of $\hat M$. 
	%	We will use the Gromov boundary with 
	%	the ideal boundary by Theorem 11.104 of \cite{DK2018}. 
	%	Let $x \in l$ be a point projected from 
	%	$\vec{u}$. Let $y_i \in l_i$ be a point  
	%	projected from $\vec{u}_i$ and the forward endpoint 
	%	$z_i \in \partial_\infty \hat M$. 
	
	
	Champetier \cite{Champetier94} constructs the space 
	$G\hat{\Gamma}_M$ from $\hat{\Gamma}_M$ quasi-isometric to $\Gamma_M$ 
	following Gromov. 
	% 	By Proposition	\ref{prop:Ucdense}, 
	% 	for each element of $\hat M$, 
	% 	there is a complete isometric geodesic within 
	% 	a distance of $C$.  
	%Since $\hat M$ is quasi-isometric with $\Gamma_M$, 
%	Let $x_0$ denote a base point of $\hat M$. 
Recall from Section \ref{sub:Gromov} that $\partial_\infty \hat{M} = \partial_\infty \Gamma_M$. Using the orbit map $\Gamma_M \ra \hat{M}$ sending $\gamma$ to $\gamma(x_0)$, we extend the map to $\hat{\Gamma}_M \ra \hat{M}$ by sending each edge to a shortest geodesic.

Let $I: G\hat{\Gamma}_M \ra G\hat{M}$ be a map that sends a geodesic $g$ in $\hat{\Gamma}_M$ to a complete isometric geodesic $g'$ in $\hat{M}$ with the same endpoints in $\partial_\infty \hat{M}$, and maps $g(0)$ to the nearest point on the image of $g'$.

The map $I': G\hat{M} \ra G\hat{\Gamma}_M$ can be defined by taking a complete isometric geodesic $g$ in $\hat{M}$ to a geodesic $g'$ in $\hat{\Gamma}_M$ in a similar manner, with $g(0)$ going to one of the elements of $\Gamma_M(x_0)$ nearest to it.

By Theorem 11.72 of \cite{DK2018}, which states that every isometric geodesic in $\hat{\Gamma}_M$ is a quasi-geodesic in $\hat{M}$ and vice versa, and considering that every complete isometric geodesic in $\hat{M}$ is uniformly bounded away from one in $\hat{\Gamma}_M$ in the Hausdorff distance, it follows that $G\hat{M}$ is quasi-isometric to $G\hat{\Gamma}_M$ via the maps $I$ and $I'$, using \eqref{eqn:dGX}.

%	Recall $\partial_\infty \hat M = \partial_\infty \hat \Gamma_M$. 
%   Using the orbit map $\Gamma_M \ra \hat M$ sending $\gamma$ to $\gamma(x_0)$,
%   we extend the map to $\hat \Gamma_M \ra \hat M$ by sending each 
%   edge to a shortest geodesic. 
%	Let $I: G\hat \Gamma_M \ra G\hat M$ be a map 
%	given by sending a geodesic $g$ in $\hat \Gamma_M$ 
%	to a geodesic $g'$ in $\hat M$ with the same endpoints in
%	$\partial_\infty \hat M$
%	with $g(0)$ to a nearest point to the image of $g'$. 
%	The map $I': G\hat M \ra G\hat \Gamma_M$ can be defined 
%	by taking the geodesic $g$ of $\hat M$ 
%	to a geodesic $g'$  of $\hat \Gamma_M$ 
%	in the same manner with $g(0)$ going to one of the elements of
%	$\Gamma_M(x_0)$ nearest to it. 
%	Since a geodesic in $\hat \Gamma_M$ is a quasi-geodesic in 
%	$\hat M$, and every geodesic in $\hat M$ is uniformly bounded 
%	away from one in $\hat \Gamma_M$ in the Hausdorff distance
%	by Theorem 11.72 of \cite{DK2018}, 
%	%Proposition \ref{prop:GtildeM} implies that  
%	$G\hat M$ is quasi-isometric with 
%	$G\hat \Gamma_M$ by the maps $I$ and $I'$ 
%	using \eqref{eqn:dGX}.  
	
	%Since $\hat M$ is quasi-isometric to $\Gamma_M$, 
	%$\hat \Gamma_M$ is quasi-isometric to $G\hat M$ 
	%by our construction. 
	%shows that $\hat \Gamma_M$ is quasi-isometric to 
	%$\Gamma_M$. 
	% 	Since $\Uc \hat M$ is quasi-isometric 
	% 	to $\hat M$ by Proposition \ref{prop:Ucdense}, 
	% 	$G\hat M$ is quasi-isometric to $\hat \Gamma_M$. 
	We define 
$\mathcal{E}': G\hat{\Gamma}_M \ra \partial_\infty^{(2)} \Gamma_M \times \bR$ 
	as we did for $\mathcal{E}$ above. 
	By Proposition 4.8 and (4.3) of \cite{Champetier94}, 
	$\mathcal{E}'$ sends 
	$G\hat{\Gamma}_M$ to 
	$\partial_\infty^{(2)} \Gamma_M \times \bR
	= \partial_\infty^{(2)} \hat{M}\times \bR$ as a quasi-isometric homeomorphism. 
	Since $\mathcal{E}(g, t)$ and $\mathcal{E}'\circ I'(g, t)$ are uniformly bounded away from 
	each other,  the result follows. 
	%
	%	Using the shadow topology of the ideal boundary 
	%as in Section 11.11 of \cite{DK2018}, we obtain that 
	%
	%	 $\mathcal{E}$ composed with the projection 
	%	 to the $\partial^{(2)}$-factor is continuous. 
	%
	%	The map $\mathcal{E}$ is given by taking endpoints of a complete isometric geodesics. 
	%	The map is surjective by Proposition 2.1 of Chapter 2 of \cite{CDP90}. 
	%The compactness of the fiber follows by Lemma \ref{lem:twogeo}. 
	%by the Morse lemma as written in 
	%Proposition 3.1 of Chapter 3 of \cite{CDP90}. 
	%	See \cite{GLM09} for connectedness with anologous proof. 
	%\qed
\end{proof}

%\marginpar{This quasi-isometry clear from $G_M$? Check... }

The following shows that $\Uc\hat{M}$ can be 
used as our Gromov flow space. 

\begin{theorem}\label{thm:identify} 
	%	There is a metric on $\partial^{(2)} M \times \bR$ so that 
	The map	
	$\mathcal{E}\circ \mathcal{F}:\Uc \hat M \ra \partial_\infty^{(2)} \Gamma_M \times \bR$ 
	is a quasi-isometry where each complete isometric 
	geodesic in $\Uc \hat M$ goes to 
	$(t_1,t_2)\times \bR$ for its pair $(t_1, t_2)\in \partial_\infty^{(2)} \Gamma_M$ 
    of endpoints. 
\end{theorem} 
\begin{proof} 
	Proposition \ref{prop:Uc} and Lemma \ref{lem:identify} imply
	the result. %\qed
	% 	We construct a metric on  $\partial^{(2)} M \times \bR$ the space 
	% 	as follows: 
	% 	For each pair of complete isometric geodesics going to the same 
	% 	sets of form $(t_1, t_2)\times \bR$, we use one of 
	% 	the metric restricted to it. 
	% 	For pairs of points $\partial^{(2)} M \times \bR$
	% 	in different factors of $\bR$, 
	% 	we take the inverse images of $\mathcal{E}\circ \mathcal{F}$ of 
	% 	the pair of points and taking the maximal distance in $\Uc \hat M$. 
	% 	The map restrict to a complete isometric geodesic to an $\bR$-factor is an
	% 	isometry by construction in the proof of Proposition \ref{prop:Uc}. 
\end{proof} 




%%% Sept 27 6:00pm 2019
%\part{title}
%For each element $l$ of $\partial^{(2)}_G \hat M$, i.e, the lift 
%$\vec{l}$ of an oriented complete 
%geodesic $l$ in $\Uc \hat M$, 
%the subspace $V_+(y), V_0(y), V_-(y)$ 
%are independent of points of $\Uc \hat M$ on the line $\vec{l}$. 
%We obtain three subspaces $V_+(\vec{l}), V_0(\vec{l}), V_-(\vec{l})$.  

%\subsection{Flow space $G\hat M$ quasi-isometric to $\partial_\infty^{(2)}\hat M \times \bR$}

%	We give the space 
%	$\partial_\infty^{(2)}\hat M \times \bR$ the metric as in the paragraph after Theorem \ref{thm:Mineyev}. 

 
% \begin{proposition} \label{prop:Uc} 
% 	Let $M$ be a closed aspherical manifold.
% 	Suppose that $\Gamma_M$ is word-hyperbolic.
% 	Then there is 
% 	a quasi-isometric surjective map 
% 	$\mathcal{E}: G\hat M \ra \partial^{(2)}_\infty \hat M \times \bR$ whose fibers are compact. 
% 	%Here $C = 4\delta$ if $\hat M$ is $\delta$-hyperbolic, and 
% 	The map obtained from composing with 
% 	a projection to the first factor is a map given by taking the endpoints of complete isometric geodesics. 
% \end{proposition} 
% \begin{proof} 
% 	Choose a basepoint $x_0$ in $\hat M$. 
% 	For each complete isometric geodesic $g: \bR \ra \hat M$, 
% 	let $g_{x_0}$ denote the projection point on $g(\bR)$ that is of minimal distance from $x_0$
% 	and let $\partial_+ g, \partial_-g \in \partial_\infty \hat M$ denote
% 	the forward and backward endpoints,  
% 	and let $t(g) \in \bR$ denote $\pm d(g(0), g_{x_0})$ where we use $+$ 
% 	if $g(0)$ is ahead of $g_{x_0}$.  
% 	
% 	We define 
% 	$\mathcal{E}(g) = (\partial_+ g, \partial_- g, t(g))$. 
% 	The surjectivity follows from Proposition 2.1 of Chapter 2 of 
% 	\cite{CDP90} since $t$ is a proper map on each complete isometric  
% 	geodesic. 	The compactness of the fiber follows by Lemma \ref{lem:twogeo}. 
% 	
% 	%	 The $C$-continuity of $\mathcal{E}$ composed with 
% 	%	 a projection to the $\bR$-factor in 
% 	%	 $\partial^{(2)}_\infty \hat M \times \bR$ 
% 	%	 follows from the fact that a closest point from 
% 	%	 $x_0$ to $g(\bR)$ is determined up to a distance $\leq %4\delta$
% 	%	 by Lemma 11.52 of \cite{DK2018}. 
% 	
% 	%	By Lemma \ref{lem:identify}, we use $\Uc M$ instead of $G\hat M$. 
% 	%	Given a sequence of elements $\vec{u}_i$ of $\Uc M$ converging to an 
% 	%	element $\vec{u}$ of $\Uc M$, the sequence 
% 	%	of corresponding geodesics $l_i$
% 	%	converges to a geodesic $l$ by Theorem 11.45 of \cite{DK2018}
% 	%	in every compact subset of $\hat M$. 
% 	%	We will use the Gromov boundary with 
% 	%	the ideal boundary by Theorem 11.104 of \cite{DK2018}. 
% 	%	Let $x \in l$ be a point projected from 
% 	%	$\vec{u}$. Let $y_i \in l_i$ be a point  
% 	%	projected from $\vec{u}_i$ and the forward endpoint 
% 	%	$z_i \in \partial_\infty \hat M$. 
% 	
% 	
% 	Champetier \cite{Champetier94} constructs the space 
% 	$G\hat \Gamma_M$ from $\hat \Gamma_M$ quasi-isometric to %$\Gamma_M$ 
% 	following Gromov. 
%% 	By Proposition	\ref{prop:Ucdense}, 
%% 	for each element of $\hat M$, 
%% 	there is a complete isometric geodesic within 
%% 	a distance of $C$.  
%%Since $\hat M$ is quasi-isometric with $\Gamma_M$, 
%Let $x_0$ denote a base point of $\hat M$. 
%Recall $\partial_\infty \hat M = \partial_\infty \hat \Gamma_M$. 
%Let $I: G\hat \Gamma_M \ra G\hat M$ be a map 
%given by sending a geodesic $g$ in $\hat \Gamma_M$ 
%to a geodesic $g'$ in $\hat M$ with the same endpoints in
%$\partial_\infty \hat M$
%with $g(0)$ to a nearest point to the image of $g'$. 
%The map $I': G\hat M \ra G\hat \Gamma_M$ can be defined 
%by taking the geodesic $g$ of $\hat M$ 
%to a geodesic $g'$  of $\hat \Gamma_M$ 
%in the same manner with $g(0)$ going to one of the element of
%$\Gamma_M(x_0)$ nearest to it. 
%Since geodesic in $\hat \Gamma_M$ is a quasi-geodesic in 
%$\hat M$, and every geodesic in $\hat M$ is uniformly bounded 
%away from one in $\hat \Gamma_M$ in the Hausdorff distance
%by Theorem 11.72 of \cite{DK2018}, 
%%Proposition \ref{prop:GtildeM} implies that  
% $G\hat M$ is quasi-isometric with 
% 	$G\hat \Gamma_M$ by the maps $I$ and $I'$ 
% 	using \eqref{eqn:dGX}.  
% 	
% 	%Since $\hat M$ is quasi-isometric to $\Gamma_M$, 
% 	%$\hat \Gamma_M$ is quasi-isometric to $G\hat M$ 
% 	%by our construction. 
% 	%shows that $\hat \Gamma_M$ is quasi-isometric to 
% 	%$\Gamma_M$. 
%% 	Since $\Uc \hat M$ is quasi-isometric 
%% 	to $\hat M$ by Proposition \ref{prop:Ucdense}, 
%% 	$G\hat M$ is quasi-isometric to $\hat \Gamma_M$. 
% We define $\mathcal{E}': G\hat \Gamma_M \ra \partial_\infty^{(2)} %\hat \Gamma_M \times \bR$ 
% as above. 
% 	By Proposition 4.8 of \cite{Champetier94}, 
% 	$\mathcal{E}'$ sends 
% 	$G\hat \Gamma_M$ to 
% 	$\partial_\infty^{(2)} \hat \Gamma_M \times \bR
% 	= \partial_\infty^{(2)}\hat M\times \bR$ as a quasi-isometric %homeomorphism. 
% Since $\mathcal{E}= \mathcal{E}'\circ I'$, the result follows. 
% 	
% 	%	Using the shadow topology of the ideal boundary 
% 	%as in Section 11.11 of \cite{DK2018}, we obtain that 
% 	
% 	%	 $\mathcal{E}$ composed with the projection 
% 	%	 to the $\partial^{(2)}$-factor is continuous. 
% 	
% 	%	The map $\mathcal{E}$ is given by taking endpoints of a %complete isometric geodesics. 
% 	%	The map is surjective by Proposition 2.1 of Chapter 2 of %\cite{CDP90}. 
% 	%The compactness of the fiber follows by Lemma \ref{lem:twogeo}. 
% 	%by the Morse lemma as written in 
% 	%Proposition 3.1 of Chapter 3 of \cite{CDP90}. 
% 	%	See \cite{GLM09} for connectedness with anologous proof. 
% 	
% \end{proof}
%
%\marginpar{This quasi-isometry clear from $\Gamma_M$? Check... }
% 
% The following shows that $\Uc\hat M$ can be 
% used as our Gromov-Mineyev flow space up to some collapsing. 
% 
% \begin{theorem}\label{thm:identify} 
% %	There is a metric on $\partial^{(2)} M \times \bR$ so that 
% The map	
% $\mathcal{E}\circ \mathcal{F}:\Uc \hat M \ra \partial^{(2)} M \times \bR$ 
% 	is a quasi-isometry so that each complete isometric 
% 	geodesic in $\Uc \hat M$ goes to 
% 	$(t_1,t_2)\times \bR$ for its endpoint pair $(t_1, t_2)\in \partial^{(2)}_\infty \hat M$ as an isometry. 
% \end{theorem} 
% \begin{proof} 
%Proposition \ref{prop:Uc} and Lemma \ref{lem:identify} imply
%the result. 
%% 	We construct a metric on  $\partial^{(2)} M \times \bR$ the space 
%% 	as follows: 
%% 	For each pair of complete isometric geodesics going to the same 
%% 	sets of form $(t_1, t_2)\times \bR$, we use one of 
%% 	the metric restricted to it. 
%% 	For pairs of points $\partial^{(2)} M \times \bR$
%% 	in different factors of $\bR$, 
%% 	we take the inverse images of $\mathcal{E}\circ \mathcal{F}$ of 
%% 	the pair of points and taking the maximal distance in $\Uc \hat M$. 
%% 	The map restrict to a complete isometric geodesic to an $\bR$-factor is an
%% 	isometry by construction in the proof of Proposition %\ref{prop:Uc}. 
% \end{proof} 
% 

 



\section{Partial hyperbolicity and $P$-Anosov properties} \label{sec:phyp} 

%From now on, we shift our attention to $P$-Anosov properties. 

\reD{The purpose of this section is to set notations and recall standard facts since some of these are not too well-known to many geometric topologists not in this specific area.}

\subsection{Some background} \label{sub:background} 
We have to repeat some of the standard materials from
Section 2.2 of \cite{GGKW17ii} and Chapter 2 of \cite{GJT}. 
Consider a real reductive Lie group $G$ with a maximal compact subgroup $K$. 

We assume that $\mathrm{Ad}(G) \subset \Aut(\mathfrak{g})$ for simplicity. 
For example, this holds if $G$ is connected. We can usually reduce everything here to 
the case where $G$ is connected. 
Here, $G$ is an almost product of $Z(G)_0$ and $G_s$ 
where $Z(G)_0$ is the identity component of the center $Z(G)$ of $G$ 
and $G_s = D(G)$ the derived subgroup of $G$, which is semisimple. 

A parabolic subgroup $P$ of $G$ is a subgroup $P$ of the form 
$G \cap \mathbf{P}(\bR)$ for some algebraic subgroup $\mathbf{P}$ of an algebraic group $\mathbf{G}$ 
where $\mathbf{G}(\bR)/\mathbf{P}(\bR)$ is compact
and $\mathbf{G}(\bR) = G$. 

Let $\mathfrak{z}(\mathfrak{g})$ denote the Lie algebra of $Z(G)$, 
and $\mathfrak{g}_s$ the Lie algebra of $G_s$. 
Then $\mathfrak{g} = \mathfrak{z}(\mathfrak{g}) \oplus \mathfrak{g}_s$.
Letting $\mathfrak{t}$ denote the Lie algebra of $K$, 
we have $\mathfrak{g} = \mathfrak{t} \oplus \mathfrak{q}$ 
where $\mathfrak{q}$ is the $-1$-eigenspace of an involution $\theta$ preserving 
$\mathfrak{t}$ as the $1$-eigenspace.  
Let $\mathfrak{a} \subset \mathfrak{q}$ be the maximal abelian subalgebra, i.e., the Cartan subalgebra of 
$\mathfrak{g}$. Then $\mathfrak{a} = \mathfrak{z}(\mathfrak{g}) \cap \mathfrak{q} \oplus \mathfrak{a}_s$ 
where $\mathfrak{a}_s = \mathfrak{a} \cap \mathfrak{g}_s$, a maximal \green{abelian} subalgebra of 
$\mathfrak{q} \cap \mathfrak{g}_s$. 

We have a decomposition to $ad(\mathfrak{a})$-eigenspaces
 $\mathfrak{g}=\mathfrak{g}_{0} \oplus \bigoplus_{\alpha \in \Sigma} \mathfrak{g}_{\alpha}$.
 Here, 
 $\mathfrak{g}_{0}$ is the direct sum of $z(\mathfrak{g})$ and the centralizer of $\mathfrak{a}$ in $\mathfrak{g}_{s} $. 
 The set $\Sigma \subset \mathfrak{a}^{*}= \Hom_{\bR}(\mathfrak{a}, \bR)$ projects to a (possibly nonreduced) root system of $\mathfrak{a}_{s}^{*},$ and each $\alpha \in \Sigma$ is called a restricted root of $\mathfrak{a}$ in $\mathfrak{g}$.
 There is a subset 
$\Delta \subset \Sigma$ called a {\em simple system},  i.e. a subset such that any root is expressed 
uniquely as a linear combination of elements of $\Delta$ with coefficients all of the same sign; 
the elements of $\Delta$ are called simple roots.
There is a subset
$\Sigma^{+} \subset \Sigma$ which is the set of positive roots, i.e. 
 roots that are nonnegative linear combinations of elements of $\Delta$, where 
 $\Sigma=\Sigma^{+} \cup\left(-\Sigma^{+}\right)$ holds additionally. 

We define $P_\theta$ as the parabolic subgroup with the Lie algebra
\[\operatorname{Lie}\left(P_{\theta}\right)=\mathfrak{g}_{0} \oplus \bigoplus_{\alpha \in \Sigma^{+}} \mathfrak{g}_{\alpha} \oplus \bigoplus_{\alpha \in \Sigma^{+} \backslash \Sigma_{\theta}^{+}} \mathfrak{g}_{-\alpha}.\]

%Recall the Cartan decomposition $\mathfrak{g} = \mathfrak{t} \oplus 
%\mathfrak{p}$ and $\theta$ the associated Cartan involution. 
%Let $\mathfrak{a}$ denote a maximal abelian subaglebra of $\mathfrak{p}$. 
%Let $\Sigma$ denote the set of roots of $\mathfrak{g}$.
%Let $A$ denote the  subalgebra of $\sl(n, \bR)$. 
%The Lie algebra 
%\[  \mathfrak{g} = \mathfrak{g}_0 + \sum_{\alpha \in \Sigma} \mathfrak{g}_{\alpha} \]
%where 
%\[ \mathfrak{g}_\alpha := \{X \in \mathfrak{g}| [H, X] = \alpha(H) X 
%\hbox{ for all } H \in \alpha\}.\]
%And $\mathfrak{g}_0 = \mathfrak{m} + \mathfrak{a}$
%where $\mathfrak{m} $ is the centralizer in $\mathfrak{t}$ of $\mathfrak{a}$. 
The connected components of $\mathfrak{a} \setminus \cup_{\alpha \in \Sigma} \ker(\alpha)$ 
are called {\em Weyl chambers} of $\mathfrak{a}$. 
The component where every $\alpha \in \Sigma_+$ is positive is denoted by $\mathfrak{a}^+$. 

%Let $\Sigma^+$ denote the set of positive roots. 
%Let $\Delta = \Delta(\mathfrak{g}, \mathfrak{a}^+)$ denote the set of simple roots. 
%Let $B$ denote the Killing form of $\mathfrak{p}$. 
%We define $H_\alpha$ to be the vector in $\mathfrak{a}$
%so that $\alpha(H) = B(H, H_\alpha)$ for $\alpha \in \Sigma$. 
%Each subset $I$ of $\Delta$ determines an orthogonal decomposition 
%$\mathfrak{a} = \mathfrak{a}^I \oplus \mathfrak{a}_I$
%where $\mathfrak{a}^I$ is the span of $H_{\alpha}, \alpha \in I$ 
%and $\mathfrak{a}_I$ the orthogonal complement of $\mathfrak{a}^I$
%in $\mathfrak{a}$. 

%We denote by $\Sigma^I$ the subset of $\Sigma$ vanishing on $\mathfrak{a}_I$. 
%Let $\Sigma_I := \Sigma - \Sigma^I$. 
%Let $\Sigma^{I, +}$ denote $\Sigma^I \cap \Sigma^+$
%and $\Sigma_{I, +}$ denote $\Sigma_I \cap \Sigma^+$.  
%We define \[\mathfrak{n} := \sum_{\alpha \in \Sigma^+} \mathfrak{g}_{\alpha}, 
%\mathfrak{n}^I := \sum_{\alpha \in \Sigma^{I, +}} \mathfrak{g}_{\alpha}, 
%\mathfrak{n}_I := \sum_{\alpha \in \Sigma^+_I} \mathfrak{g}_{\alpha}.\]
%Let $N$ denote $\exp(\mathfrak{n})$, and we define  
%$N^I := \exp(\mathfrak{n}^I)$, and $N_I := \exp(\mathfrak{n}_I)$. 
%
%We denote $A$ as $\exp \mathfrak{a}$, $N$ the nilpotent subgroup with Lie algebra $\mathfrak{n}$, 
%and $M$ the centralizer of $\mathfrak{a}$ in $K$. 
%
%
%If $I$ is a subset of $\Delta$, let $P^I$ denote the normalizer in $G$ of 
%$\mathfrak{n}_I$. 

%\begin{theorem}[Theorem 2.8 \cite{GJT}] \label{thm:GJT} 
%	$P^I$ is a standard parabolic subgroup. Every standard parabolic subgroup $P'$ is 
%	of this form. 
%	The Lie algebra of $P^I$ is 
%	\[\mathfrak{b}_I = \mathfrak{b} + \sum_{\alpha \in \Sigma^I} \mathfrak{g}_{\alpha} = \mathfrak{b} + \bar{\mathfrak{n}}^I,\] 
%	where $\mathfrak{b}$ is the Lie algebra of $P = MAN$
%	and $\bar{\mathfrak{n}}^I = \theta(\mathfrak{n}^I)$.
%	Every Lie algebra containing $\mathfrak{b}$ is of this form. 
%	
%	$P^I$ is the unique closed subgroup of $G$ containing $P$ with the Lie algebra 
%	$\mathfrak{b}_I$. In particular, $P^{I_1} \subset P^{I_2}$ if and only if 
%	$I_1 \subset I_2$. 
%\end{theorem} 

Let $G = \SL_\pm(n, \bR)$. 
The maximal abelian subalgebra 
$A_n$ of $\mathfrak{p}$ in this case
is the subspace of diagonal matrices with 
entries $a_1, \dots, a_n$ where $a_1 + \dots + a_n = 0$. 
Let $\log \lambda_i: A_n\ra \bR$ denote the projection to the $i$-th factor. 
$\Sigma$ consists of $\alpha_{ij} :=\log \lambda_i -\log \lambda_j$ for $i\ne j$, 
and \purple{the root space} $\mathfrak{sl}(n, \bR)_{\alpha_{ij}} = \{c E_{ij}| c\in \bR\}$ where $E_{ij}$ is a matrix with \purple{an} entry $1$ at $(i, j)$ and zero entries everywhere else. 

A positive Weyl chamber is given by
\[A_n^{+} =\{(a_1, \dots, a_n)| a_1\geq \dots \geq a_n, a_1+\cdots + a_n = 0\},\]
and we are given 
\begin{gather*} 
\Sigma^+ = \{\log \lambda_i -\log \lambda_j | i < j\} \hbox{ and }\\ 
\Delta = \{\alpha_i:= \log \lambda_i -\log \lambda_{i+1} | i=1, \dots, n-1\}.
\end{gather*}  

We recall Example 2.14 of \cite{GGKW17ii}: 
Let $G$ be $\GL(n, \bR)$ seen as a real Lie group. Its derived group is $G_{s}=D(G)=\SL(n, \bR) .$ 
We can take %$K$ to be $\mathrm{O}(n)$ and 
$\mathfrak{a} \subset \mathfrak{g l}(n, \bR)$ to 
be the set of real diagonal matrices of size $n \times n$. 
As above, for $1 \leq i \leq n,$ let $\log \lambda_{i} \in \mathfrak{a}^{*}$ 
be the evaluation of the $i$-th diagonal entry. 
Then 
\[\mathfrak{a} =\mathfrak{z}(\mathfrak{g}) \cap \mathfrak{a} \oplus \mathfrak{a}_{s} 
\hbox{ where }
\mathfrak{z}(\mathfrak{g}) \cap \mathfrak{a}=\bigcap_{1 \leq i, j \leq n}  \operatorname{Ker}\left(\log \lambda_{i}-\log \lambda_{j}\right)\] 
is the set of real scalar matrices, and $\mathfrak{a}_{s}=
\operatorname{Ker}\left(\log \lambda_{1}+\cdots+\log \lambda_{n}\right)$ 
is the set of traceless real diagonal matrices. The set of restricted roots of $\mathfrak{a}$ in $G$ is
$$
\Sigma=\left\{\log \lambda_{i}-\log \lambda_{j} \mid 1 \leq i \neq j \leq n\right\}.
$$
We can take $\Delta=\left\{\log \lambda_{i}-\log \lambda_{i+1} \mid 1 \leq i \leq n-1\right\}$ so that
$$
\Sigma^{+}=\left\{\log \lambda_{i}-\log \lambda_{j} \mid 1 \leq i<j \leq n\right\},
$$
and $\mathfrak{a}^{+}$ is the set of elements of $\mathfrak{a}$ whose entries are increasing. 
Let $\clo({\mathfrak{a}}^+)$ denote its closure. 
For 
\begin{multline} 
\theta=\left\{\log \lambda_{i_{1}}-\log \lambda_{i_{1}+1}, \dots, \log \lambda_{i_{m}}-\log \lambda_{i_{m}+1}\right\} \\ 
\hbox{ with } 1 \leq i_{1}<\cdots<i_{m} \leq n-1,
\end{multline} 
the parabolic
subgroup $P_{\theta}$ (resp. $\left.P_{\theta}^{-}\right)$ is the group of block upper (
resp. lower) triangular matrices in $\GL(n,\bR)$ with square diagonal blocks 
of sizes $i_{1}, i_{2}-i_{1}, \dots, i_{m}-i_{m-1}, n-i_{m}$. 
In particular, $P_{\Delta}$ is the group of upper triangular matrices in $\GL(n, \bR)$.

From now on, we will fix $\mathfrak{a}_n \subset \GL(n, \bR)$ 
as the set of real diagonal matrices of size $n \times n$. 
We will also regard $\theta$ as a subset of $\{1, \dots, n-1\}$ where 
each $i$ corresponds to $\log \lambda_i - \log \lambda_{i+1}$. 
Also, let $\mathfrak{a}_n^+ $ denote the open positive Weyl chamber. 


\subsection{$P$-Anosov representations}  \label{sub:Panosov} 
Let $(P^{+}, P^{-})$ be a pair of opposite parabolic subgroups of $\GL(n, \bR)$, and 
let $\mathcal{F}^{{\pm}}$ denote the flag spaces $\GL(n, \bR)/P_{\pm}$.
Let $\chi$ denote the unique open $\GL(n, \bR)$-orbit of the product 
$\mathcal{F}^{+}\times \mathcal{F}^{-}$. 
The product subspace $\chi$ has two $\GL(n, \bR)$-invariant 
distributions $E^{\pm} := T_{x_{\pm}} \mathcal{F}^{\pm}$ for $(x^{+}, x^{-}) \in \chi$. 
\purple{Suppose that $\Gamma_M$ is a finitely generated word-hyperbolic group. }
Let $\phi_t$ be given as in Theorem \ref{thm:Mineyev} as the $\bR$-action on 
$F \Gamma_M$  whose orbits are quasi-geodesics. 
We denote by $(\tau_+, \tau_-): F\Gamma_M \ra F\Gamma_M/\bR \cong \partial_\infty \Gamma_M^{(2)}$ the maps associating to a point the \green{endpoints} of its $\bR$-orbit.

%Recall the composition $f$ of maps in Lemma \ref{lem:identify}. 
%By Theorem \ref{thm:identify}, 
%$\partial_\infty^{(2)} \hat M \times \bR$ is a quotient 
%metric space of $\Uc \hat M$ by a quasi-isometry. 

\begin{definition}[Definition 2.10 of \cite{GW12}]\label{defn:Anosov} 
\purple{Suppose that $\Gamma_M$ is a finitely generated word-hyperbolic group. }
	A representation $\rho: \Gamma_M \ra \GL(n, \bR)$ is {\em $(P^{+}, P^{-})$-Anosov} if 
	there exist continuous $\rho$-equivariant maps $\xi^+: \partial_\infty \Gamma_M \ra {\mathcal{F}}^+$, 
	$\xi^-:\partial_\infty \Gamma_M \ra {\mathcal{F}}^-$ such that:  
	\begin{enumerate} 
		\item[(i)] for all $(x, y)\in \partial_\infty^{(2)} \Gamma_M$, the pair 
		$(\xi^+(x), \xi^-(y))$ is transverse, and  
		\item[(ii)] for one (and hence any) continuous and equivariant family of norms 
		$(\llrrV{\cdot}_{m})_{m\in F\Gamma_M} $ on 
		\[(T_{\xi^+(\tau^+(m))} {\mathcal{F}}^+)_{m\in F\Gamma_M}
		\Big(\hbox{resp. } 
		(T_{\xi^-(\tau^-(m))} {\mathcal{F}}^-)_{m\in F\Gamma_M}\Big):\]
		\[ \llrrV{e}_{\phi_{-t}m} \leq A e^{-at} \llrrV{e}_{m}  \Big(\hbox{resp.} 
		\llrrV{e}_{\phi_{t}m} \leq A e^{-at} \llrrV{e}_{m} \Big).
		   \]
	\end{enumerate}
We call 	$\xi^\pm:\partial_\infty \green{\Gamma_M} \ra {\mathcal{F}}^\pm$ the {\em Anosov maps} associated with 
$\rho: \Gamma_M \ra G$. 
	(See also Section 2.5 of  \cite{GGKW17ii} for more details.)
\end{definition} 



It is sufficient to consider only the case where $P^{+}$ is conjugate to $P^{-}$
by Lemma 3.18 of \cite{GW12}. (See also Definition 5.62 of \cite{KLP17}). 
In this case, $\rho$ is called {\em $P$-Anosov} for $P = P^+$. 
%Hence, the flag space $\mathcal{F}^{+}$ consists of flags of form 
%$F_{1}\subset F_{2}\subset \dots \subset F_{m}$ denoting $F_{m+1}=\bR^{n}$ where 
%$\dim F_{i}= f_{i}$ where $0< f_1 < f_2 < \dots < f_m < n$. 
%%satisfying $d_{i} = n- d_{m+1-i}$. 
%%We assume that $m \geq 3$. We will not study the case 
Note that this means that 
\begin{equation}  \label{eqn:theta*} 
\theta = \theta^*
\end{equation} 
where $\theta^*$ is the image $\theta$ under the map $i \mapsto n-i$, $i=1, \dots, n-1$.
Hence, this will always be true for the set $\theta$ occurring for $P_\theta$-Anosov 
representations below.
(See the paragraph before Fact 2.34 of \cite{GGKW17ii}.)


%%Since $\hat M$ is negatively curved, it has a well-defined ideal boundary 
%%$\partial_\infty \hat M$ (see Eberlein \cite{Eberlein}). 
%%We denote by $\partial_\infty^{(2)} \hat M = \partial_\infty \hat M \times \partial_\infty \hat M - \Delta$ 
%%where $\Delta$ is the diagonal set. The space is homeomorphic to the space of 
%%complete geodesics in $\hat M$. 

%A flag space is given by a strictly increasing sequence $F=\{ F_{1}\subset F_{2} \subset \cdots \subset F_{m}\}$ of 
%nonzero proper subspaces of $\bR^{n}$. The dimension $f=\{f_{1}, f_{2}, \dots, f_{m}\}$ is the invariant of $F$. 
%where $f_{j}= \dim F_{j}$ for $j=1, \dots, m$. 
%Let $P_{F}$ be the subgroup of $\SL_{\pm}(n, \bR)$ preserving $F$. The conjugacy class of $P_{F}$ is uniquely determined by the sequence $f$. 



%The parabolic group $P^{+}$ is of such a form. The opposite one $P^{-}$ can be considered a 
%stabilizer of the dual flags $F^{\ast} =\{ F^{\ast}_{1}\supset F^{\ast}_{2} \supset \cdots \supset F^{\ast}_{m}\}$ for 
%$F^{\ast}_i\subset \bR^{n\ast}$ where 
%$F^{\ast}_{j}$ is the nullifier of $F_{j}$ for each $j$. 
%%We assume that $P^+$ and $P^-$ are conjugate to a parabolic subgroup $P$. 

%We define $m_{P}$ to be the number of proper subspaces in the flag $F_{P}$ stabilized by $P$. 
%Let $f_{P}=\{ f_{1, P}, \dots, f_{m_P, P}\}$ denote the flag dimension of $F_{P}$. 
%A pair of points $(x^{+}, x^{-}) \in \mathcal{F}^{+}\times \mathcal{F}^{-}$ is {\em transverse} if $(x^{+}, x^{-}) \in \chi$. 
%%According to the definition \cite{GW12}, there exist continuous functions 
%%$\xi^{\pm}: \partial_{\infty} \hat M \ra \mathcal{F}^{\pm}$ with following properties: 
%%\begin{enumerate} 
%%	\item for each $(t^{+}, t^{-}) \in \partial_{\infty}^{(2)} \hat M$, the pair 
%%	$\xi^{+}(t^{+})$ and $\xi^{-}(t^{-})$ are transverse. 
%%	
%%	\item For a continuous and equivariant family of norms $(\llrrV{\cdot}_{\hat m \in %\hat M})$ on 
%%	the fibers of $\sigma^{\ast} E^{+} \ra M$ and $\sigma^{\ast} E^{-}$ and a positive %constant $A, a$ such that 
%%	for any $t > 0, e^{\pm} \in \sigma^{\ast} E^{\pm}$ with $\pi_\rho(e^{\pm}) = m$ one %has 
%%	$\llrrV{\phi_{-t}^\ast e^{+}}_{\phi_{-t} m} \leq A e^{-at} \llrrV{e^{+}_{m}}$ 
%%	and $\llrrV{\phi_{t}^\ast e^{-}}_{\phi_{t} m} \leq A e^{-at} \llrrV{e^{-}_{m}}$.
%%\end{enumerate} 
%
%Assume that $\rho$ is a $(P^{+}, P^{-})$-Anosov representation, and let 
%$\xi^{\pm}: \partial_{\infty} \Gamma_M \ra {\mathcal{F}^{\pm}}$ be the associated Anosov %map. 
%Let $\gamma$ be a nontorsion element. 
%Let $t^{\pm}_\gamma$ denote the attracting and repelling fixed point of $\gamma$ in %$\partial_{\infty} \Gamma_M$ respectively. 
%Then $\xi^{+}(t^{+}_{\gamma})$ is the unique attracting fixed point of $\rho(\gamma)$ on %$\mathcal{F}^{+}$, and 
%the basin of attraction is the set of points of $\mathcal{F}^{+}$ transverse to $\xi^{-}(t^{-}_{\gamma})$. 
%(See Section 2 of \cite{GW12} for more details)

%\subsection{$P$-Anosov representations for $\GL(n, \bR)$.}  
%\label{sub:PI} 

%Here, we will discuss on various definitions involved in defining 
%$P$-Anosov property using singular values. 
%We will not give the precise definition of the $P$-Anosov property due to  Kapovich-Leeb-Porti 
%because of the length but refer to Definition 5.43 of \cite{KLP17} and various equivalent 
%properties in Theorem 1.1 of \cite{KLP17}. 
%We will use the notation of \cite{KLP17} here.
%Let $\sigma_{mod}$ denote the spherical Weyl chamber of the semisimple Lie group $G$. Recall the notion of opposition involution $\iota: \sigma_{mod} \ra \sigma_{mod}$. (See Kapovich-Lee-Porti \cite{KLP18} and Benoist \cite{Benoist97}.)

%Proposition 5.61 of 
%Appendix 5.11 of \cite{KLP18} explains the equivalence of 
%the definition of Guichard-Wienhard to that of Kapovich-Leeb-Porti.
%A $(P_+, P_-)$-Anosov property for $\rho$ is equivalent to 
%that of a $P_{\tau_{mod}}$-Anosov property (resp. $P_{\iota(\tau_{mod})}$-Anosov property) for 
%$P_+ = P_{\tau_{mod}}$ (resp. $P_- = P_{\iota(\tau_{mod})}$) for $\rho$ 
%and $\tau_{mod}$ a side in the spherical Weyl chamber $\sigma_{mod}$ determined by $P_+$
%and $\iota(\tau_{mod})$ is a side in the Weyl chamber $\sigma_{mod}$ determined by $P_-$. 
%according to Kapovich-Lee-Porti. 
%Since we assume $P^+$ is conjugate to $P^-$, 
%we have $\iota(\tau_{mod}) = \tau_{mod}$. 






%\begin{lemma}
%Let $g \in \Gamma$ and let $\mathcal{L}(\Gamma)$ be partially hyperbolic.
%Then $g$ has an eigenvalue equal to $1$ and 
%$\dim V_+ + \dim V_- < n$. 
%\end{lemma}
%\begin{proof} 
%Since $g$ acts without a fixed point in $\mathds{A}^n$, there is an eigenvalue $1$ of $g$. 
%%The Zariski closure $G$ of $\mathcal{L}(\Gamma)$ is an algebraic group in 
%%$\SL_{\pm}(n, \bR)$ with every element having eigenvalue equal to $1$. 
%%Every element of the maximal torus of $G$ has $1$ as the eigenvalue. 
%
%%Suppose that $n$ is even. 
%%It follows that there must be at least two $1$ since otherwise by taking inverses we may 
%%obtain elements without eigenvalue $1$. 
%
%Suppose that $m=1$. Then $g$ acts on the flag consists of $\{F_{+}\}$
%and a transversal flag $\{F_{-}\}$ and $\dim F_{+} = \dim F_{-}$. 
%Hence, $n$ is even. 
%Let $L_{g}$ be the eigensubspace of dimension $\geq 2$ corresponding to the eigenvalue $1$. 
%$\mathcal{L}(g)$ has a block forms of dimension $n, n$. 
%A diagonalizable maximal torus $M$ acts on $F_{+}$ and $F_{-}$. 
%%Then the dimension of $L_{k} \cap F_{+}$ and $L_{k}\cap F_{-}$ are equal
%%for each element of $k$ of $M$. 
%Hence, $L_{g} \cap F_{+}$ and $L_{g} \cap F_{-}$ are both $\geq 1$. 
%Let $V$ be a span of  $L_{g} \cap F_{+}$ and $L_{g} \cap F_{-}$. 
%In this case, we may obtain a transverse subspace  $F'$ of dimension $n$ 
%to $F_{+}$ and $F_{-}$ meeting $V - F_{+} - F_{-}$.
%Then $g^{n}(F')$ does not converge to $F_{+}$ as $n \ra \infty$. 
%This contradicts the definition of the Anosov representations. 
%
%%Suppose $m=1$. Then each nonidentity element $g$ fixes 
%%to subspaces $W$ and its complement $W'$ where $\dim W = \dim W'$, 
%%and they both have nonunit norm for eigenvalues all greater than $1$ or 
%%smaller than $1$. 
%%Hence $g$ acts on two affine spaces of complementary dimensions
%%meeting on a point. Hence, $g$ does not act freely. 
%%From now on assume $m \geq 3$. 
%
%
%\end{proof} 
%
%\marginpar{This is a bit strange place to be}
%\marginpar{The KLP definition is not really here... What should I say? }

%We will be interested in the subset $I$ in $\Delta$ containing 
%$\{\alpha_{k+1}, \dots, \alpha_{n-k-1}\}$ 
%not containing $\lambda_k$ and $\lambda_{n-k}$ for some $k$. 
%We will denote $P^I$ by $P^J$ also where $J$ is the subset of 
%$\{1, 2, \dots, n-1\}$ corresponding to $I \subset \Delta$
%by the correspondence $j \mapsto \alpha_j= \lambda_j - \lambda_{j+1}$. 
%A {\em sequential subset} of $\{1, 2, \dots, n-1\}$ is 
%a sequence increasing by $1$ at each step. 
%Let $J$ be a subset of $\{1, \dots, n-1\}$.
%Let $S(J)$ denote the set of all maximal sequential subset of $J$.
%
%
%For each sequential subset $K$ of $J$, 
%we denote by block $B^K$ in $K'\times K'$ where $K'$ is $K \cup \{k+1\}$ for 
%the last element $k$ of $K$. 
%Then we consider a matrix of block diagonal form
%$B^J:= \oplus_{K \in S(J)} B^K$ for each maximal sequential subset $K$ of $J$. 
%Note that $P^J$ can be understood as the matrix subgroup of $\SL_\pm(n, \bR)$ 
%where only nonzero entries are in $B^J$ or on the upper triangular matrix. 
%These follow easily from Theorem \ref{thm:GJT}. 
%
%Using the terminology of Kapovich-Lee-Porti \cite{KLP18}, 
%$P^J = P_{\tau_{mod}}$ where $\tau_{mod}$ is given 
%as  the spherical projection of 
%$\clo(\mathfrak{a}^+) \cap \mathfrak{a}_J$, i.e., 
%it is given by 
%$a_1 \geq \dots \geq a_n$ 
%where $a_i = a_{i+1}$ for $i\in J$. 
%We denote  this simplex by $\tau_{mod}^J$. 
%
%Since we require $P_+$ to be conjugate to $P_-$, 
%we have $\tau_{mod}^J$ is equivalent to $\iota(\tau_{mod}^J)$	
%under the longest length 
%product of Weyl reflections of the Weyl chamber. 
%(See Definition 5.62 in \cite{KLP17}.)
%We may restrict to this case only by Lemma 3.18 of \cite{GW12} as stated above.
%This means that the map $\{ 1, 2, \dots, n-1\}$ to itself 
%given by $i \ra n-i$ sends $J$ to itself. 

%Let $J =\{ 1, \dots n-1\} - \{k, n-k\}$. 
%Then we say that $P^J$-Anosov representation is 
%{\em $k$-Anosov}. 
%If $I$ does not contain $k$ and $n-k$, 
%$P^I$-Anosov representation is said to be $k$-Anosov
%(in the singular value sense).

%\marginpar{All these distinctions.. too many confusing also.}

\subsection{$P$-Anosov representations and dominations}

%The purpose of this subsection is to prove Corollary \ref{cor:main}.

%We will show that partially hyperbolic representations are $P$-Anosov for some parabolic group $P$.
%We also showed that a $P$-Anosov reprentations are partially hyperbolic. 

Let $\rho: \Gamma_M \ra \GL(n, \bR)$ denote a representation. 
%We consider Sambarino's criterion for $P$-anosov representations: 
%We order singular values 
%\[a_1(\rho(g)) \geq a_2(g) \geq \dots \geq a_n(g)
%\hbox{ of } \rho(g) \hbox{ for } g \in \rho(\Gamma_M).\] 
%
%
% 
%Bochi-Potrie-Sambarino \cite{BPS} says that for P-Anosov representation $\rho$, 
%there exists  a fixed $j$ such that 
%$a_j(g)/a_{j+1}(g) \geq C \exp(a |g|)$ for the word length $g$ in $\Gamma_M$. 
%This is equivalent to the P-Anosov condtion for a choice of a parabolic group $P$. 
%
%
%
%Notice $a_i(g^{-1}) = a_{n-i+1}(g)^{-1}$ for $i=1, \dots, n$. 
%Hence 
%\[ \frac{a_j(g^{-1})}{a_{j+1}(g^{-1})} = \frac{a_{n-j}(g)}{a_{n-j+1}(g)} \geq C \exp(a|g|). \]
%Since there is a map $g \mapsto g^{-1}$ 
%with $|g| = |g^{-1}|$, 
%if $j$ satisfies the above then $n-j$ also satisfies above. 
%
%	Suppose that $j$ is the largest minimal number $< n/2$ where 
%$a_j(g)/a_{j+1}(g)$ form an exponential function. 
%Then $j$ is called a {\em BPS-number}.
%
%We call the least number satisfying the 
%property below for $i_0$ the {\em neutral direction index} for $g\in \Gamma_M$
%or its holonomy element. 
%
%
%
%
%\marginpar{A problem,, there it is an eigenvalue...}
%A complete affine $n$-manifold is a $n$-manifold $M$ 
%with a diffeomorphism $M \ra \mathds{A}^n/\Gamma$. 
%
Let $w(g)$ denote the word length of $g$. 
Suppose that there exists an integer $k$,
$1 \leq k \leq n-k+1 \leq n$, 
so that the following hold for a constant $A, C > 0$: 
%\begin{multline} 
%\{a_1(\rho(g)), \cdots, a_k(\rho(g))\}, 
%\{a_{k+1}(\rho(g)), \cdots, a_l(\rho(g))\}, \\
%\{ a_{l+1}(\rho(g)), \dots, a_n(\rho(g)) \}
%\hbox{ with } k+1 \leq i_0(\rho(g)) \leq l 
%\end{multline} 
%\begin{itemize}  
	%%\item $k+1(\rho(g)) \leq i_0 \leq l$
	%% and $1/C< a_{k+1}(\rho(g)), \dots, a_{l}(\rho(g)) < C$
	%%for some uniform $C > 1$
%	\item 
\begin{equation*}
 \frac{a_k(\rho(g))}{a_{k+1}(g)} \geq C\exp(A w(g)) \hbox{ and }
%	\item 
  \frac{a_{n-k}(\rho(g))}{a_{n-k+1}(g)} \geq C\exp(A w(g)).
\end{equation*} 
	%$a_{k'}(\rho(g))$ for at least one index $k'(\rho(g))$, $k  < k'(\rho(g)) \leq l$, 
	%forms subexponential functions of $\Gamma_M$. 
%\end{itemize} 
%Since we can use $g^{-1}$ also, 
%We may restrict to the cases $l+1 = n-k+1$. 
%So, it must be that $l= n-k$. 
%Since there is $i_0$ such that $a_{i_0} = 1$, 
%we have $k < n/2$ 
In this case, we say that $\rho$ is {\em $k$-dominated}
%	in the singular value sense} 
for $k \leq n/2$ 
(see Bochi-Potrie-Sambarino \cite{BPS}).


%We will relate this with the definition in the bundle sense. 




%\marginpar{I need to correct this part... not right now..}

%In Section \ref{sub:notation}, we identify $G\hat M$ with 
%$\Uu \hat M$. 
%and Corollary 8.3H imply that 
%We can identify $\Uu \hat M$ with $\partial^{(2)}\hat M \times \bR$ 
%by Theorem 8.3C of Gromov \cite{Gromov87} (see Theorem 1 of Matheus \cite{Matheus90} or %Champetier \cite{Champetier94}).

%We can identify their space denoted $\widetilde{U\Gamma}$ with $\Uu \hat M$
%since the Gromov boundary $\partial \Gamma$ is identifiable with 
%the Gromov boundary of $\hat M$ with a hyperbolic metric 
%and $\hat M$ and the Cayley graph of $\Gamma$ are quasi-isometric
%by Theorem 11.108 and Corollary 11.29 of \cite{DK2018}.  
%\marginpar{Precise reference for it? The identification is obivous but precise? }

%The Gromov flow space $\widetile{U\Gamma_M}$ is defined 
%as $\partial^{(2)} \Gamma_M \times \bR$ 
%and $\phi_t((x, y, x)) = (x, y, s+t)$ for $x, y \in \partial \Gamma_M, x \ne y$, 
%and $s, t \in \bR$. 
%
%The Cayley graph $G$ of $\Gamma$ embed quasi-isometrically to $\hat M$. 
%By Theorem 11.108 of \cite{DK2018}, 
%
%We will identify $\Uu \hat M$ with the Gromov space 
%$\widetilde{U\Gamma}$. 

%By Proposition 4.9 of Bochi-Potrie-Sambarino \cite{BPS}, 
%if $\rho$ is partially hyperbolic with index $k$, then 
%$\rho$ is $k$-Anosov in their sense.

We will use $\llrrV{\cdot}$ to indicate the Euclidean norm in a maximal flat 
in a symmetric space $X$. (See Example 2.12 of \cite{KL18}.)
We will use the following notation: 
\begin{multline*} 
\begin{split} 
\vec{a}(g) & :=(a_1(g), \dots, a_n(g)), \\
\log \vec{a}(g) & := (\log a_1(g), \dots, \log a_n(g)) \in \clo(\mathfrak{a}^+_n)  \red{\hbox{(Cartan projection)}},\\ 
\log \vec{\lambda}(g)  &:= (\log \lambda_1(g), \dots, \log \lambda_n(g)) 
\in \clo(\mathfrak{a}^+_n) \red{\hbox{(Jordan projection)}}
\end{split} 
\end{multline*} 
for the singular values $a_i(g)$ and 
the modulus $\lambda_i(g)$ of the eigenvalues of $g$
for $g \in \GL(n, \bR)$.


%
%\begin{equation}\label{eqn:delta} 
% \log \vec{a}(\rho(g)), g \in G,
%\end{equation} 
%measures the Weyl-chamber-valued distance 
%$d_{\Delta}(x, \rho(g)(x))$ from 
%$x$ to $\rho(g)(x)$ in the maximal flat in $X$ where $g$ acts on 
%(see Section 2 of \cite{KL18}.)
% for the first statement. 

Given a set $A$, 
we say that two functions $f, g:A \ra \bR$ are {\em compatible} if 
there exists a uniform constant $C>1$ such that 
$C^{-1} (x) < f(x) < C g(x)$ for all $x\in A$. 

%We show a partial converse to 
\begin{lemma}\label{lem:Equiv} 
	Let a finitely generated \purple{word-hyperbolic} group $G$ have a representation $\rho:G \ra \GL(n, \bR)$.
	% whose linear part $\rho$ is 

	Then the following are equivalent\/{\em :} 
	\begin{itemize} 
	\item  $\rho$ is $P$-Anosov in the Gu\'eritaud-Guichard-Kassel-Wienhard sense \cite{GGKW17ii}
	for some parabolic group $P = P_\theta$ for an index set $\theta$ containing  $k$
	and $n-k$. 
	%	Suppose that the image of $\rho$ has the reducive Zariski closure. 
	%	and let $j$ be a BPS-number of the representation. 
	%	Let $\rho$ be a linear part of the holonomy representation of 
	%	a complete affine $n$-manifold $M$. 
	%	Suppose that for every $g$, $g \in \Gamma_M$,  
	%	we have $j < i_0(\rho(g)) < n-j+1$
	%	$i_0(\rho(g))$ such that $a_{i_0(\rho(g))}=1$
	%	and $a_q(\rho(g))$ for each $q$, $j < q < n-j +1$, form a subexponential
	%	function. 
%	\item $\rho$ is $Q$-Anosov in the sense of Kapovich-Lee-Porti
%	for some parabolic group $Q = P^I$ for an index set $I$ missing $k$
%	and $n-k$. 
	\item  $\rho$ is a $k$-dominated representation 
	%in the singular value sense 
	for $1 \leq k \leq n/2$. 
	\end{itemize} 
\end{lemma} 
\begin{proof} 
	This follows \purple{from} Theorem 1.3(3) of \cite{GGKW17ii}
	since the set $\theta$ contains $k$ and $n-k$.  
\purple{Or one can use Propositions 4.5 and 4.9 of \cite{BPS}.}
%	\qed
	%	Let $X$ denote the symmetric space associated with $\SL_\pm(n, \bR)$. 
%	We will use the notation of \cite{KLP17} for this proof. 
	%By Proposition 5.61 of \cite{KLP17}, 
%The first two items are equivalent by Proposition 5.61 of 
%	Appendix 5.11 of \cite{KLP18}. 
%	Let $\tau_{mod}$ be a side of the spherical Weyl chamber $\sigma_{mod}$ so that 
%	$P = P_{\tau_{mod}}$. Let $\Delta_{mod}$ denote the Euclidean Weyl chamber 
%	corresponding to $\sigma_{mod}$. 
%	Theorem 1.1 of \cite{KLP17} implies that 
%	the $P$-Anosov property implies the $\tau_{mod}$-URU property and 
%	hence the uniform $\tau_{mod}$-regularity of $\rho$.
%	Definition 4.27 of \cite{KLP17} says that 
%	every sequence $(g_n)$ of distinct elements of $\rho(G)$ satisfies 
%	\begin{equation} \label{eqn:dDelta} 
%		\liminf_{n\ra +\infty} \frac{d(\delta_n, \partial_{\tau_{mod}}\Delta_{mod})}{\llrrV{\delta_n}} > 0   
%		\hbox{ for } \delta_n =   d_{\Delta}(x, \rho(g_n)(x)), x\in X.
%	\end{equation}  
%using \eqref{eqn:delta}. 
%	while $\delta_n = \delta(g_n)$. 
	%	where $o$ is the origin of $X$.
%	Since $\rho$ is $\tau_{md}$-URU, $\rho$ is undistorted.  	
%	The map $g \mapsto w(g)$, $g\in \Gamma_M$, is compatible with 
%	the map $g \mapsto \delta(g)$, 
%	$g\in \Gamma_M$. 
%	Since $\tau_{mod}$ is given by setting $\log a_i = \log a_{i+1}$ for $i\in I$
%	on $\mathfrak{a}^+$, and $k$ is missing from $I$,
%	Since $P = P^I$ for some index set $I$, and choosing $P^I$ conjugate to 
%	its opposite parabolic group, 
%the $k$-domination of Bochi-Potrie-Sambarino \cite{BPS}
%	holds for some $k, 1 \leq k \leq  n/2$.
%	%
%	%
%	Finally, a $k$-dominated representation is $n-k$-dominated as we can 
%	obtain from taking an inverse map $g \mapsto g^{-1}$. 
%	Theorem 8.4 of \cite{BPS} implies that 
%	$\rho$ is $P^I$-Anosov in the Guichard-Wienhard sense with 
%	$I$ missing $k$ and $n-k$. 
%	%\qed 
\end{proof} 

%\marginpar{Dropped a URU lemma here}
%\begin{lemma} \label{lem:URU} 
%%	Let $X$ be the symmetric space 
%%	associated with $\SL_\pm(n, \bR)$. 
%	Suppose that	$\rho:G\ra \GL(n,\bR)$
%	is a representation of a finitely generated discrete group $G$. 
%	%\item $\delta(g) = \log \vec{a}(\rho(g))$, $g \in G$, 
%	%measures the Weyl-chamber-valued distance 
%	%$d_{\Delta_{mod}}(x, \rho(g_n)(x))$ from 
%	%$x$ to $\rho(g)(x)$ in $X$. 
%	Suppose that $\rho$ is $P$-Anosov and
%	the  image has the reducive Zariski closure. 
%	%I f $\rho$ is $P$-Anosov, then 
%	Then three functions $\Gamma \ra \bR$
%	given as follows  are compatible: 
%	\begin{itemize} 
%		\item $g \mapsto\llrrV{\log \vec{a}(g)}$, %$g \mapsto \llrrV{d_{\Delta_{mod}}(x, \rho(g)(x))}, g\in \Gamma$, 
%		\item $g\mapsto w(g)$, $g \in \Gamma$, and 
%		\item $g \mapsto \llrrV{\log\vec{\lambda}(g)}, g\in \Gamma$. 
%	\end{itemize} 
%\end{lemma} 
%\begin{proof} 
%	%	By Proposition 5.61 of \cite{KLP17}, the original Guichard-Wienhard definition is equivalent to the Kapovich-Lee-Porti definition.
%	The compatibility of the first two follows 	by %Remarks 1.5 (or Theorem 1.3(3)) 
%	(2.11) and 2.19 of \cite{GGKW17ii}.  %since it is easy to see that the first item is bounded above by a constant times the word length. 
%	%by Lemma \ref {lem:Equiv} 
%	%since the domination of $w(\cdot)$ to $\log \vec{a}(\cdot)$ is clear from matrix multiplications 
%	%and the reverse domination follows from the $k$-domination. 
%%	by Section 2.5.3 of \cite{GGKW17ii}. 
%	%since
%	%	Theorem 1.1 of \cite{KLP17} implies that 
%	%	the respresentation is quasi-isometric. 
%
%We have
%$\llrrV{\log\vec{\lambda}(g)} = \lim_{i\ra \infty}\frac{1}{i} \llrrV{\log\vec{a}(g^i)}$
%by Lemma 2.2 of \cite{BS2018p}. %\cite{Benoist98}, 
%	Also, the stable length $\left| g\right|_\infty := \lim_{i \ra \infty} \frac{1}{i} w(g^i)$ is compatible 
%	with $w(g)$ by Proposition 6.4 of Chapter 10 of \cite{CDP90}. 
%	Since the right sides are compatible by the first two items, 
%	the compatibility follows. 
%%	\qed
%\end{proof} 
%


\section{The proof of Theorem  \ref{thm:main2}} \label{sec:PAnosov} 



\subsection{$k$-convexity} \label{sub:kconvexity} 
%%Since $\hat M$ is negatively curved, $\Gamma_M$ is hyperbolic. 
\reD{We will relate the partial hyperbolic property in the singular-value sense to that in the bundle sense
	following \cite{BPS}.  We need this modification since we have to have actual expanding and shrinking under the flow.}

We continue to assume that $\Gamma_M$ is \purple{word-hyperbolic}. 
Let $\rho:\Gamma_M \ra \GL(n, \bR)$ be a representation. 
%Let $\mathcal{G}_k(\bR^n)$ denote the space of $k$-dimensional subspaces of 
%$\bR^n$. 
%\marginpar{What am I assuming here? }
%Recall that we can identify $\partial_\infty \hat \Gamma_M$ for the Cayley graph $\hat \Gamma_M$ of $\Gamma_M$ with 
%$\partial_\infty \hat M$ by Theorem 11.108 of \cite{DK2018}. 
%Hence, we can identify $\partial_\infty^{(2)} \hat \Gamma_M$ with $\partial^{(2)}_\infty \hat M$. 
%The Gromov flow space $\partial_\infty^{(2)} \hat \Gamma_M \times \bR$ is identified with 
%$\partial^{(2)}_\infty \hat M \times \bR$. 
%
%Let $\mathcal{G}_k(\bR^n)$ denote the space of $k$-dimensional subspaces of 
%$\bR^n$. 
%By Proposition 4.9 of Bochi-Potrie-Sambarino \cite{BPS}, $\rho$ is $k$-Anosov in 
%their sense. 
%We can identify $\partial G$ for the Cayley graph $G$ of $\pi_1(\hat M)$ with 
%$\partial_\infty \hat M$ by Theorem 11.108 of \cite{DK2018}. 
%Hence, we can identify $\partial^{2} G$ with $\partial^{2} \hat M$. 
%
%%We can identify their space denoted $\widetilde{U\Gamma}$ with $\Uu \hat M$
%%since the Gromov boundary $\partial \Gamma$ is identifiable with 
%%the Gromov boundary of $\hat M$ with a hyperbolic metric 
%%and $\hat M$ and the Cayley graph of $\Gamma$ are quasi-isometric
%%by Theorem 11.108 and Corollary 11.29 of \cite{DK2018}.  
%%\marginpar{Precise reference for it? The identification is obivous but precise? }
%
%%The Gromov flow space $\widetile{U\Gamma_M}$ is defined 
%%as $\partial^{(2)} \Gamma_M \times \bR$ 
%%and $\phi_t((x, y, x)) = (x, y, s+t)$ for $x, y \in \partial \Gamma_M, x \ne y$, 
%%and $s, t \in \bR$. 
%%
%%The Cayley graph $G$ of $\Gamma$ embed quasi-isometrically to $\hat M$. 
%%By Theorem 11.108 of \cite{DK2018}, 
%%
%%We will identify $\Uu \hat M$ with the Gromov space 
%%$\widetilde{U\Gamma}$. 
%
We say that the representation $\rho$ is {\em  $k$-convex} if there exist continuous maps 
$\zeta: \partial_\infty \Gamma_M \ra \mathcal{G}_k(\bR^n)$ and $\theta:\partial_\infty \Gamma_M\ra \mathcal{G}_{n-k}(\bR^n)$ such that the following hold:
\begin{description} 
	\item[(transversality)] for every $x,y\in\partial_\infty \Gamma_M,x\ne y$, we have $\zeta(x)\oplus\theta(y)=\bR^n$, and
	\item[(equivariance)] for every $\gamma \in \Gamma_M$, 
	we have \[\zeta(\gamma x) = \rho(\gamma)\zeta(x), \theta(\gamma x) = \rho(\gamma)\theta(x), x \in \partial_\infty \Gamma_M.\]
\end{description}

%By Theorem \ref{thm:Mineyev}, %and Theorem \ref{thm:identify}, 
%we may let $\hat \Gamma_M$ to be  % $\Uc \hat M$ coarsely equivalent to 
%$\partial_\infty^{(2)} \Gamma_M \times \bR$ with some metric. 

Using the representation $\rho$, it is possible to construct a linear 
flow $\psi_t$ over the geodesic flow $\phi_t$ of 
the Gromov flow space $\partial^{(2)}_\infty \hat M \times \bR$ as follows. 
Consider the lifted geodesic flow $\hat \phi_t$ on 
$\partial^{(2)}_\infty \hat M \times \bR$, and
define a linear flow on $\hat E := 
(\partial^{(2)}_\infty \hat M \times \bR)\times\bR^n$ by:
$\hat\psi_t(x,\vec{v})=(\phi_t(x), \vec{v})$ for $x\in\partial^{(2)}_\infty \hat M \times \bR$.
Now consider the action of $\Gamma_M$ on $\hat E$ given by:
\[ \gamma \cdot \left(x, \vec{v}\right):= 
\left(\gamma(x), \rho(\gamma)(\vec{v})\right). \] 
It follows that $\hat \psi_t$ induces a flow on $E_\rho:=\bR^n_{\rho}= \hat E/\Gamma_M$. 

%\marginpar{What is the geodesic here again? Check out}
When the representation $\rho$ is $k$-convex, by equivariance, there exists a $\phi_t$-invariant splitting of the form $E_\rho = Z \oplus \Theta$; 
it is obtained by taking the quotient of the bundles
\[\hat Z(x) = \zeta(x_+) \text{ and }
\hat \Theta(x) = \theta(x_-)
\text{ for } x \in \partial^{(2)}_\infty \hat M \times \bR\]
where $x_+$ is the forward endpoint of 
the complete isometric geodesic through $x$ and 
$x_-$ is the backward endpoint of the complete isometric geodesic through $x$. 

We say that a $k$-convex representation is {\em $k$-Anosov
in the bundle sense} if the splitting 
$E_\rho=Z\oplus\Theta$ is a dominated splitting for the linear bundle automorphism $\psi_t$, with $Z$ dominating $\Theta$ in the terminology of \cite{BPS}. 
This is equivalent to the fact that the bundle $\Hom(Z,\Theta)$ is uniformly contracted
by the flow induced by $\psi_t$. 


%
%Directly following Bochi-Potrie-Sambarino \cite{BPS}, we obtain: 
%\begin{theorem} \label{thm:bundle} 
%	Let $\rho$ be a partially hyperbolic representation with an index $k$. 
%	%	Suppose that there is a neutrality function 
%	%	$\Gamma_M \ra [k+1, n-k]$. 
%	% neutrality index for each holonomy element is in $[k+1, n-k]$. 
%	Then 
%	$\rho$ is partially hyperbolic in the bundle sense with index $k$.
%\end{theorem} 
%\begin{proof} 
%	The proof is same as that of Theorem 2.4 in \cite{AF1} 
%	where we work over $\Uc \hat M$ and $\Uc M$. 
%\end{proof} 

%\begin{proof}[Proof of Corollary \ref{cor:main}] 
%	The proof is the same as in Corollary 1.3 of \cite{AF1} by Lemma 2.5 of \cite{AF1} 
%	and Theorem \ref{thm:bundle}.  
%	
%	\end{proof} 
%
%\begin{proof}[Proof of Corollary \ref{cor:main2}]
%	 The proof is the same as that of Corollary 1.4 of \cite{AF1} by Lemma 2.5 of \cite{AF1} 
%	 and Corollary \ref{cor:main}. 
%	\end{proof} 


%We say that the representation $\rho$ is {\em  $k$-convex} if there exist continuous maps $\zeta: \partial \Gamma_M \ra \mathcal{G}_k(\bR^n)$ and $\theta:\partial\Gamma_M\ra \mathcal{G}_{n-k}(\bR^n)$ such that:
%\begin{description} 
%	\item[(transversality)] for every $x,y\in\partial\Gamma_M,x\ne y$, we have $\zeta(x)\oplus\theta(y)=\bR^n$, and
%	\item[(equivariance)] for every $\gamma \in \Gamma_M$, 
%	we have \[\zeta(\gamma x) = \rho(\gamma)\zeta(x), \theta(\gamma x) = %\rho(\gamma)\theta(x), x \in \partial \Gamma_M.\]
%\end{description}
%Using the representation $\rho$, it is possible to construct a linear 
%flow $\psi_t$ over the geodesic fow $\phi_t$ of $\Uu M$ as follows.
%By the identification $\Uc \hat M$ with $\partial^{(2)}\hat M \times \bR$, we
%consider the lifted geodesic fow $\hat \phi_t$ on $\Uc \hat M$, and
%define a linear flow on $\hat E := \Uc \hat M \times\bR^n$ by:
%$\hat\psi_t(x),v)=(\phi_t(x),v)$ for $x \in \Uc \hat M$.
%Now consider the action of $\Gamma_M$ on $\hat E$ given by:
%\[ \gamma \cdot (x, v):= 
%(\gamma(x), v). \] 
%It follows that $\phi_t$ induces in a flow on $E:=\bR^n_{\rho}= \hat E/\Gamma_M$. 
%
%When the representation $\rho$ is $k$-convex, by equivariance there exists a $\phi_t$-invariant splitting of the form $E_\rho = Z \oplus \Theta$; 
%it is obtained taking the quotient of the bundles
%$\hat Z(x) = \zeta(x_+)$ and 
%$\hat \Theta(x) = \theta(x_-)$
%for $x \in \Uc \hat M$ and $x_+$ is the forward endpoint of 
%the complete geodesic through $x$ and 
%$x_-$ is the backward endpoint of the complete isometric geodesic through $x$. 
%\marginpar{geodesic here? Check the meaning...}
%
%We say that a $k$-convex representation is {\em $k$-Anosov} if the splitting 
%$E_\rho=Z\oplus\Theta$ is a dominated splitting for the linear bundle automorphism $\psi_t$, with $Z$ dominating $\Theta$ in the terminology of \cite{BPS}. 
%This is equivalent to the fact that the bundle $Hom(Z,\Theta)$ is uniformly contracted
%by the flow induced by $\psi_t$. 


%Let $P_{k, n-k}$ be the subgroup of  $GL(n, \bR)$ where 
%only nonzero entries $a_{ij}$ satisfy one of the following holds.
%\begin{itemize} 
%	\item $i \leq k$, 
%	\item $k < i \leq n-k$ and $k < j$, or
%	\item $n-k < i \leq n$ and $n-k < j$. 
%\end{itemize} 

%Hence, there exists some $m$ so that 
%$a_{m+1}(\rho(g))/a_{m}(\rho(g)) \leq C \exp(a''w(g))$ for some $a''> 0$. 
%But $m$ may depend on $g$. 



Suppose that we further require for a $k$-dominated representation $\rho$
 \purple{where} $k \leq n/2$\/: 
\begin{itemize}  
	%%\item $k+1(\rho(g)) \leq i_0 \leq l$
	%% and $1/C< a_{k+1}(\rho(g)), \dots, a_{l}(\rho(g)) < C$
	%%for some uniform $C > 1$
	\item $a_p(\rho(g)) \geq C^{-1} \exp(A w(g))$ for $p \leq k$, and 
	\item $a_r(\rho(g)) \leq C \exp(-A w(g))$ for $r \geq n-k+1$
	\purple{and} some constants $C> 1, A> 0$. 
	%	\item $a_k(\rho(g))/a_{k+1}(g) \geq C\exp(A w(g))$, and
	%	\item $a_{n-k}(\rho(g))/a_{n-k+1}(g) \geq C\exp(A w(g))$.
	%$a_{k'}(\rho(g))$ for at least one index $k'(\rho(g))$, $k  < k'(\rho(g)) \leq l$, 
	%forms subexponential functions of $\Gamma_M$. 
\end{itemize} 
Then we say that $\rho$ is {\em partially hyperbolic
	in the singular-value sense with the index $k$}. 
 
We obtain by Lemma \ref{lem:Equiv}:
\begin{lemma}\label{lem:phypAnos} 
	If $\rho$ is partially hyperbolic for an index $k$, $k \leq n/2$, in the singular-value sense, 
	then $\rho$ is $P_\theta$-Anosov for the index set $\theta$ containing $k$ and $n-k$. 
	\qed
\end{lemma} 
%\begin{proof} 
%	Our condition implies that $\rho$ is $k$-dominated in the sense of 
%	\cite{BPS},
%	implying that $\rho$ is $k$-Anosov. 
%	By Proposition 4.9 of \cite{BPS}, 
%	$\rho$ is $k$-Anosov. 
%Here, $P^I = P_{\tau_{mod}^I}$ by the terminology of \cite{KLP18}.  
%Since $\rho$ is $P^I$-Anosov, 
%$\rho$ is $\tau_{md}^I$-regular by Theorem 5.47 (iii) and (v). 
%$a_k(\rho(g))/a_{k-1}(\rho(g))$ form an exponential function $\Gamma_M \ra \bR^+$.
%Basically use \cite{KL18} 
%where the vector 
%$(a_1(\rho(g)), \dots, a_n(\rho(g)))$ converges. Details... to do...
%The flag variety for subspaces of dimensions $k, n-k, n$ is needed. 
%\end{proof} 

%We define 
Recall $\hat M$ the cover of $M$ with the deck transformation group $\Gamma_M$. 
Recall from Section \ref{sub:Gromov}  that we can identify $\partial_\infty \Gamma_M$ for the Cayley graph $\hat \Gamma_M$ of $\Gamma_M$ with 
$\partial_\infty \hat M$ by Theorem 11.108 of \cite{DK2018}. 
Hence, we can identify $\partial_\infty^{(2)} \Gamma_M$ with $\partial^{(2)}_\infty \hat M$. 
The Gromov flow space $\partial_\infty^{(2)} \Gamma_M \times \bR$ is identified with 
$\partial^{(2)}_\infty \hat M \times \bR$. 
By directly following Proposition 4.9 of 
Bochi-Potrie-Sambarino \cite{BPS}, we obtain: 
\begin{theorem} \label{thm:bundle} 
Let $\rho: \Gamma_M \ra \GL(n, \bR)$ be a representation.
	Let $\rho$ be partially hyperbolic with an index $k$, $k < n/2$, 
	in the singular-value sense. 
	%	Suppose that there is a neutrality function 
	%	$\Gamma_M \ra [k+1, n-k]$. 
	% neutrality index for each holonomy element is in $[k+1, n-k]$. 
	Then 
	$\rho$ is partially hyperbolic in the bundle sense with the index $k$.
\end{theorem} 
\begin{proof} 
By Lemmas \ref{lem:phypAnos}  and \ref{lem:Equiv}, 
	$\rho$ is $k$-dominated.

%	they \cite{BPS} prove this on the Gromov flow space 
%	$\partial^{(2)}\hat M \times \bR$. 
	% as in
	%Lemma \ref{lem:phypAnos}. 
	Bochi, Potrie, and Sambarino \cite{BPS} 
	construct subbundles $Z$ and $\Theta$ of $E_\rho$ 
	where $Z$ dominates $\Theta$
	and $\dim Z = k, \dim \Theta = n-k$. 

	Now, $\rho$ is also \purple{$(n-k)$-dominated}. 
	By Lemma \ref{lem:phypAnos}, 
	we obtain a new splitting of the bundle $E_\rho = Z'\oplus \Theta'$ 
	over $\partial_\infty^{(2)}\hat M \times \bR$
	where $\dim Z' = n-k$ and $\dim \Theta'$ is $k$. 
	Furthermore, $Z'$ dominates $\Theta'$. 
	Then $Z$ is a subbundle of $Z'$, and 
	 $\Theta'$ is a subbundle of $\Theta$ by 
	Proposition 2.1 of \cite{BPS}. 
	%
	%This is true over each closed orbit of $\phi_t$
	%since we can consider a single element corresponding to the orbit
	%and the fact that $\Theta'$ is also parallel under the connection. 
	%Since the union of closed orbits are dense in $\Uu M$
	%we have the result. 
	We now form bundles $Z, Z'\cap \Theta, \Theta'$ over $\partial_\infty^{(2)}\hat M \times \bR$. 
	%The exponential properties are obvious now
	%since $Z' \cap \Theta$ has a subbundle over each closed 
	%geodesic that has a $\phi_t$-invariant fiber-wise Euclidean metric under the flow. 
	%Over the closed orbits of $\Uu M$, $Z'\cap \Theta$ expands 
	%or contracts subexponentially. 
	The dominance property \eqref{eqn:dominance} of Definition \ref{defn:phyp}  follows from those of $Z, \Theta$ and $Z', \Theta'$. 
		
	There are exponential expansions 
	of $a_i(g)$ for $1\leq i \leq k$ by the partially 
	hyperbolic condition of the premise. 
	%since $a_k(g)$ is bounded below by downward 
	%subexponential function. 
	Now (iii)(a) of Definition \ref{defn:phyp} follows: 
	%since we can 
	%prove the exponential expansion in $Z$ under the flow 
	%by the characterizing equations (2.3) and (2, 4) 
	%of $Z$ and $\Theta$ in Theorem 2.2 of \cite{BPS},
	%and the expanding subspace has at least dimension $k$: 
	Define $U_p(A)$ as the sum of eigenspaces of $\sqrt{AA^*}$ corresponding to 
    \purple{the} $p$ largest eigenvalues. 
	Define $S_{n-p}(A):= U_{n-p}(A^{-1})$. 
	Since \reD{$\hat M/\Gamma_M$} is compact, Theorem 2.2 of \cite{BPS} shows 
	\[U_k(\psi^t_{\phi^{-1}_t(x)})\ra \hat Z(x) \hbox{ and } 
	S_{n-k}(\psi^t_x) \ra \hat \Phi(x) \hbox{ uniformly for $x \in \hat \Gamma_M$  as } t \ra \infty.\] 

	Since by the paragraph above Theorem 2.2 of \cite{BPS}, we obtain
	\[\psi^t(S_{n-k}(\psi^t_{\phi^{-1}_t(x)})^{\perp}) = U_k(\psi^t_{\phi^{-1}_t(x)}),\]
	where the unit vectors of $S_{n-k}(\psi^t_{\phi^{-1}_t(x)})^{\perp}$ goes to 
	vectors of norm at least the $k$-th singular value for $\psi_t$. 
		The angle between $\hat Z(x)$ and $\hat \Phi(x)$ is uniformly bounded below by 
a positive constant 
		by the compactness of \reD{$\hat M/\Gamma_M$}. 
		Hence, unit vectors in $\hat Z({\phi^{-1}_t(x)})$ have 
		components in $S_{n-k}(\psi^t_{\phi^{-1}_t(x)})^\perp$ whose lengths are uniformly bounded below by a positive constant 
		for sufficiently large $t$. 

	As $t \ra \infty$, $\phi_t$ passes the images of the fundamental domain under $\Gamma_M$. 
	Let $g_i$ denote the sequence of elements of $\Gamma_M$ arising in this way. 
		Since $a_k(g_i)$ grows exponentially, the expansion properties of $Z$ follow
\red{by the above paragraph}. 
	(iii) of  Definition \ref{defn:phyp} follows up to reversing the flow. 
	
			By the quasi-isometry of $\Uc \hat M$ with $\partial^{(2)}\hat M \times \bR$ in Theorem \ref{thm:identify},
	%	as we did in the beginning of this subsection,  
	we construct our pulled-back partially hyperbolic bundles over $\Uc \hat M$.
%	\qed 
	%	Then we apply our construction above 
	%	to construct our dominated bundles using Theorem 8.3C of \cite{Gromov87}. 
%	\cite{BPS}. 
%	
	%\marginpar{Need to be more clear add stuffs....exponential stuff not done completely}
%	
%	
	%	We show that $\rho$ is $P$-Anosov with $P$ 
	%	the parabolic group of form $P_{k, n-k+1}$. 
	%	Now, we use the flag space $\mathcal{F}$ of 
	%	$ V_1 \subset V_2 \subset V_3$ where $\dim V_1 = k$, $\dim V_2 = n-k$
	%	and $\dim V_3 = n$. 
	%	We construct for each $(t^+, t^-)\in \partial^{2}_\infty \hat M$, 
	%	a pair $\zeta^+(t^+)$ and $\zeta^-(t^-)$ that are transverse 
	%	in $\mathcal{F}^+ \times \mathcal{F}^-$.
	%	So for each $(t^+, t^-)$ we obtain two subspaces 
	%	of dimension $k$ and $n-k$ which are transversal. \marginpar{I don't know transversality}
	%	This pair gives us a decomposition of $\bR^n$ for 
	%	each pair $t^+, t^-$. 
	%	The decomposition is obviously $C^0$ since $\zeta^+$ and $\zeta^-$ is 
	%	a continuous function. 
	%	
\end{proof} 



\subsection{Promoting $P$-Anosov representations to  
partially hyperbolic ones}
\label{sub:converse} 

\reD{This subsection contains the main results of this paper. The complication is that the maximal abelian subgroup of the Zariski closure of the linear part of the holonomy group sometimes embeds not linearly but piecewise linearly. In addition, we have to understand the case where the Zariski closure is not reductive.}

%%Proposition 4.6 of \cite{BPS} gives the converse to Lemma \ref{lem:phypAnos}. 
%%Assume the following for a representation $\rho:\Gamma_M \ra \SL_\pm(n, \bR)$:
%Suppose that	$\rho:G\ra \SL_\pm(\bR)$
%is a representation of a finitely generated discrete group $G$
%with following properties. 
%\begin{itemize} 
%	%\item  $\rho$ satisfies the neutrality condition for $[k+1, n-k]$. 
%	%There is a number $k$, $1\leq k \leq n/2$ so 
%	%that there is a neutrality function $\Gamma_M \ra [k+1, n-k]$. 
%	%that a neutral-index of each holonomy element is in $[k+1, n-k]$. 
%	\item $\rho$ is $P^I$-Anosov in the Guichard-Wienhard sense for 
%	a parabolic group $P^I$ where $I$ is a subset of 
%	$\{1, \dots, n\}$ not containing $k$ and $n-k$
%	for $1 \leq k \leq n/2$. 
%	%\item Each holonomy $g$ has a unit singular-value index $i_0(\rho(g))$ 
%	%in $[k+1, n-k]$. 
%	\item $a_{i}(\rho(g)) \geq C^{-1} e^{Aw(g)}$ for $i \leq k$
%	and $a_{j}(\rho(g)) \leq C e^{-Aw(g)}$ for $j \geq n-k+1$
%	for some constants $C> 1, A> 0$. 
%\end{itemize} 
%Then we say that $\rho$ is a {\em  strictly $P^I$-Anosov} representation with index $k$
%in the Guichard-Wienhard sense. 
%Here, $P^I = P_{\tau_{mod}^I}$. 
%Since $\rho$ is $P^I$-Anosov, 
%$\rho$ is $\tau_{mod}^I$-regular by Theorem 5.47 (iii) and (v). 
%By Theorem 5.22 of \cite{KLP18}, $\rho$ is $k$-dominated
%since $\Gamma_M$ is word-hyperbolic. 


%We will use $\llrrV{\cdot}$ to indicate the Euclidean norm in the maximal flat 
%in a symmetric space $X$. (See Example 2.12 of \cite{KL18}.)
%\[\vec{a}(g) =(a_1(g), \dots, a_n(g)) \in A^n
%\hbox{ and } \vec{\lambda}(g) = (\lambda_1(g), \dots, \lambda_n(g)) \in A^n\] 
%for the singular values $a_i(g)$ and 
%the modulus $\lambda_i(g)$ of the eigenvalues of $g$
%for $g \in \SL_\pm(n, \bR)$.

%We say that two functions $f, g:\Gamma \ra \bR$ are {\em compatible} if 
%there exists a uniform constant $C>1$ so that 
%$C^{-1} g(\gamma) < f(\gamma) < C g(\gamma)$ for 
%$\gamma \in \Gamma$. 



%\begin{corollary}\label{cor:Anosphyp} 
%	Let a finitely generated discrete group $G$ have  linear holonomy representation $\rho$
%	% whose linear part $\rho$ is 
%	that is strictly $P^I$-Anosov with index $k$. 
%	Suppose that the image of $\rho$ has the reducive Zariski closure. 
%	Then $\rho$ is a strictly $k$-Anosov representation 
%		in the singular value sense.  	
%	\end{corollary} 
%\begin{proof}
%	Consider Lemma \ref{lem:Equiv} and 
%	the collections: 
%	\begin{multline} 
%		\{a_1(\rho(g)), \cdots, a_k(\rho(g))\}, 
%		\{a_{k+1}(\rho(g)), \cdots, a_{n-k}(\rho(g))\}, \\
%		\{ a_{n-k+1}(\rho(g)), \dots, a_n(\rho(g)) \}, g \in \rho(G). 
%	\end{multline} 
%	Then the $k$-domination and the second condition of the
%	strict $P^I$-Anosov property imply that 
%	$\rho$ satisfies the properties for the partial hyperbolicity of index $k$. %\qed
%%	Theorem \ref{thm:bundle} implies the result. 
%%	(See Example 2.12 of \cite{KLP18}.) %Theorem \ref{thm:bundle} implies the result. 
%\end{proof} 
%
%Outline of the proof 

%\marginpar{This was your part which I fill in. Here I still cannot overcome 
%the semisimplicity. The semisimplization is not enough since 
%I need $a_i$ for all $g$ not just asymptotic part of $g^n$.}

\purple{By a {\em cone} of a subset $J$ of a vector space $V$, we mean 
the subset of nonnegative multiples of points of $J$. We denote 
it by $C(J)$. This is not necessarily convex.}. 

For a connected reductive Lie group $Z$,  
let $A_Z$ denote the maximal $\bR$-split torus of $Z$  with \purple{the} corresponding 
Lie algebra $\mathfrak{a}_Z$, and 
 let $\mathfrak{a}_Z^+$ denote the Weyl chamber of $\mathfrak{a}_Z$.
%with a closed Weyl chamber $A^+_Z$, 
Also, we denote 
the Cartan projection by $\log \vec{a}_Z: Z \ra \clo(\mathfrak{a}^+_Z)$.
%obtained from the Cartan decomposition $Z = KT^+K$ where $K$ is a compact subgroup
%and $T^+$ is the diagonal 
Let $G$ be a finitely generated Zariski dense group in $Z$. (See Section 1 of \cite{BS2021} 
and Section 2.4 of \cite{GGKW17ii}.)
We denote by $\log \lambda_{Z}: Z \ra  \clo(\mathfrak{a}^+_Z)$ the associated 
Jordan projection.
The {\em Benoist cone}  $\mathcal{BC}_Z(G)$ of 
$G \subset Z$ is the closure of the set of 
all positive linear combinations of 
$\log \vec{\lambda}_Z(g), g\in \Gamma$ in $\clo(\mathfrak{a}^+_Z)$ and $O$.
\purple{In particular,  $\mathcal{BC}_Z(G) \setminus \{O\}$ is path-connected and convex
in $\mathfrak{a}_Z$.
(See the paragraph before Theorem 1.2 of \cite{BS2021}.)}

We remark that we put $\mathfrak{a}_n$ as the diagonal matrices,   
and the Jordan projection here is same as the Cartan projection. 
We will use $\vec{\lambda}$ for $\mathfrak{a}_n$ following the convention.


We \green{worked out} the following proposition with the help of Danciger and Stecker: 
\begin{proposition}\label{prop:PAnosovStrict}  
Assume the following\/{\em :} 
\begin {itemize}
	\item  $G$ is a finitely generated group. % that is not infinite cyclic. 
	\item $\rho:G \ra \GL(n, \bR)$ is a
	$P_\theta$-Anosov representation 
	for some index set $\theta =\theta^*$ 
	containing $k$ and $n-k$ for 
	$1\leq k \leq n/2$.
	\item 
\begin{itemize} 
\item[(i)] $\rho(g)$ for each $g\in G$ has $1$ as the norm of an eigenvalue, 
or, alternatively, 
\item [(ii)]  \purple{For any sequence $g_i \in \pi_1(N)$, if $w(g_i) \ra \infty$, then $\frac{|\log a_{j_{g_i}}(\rho(g_i))|}{w(g_i)} \ra 0$ for some choices of 
the indices $j_{g_i}\in \{1, \dots, n\}$ depending on $g_i$.} 
\end{itemize} 
	\item The Zariski closure $Z$ of the image of $\rho$ is 
	a connected reductive Lie group. 
\end{itemize} 
	Then $\rho$ is partially hyperbolic with the index $k$ in the singular-value sense, 
	and $k < n/2$.   
	\end{proposition}
\begin{proof} 
	(I) We begin with some preliminary results.  % \marginpar{Show $k < n/2$... }
%	Suppose that the linear part of $\rho$ is $P^I$-Anosov.
%	Let $k$ be the index $< n/2$ missing in $I$. 
%Let $A_Z$ denote the maximal $\bR$-split maximal torus of $Z$. 
%	 Let $A_Z^+$ denote the Weyl chamber of $A_Z$.
%	 Every element of $Z$ has an eigenvalue $1$ 
%since	 $\det(g - \Idd)=0$ for $g \in \rho(G)$. 
%	 
	 	The embedding $Z \hookrightarrow \GL(n, \bR)$
	 induces a map $\iota_Z:\clo(\mathfrak{a}_Z^+) \ra \clo(\mathfrak{a}^+_n)$ for the Weyl chamber
	 $\mathfrak{a}^+_n$ of $\GL(n, \bR)$. 
	 Let $\log\lambda_1, \dots, \log\lambda_n$ denote the 
	 coordinates of the Weyl chamber $\mathfrak{a}^+_n$, 
which are actually the standard coordinates.
%	 They correspond to components of 
%	 the Jordan projection $\log\vec{\lambda}: \SL_\pm(n, \bR) \ra A_n^+$ given 
%	 by the logarithms of norms of the eigenvalues ordered by size. 
	 Of course, $\iota_Z \circ \log \vec{a}_Z| 
	 \rho(G) = \log\vec{\lambda}|\rho(G)$. 

%By Theorem 1.3 of Breuillard and Sert \cite{BS2021}, the Benoist cone can be defined 
%by singular values or eigenvalues and they coincide. 

	 	Let $B:=\mathcal{BC}_Z(\rho(G))$ denote the 
	 Benoist cone of $\rho(G)$ in $Z$. 
%\purple{As stated above, $B -\{O\}$ is path-connected.}
Let $S$ be a generating subset of $\rho(G)$. Then 
\purple{$B$ is a cone of  the joint spectrum $J(S)$}
as stated immediately after Theorem 1.3 of \cite{BS2021}. 
\purple{That is, $B = C(J(S))$. }
 The image $\iota_Z(B)$ is not convex in general. 
	 %which is the smallest closed cone containing 
	 %$\log \vec{\lambda}_Z(\rho(\Gamma_M))$
	 %for the Jordan projection $\log \vec{\lambda}_Z: Z \ra A_Z^+$. 
	 %	By Benoist \cite{Benoist97}, 
	 %and the semisimplication of Section 2.9 of \cite{KP2002p}, 
%	 Now, $B$ is a convex cone with a nonempty interior in $\clo(\mathfrak{a}_Z^+)$. 	 

%Let $z_i$ denote the standard coordinates of
%the Weyl chamber $\clo(\mathfrak{a}^+_n)$ for $i=1, \dots, n$. 
%	 By Theorem \ref{thm:HKS}, 
%The condition of a matrix $A$ having an eigenvalue $1$ is an algebraic condtion 
%given by $\det(A - \Idd)=0$ and that  of $A$ having a singular value $1$ is 
%also given by $\det(AA^T - \Idd) = 0$. 

%Since $Z$ is the Zariski closure of $\rho(G)$, every element of $B$ has an eigenvalue $1$
%or alternatively each element of $B$ has a singular value $1$. 

%The condition of a matrix $A$ having an eigenvalue $1$ is an algebraic condtion 
%given by $\det(A - \Idd)=0$. The Zariski closure of $\rho(G)$ satisfies the same condition. 
%	 For each $b \in B$, there is an index $i$ where 
%	 the log of the $i$-th coordinate of $\iota_Z(b)$ equals $0$.

Let $S \subset \rho(G)$ be a generating set. 
\purple{The condition (ii) implies (i) by Lemma 2.11 of \cite{BS2021}. 
Under the condition (i), each vector $\log \vec{\lambda}(\rho(g)), g\in G$ 
has a coordinate with the value $0$.
By Theorem 1.2 of \cite{BS2021}, we have 
\[ \clo\left(\bigcup_{m=1}^\infty \frac{1}{m}\log \vec{\lambda}_Z(S^m)\right) = J(S)\]}
%By the definition of the Benoist cone
%being the limit of $1/n \vec{\lambda}(\rho(S^n)$ where $S$ is a generating 
%subset of $\pi_1(N)$, 
\purple{Thus 
for each $x \in J(S)$, 
there is a sequence of elements of $x_j \in \frac{1}{m_j}\log\lambda_Z(S^{m_j}) $ converging to $x$. 
Since $B$ is the \purple{cone of $J(S)$} in $a^+_Z$, 
we find that there must be an index $i$ where 
$\iota_Z(x)_i = 0$ for every $x\in B$. }


From now on, the arguments are the same for both conditions.
%Theorem \ref{thm:main2}. 
We denote by 
\[\mathcal{I}(b)=\{ i \mid \log \lambda_i(\iota_Z(b)) = 0\} \text{ for } b \in B.\] 
which is a consecutive set under the first premise.
%We denote by 
%\[\mathcal{I}(b)=\{ i | \log \lambda_i(\iota_Z(b)) = 0\},\] 
%which is a consecutive set under the second premise. 
%norm of the eigenvalue equals $0$.
	 Let $S([1, ..., n])$ denote the collection of 
	 subsets of consecutive elements.
  
	 We define an integral-interval function,
	 \[\mathcal{I}: B \setminus \{O\} \ra S([1, \dots, n])\]
	 given by sending $b \in B$ to $\mathcal{I}(b)$. 
 Since the sequences of
ordered norms of eigenvalues of a convergent sequence of linear maps converge,
we have the following: 
	 \begin{equation} \label{eqn:limI} 
	 \lim_{i\ra \infty} \mathcal{I}(x_i) \subset \mathcal{I}(x) \text{ whenever } 
	 x_i \ra x, x_i, x\in B.
	 \end{equation}
 
 	
 
 %	Since $B$ is convex,  
 %	it follows that $B$ is inside the hyperplane determined by 
 %	$x_i \equiv 0$ for a fixed $i$. 
 %	Since $\rho(\Gamma_M)$ is a group, 
 %	$B$ is invariant under the opposition involution $\iota$, and  
 %	there exists an index $j$ where $1 \leq j \leq n/2$ 
 %	where $x_i \equiv 0$ on $B$ for $i=j, \dots, n-j+1$. 
 %Let $w(g)$ denote the word length of $g$. 
 %Since $\rho$ is $P^I$-Anosov, 
 %$\rho$ has a uniform $k$-gap in eigenvalues for $k < n/2$
 %by Proposition 1.5 of Kassel-Potrie \cite{KP2002p}: 
 %	and taking inverses. 
 %\begin{equation}\label{eqn:lklkp} 
 %	\log \lambda_k(\rho(g)) - \log \lambda_{k+1}(\rho(g)) > C w(g) - C' \hbox{ for } 
 %	C, C'> 0 \hbox{ and } g \in \Gamma_M.
 %\end{equation} 
 %	Since we can take $g^n$, this equation becomes 
 %	\begin{equation}\label{eqn:lklkp} 
 %	\log \lambda_k(\rho(g)) - \log \lambda_{k+1}(\rho(g)) \geq C \left|g\right|_\infty 
 %	\end{equation}
 %	for each element $g$ in $\Gamma_M$, 
 %	where $ \left|g\right|_\infty := \lim_{i \ra \infty} w(g^n)/n$. 
 %	
 %	where the last follows by 
 %Proposition 2.1 of Chapter 2 of \cite{CDP90}.	
 %
 %	By the main density theorem of Benoist \cite{Benoist98}, 
 %	each element $b \in B \cap S_n$ has a sequence 
 %	$\{g_n\}, g_n \in \rho(\Gamma_M)$ where 
 %	$\frac{\log\vec\lambda(g_n)}{\llrrV{\log\vec{\lambda}(\rho(g_n))}}  \ra b$. 
 %	We define the {\em normalized word length} $w(b)$ of elements of $b \in B\cap S_n$  
 %	to be the infimum of the set of 
 %	the infimum limits of form $\left\{\frac{w(g_n)}{\llrrV{\log\vec{\lambda}(\rho(g_n)}}\right\}$ 
 %for any sequence 
 %	$\frac{\log\vec\lambda(g_n)}{\llrrV{\log\vec{\lambda}(\rho(g_n))}}  \ra b$.
 %%	We allow $0$ and $\infty$ as values here. 
 %%	Since $\log \lambda_1(\rho(g))$ is always bounded above by a constant 
 %%	times $w(g)$ for any linear representation, 
 %%	so is $\llrrV{\log\vec{\lambda}(\rho(g)}$.
 %	By Lemma \ref{lem:URU}, 
 %$\llrrV{\log \vec{\lambda}(\rho(g)}$ is compatible 
 %with word length of $g$. 
 %	Thus, it follows that $w(b)$ is always bounded below 
 %	by a positive constant for $b\in B\cap S_n$. 
 %
 %
 
Let $S$ denote a finite set of generators of $\rho(G)$ and their inverses. 
 By \purple{Theorem 1.3} of \cite{BS2021}, the sequence
 $\frac{1}{m} \log\vec{\lambda}_Z(S^m)\subset J(S)$ geometrically converges to 
 the compact convex set $J(S) \subset\clo(\mathfrak{a}_Z^+)$  as $m \ra \infty$ 
 %and is in $J(S)$ 
 (see Section 3.1 of \cite{BS2021}).
Moreover, the \purple{cone of $J(S)$} is the Benoist cone $B$. 
This implies that the set of directions of 
 $\log \vec{\lambda}_Z(\rho(G))$
 is dense in \purple{that of} $B$. (See also \cite{Benoist98}). 
 Let $S_n$ denote the unit sphere in the Lie algebra $\mathfrak{a}_n$ of $\GL(n, \bR)$
with respect to the standard coordinates.
 We obtain
 \[\log \lambda_k(g) -\log \lambda_{k+1}(g) \geq C \llrrV{\log \vec{\lambda}(g)}, g \in \rho(G),
 \hbox{ for a constant } C> 0 \]
 by Theorem 4.2(3) of \cite{GGKW17ii}. 
Hence, we obtain
 %with $w(g)$ by Lemma \ref{lem:URU},  \eqref{eqn:lklkp} implies
 \begin{equation} \label{eqn:lklkpC} 
 	\log \lambda_k(b) -\log \lambda_{k+1}(b) \geq C 
 \end{equation} 
for every $b  \in B \setminus \{O\}$ going to an element of $\iota_Z(B) \cap S_n$ 
for a constant $C> 0$.
	 
	    (II) Our first major step is to prove
	 $\mathcal{I}(B \setminus \{O\}) \subset [k+1, n-k]$: 
	  Let $\iota_n$ denote the opposite involution of $\mathfrak{a}_n$. % of $\GL(n, \bR)$. 
First, suppose that the rank of $Z$ is $\geq 2$, and hence
$\dim \mathfrak{a}_Z^+ \geq 2$. 

   \purple{Recall from the above of 
     this proposition that $B \setminus \{O\}$ is path-connected.}
	
%	Recall our map $A_Z^+ \ra A_n$ induced by 
%	$Z(\mathcal{L}(\rho(\Gamma_M)) \ra \SL_\pm(n, \bR)$. 
A {\em diagonal subspace }
\[\Delta_{i, j} \subset \mathfrak{a}_n \text{ for } i\ne j, 
i, j=1, \dots, n,\] 
is a subspace of the maximal abelian algebra
$\mathfrak{a}_n$ given by 
$\log\lambda_i - \log\lambda_j=0$ for some indices $i$ and $j$. 
	The inverse images under $\iota_Z$ in $\clo(\mathfrak{a}_Z^+)$ of the 
	diagonal subspaces of $\mathfrak{a}_n$ may meet $B$. 
%	We let $\mathcal{I}_B$ the maximal collection of indices 
%	$(i, j), i < j$, where
%	$\iota_Z^{-1}(\Delta_{i, j})$ contains $B$, which may be empty.  
%	The components of 
%	$B$ with the union of these sets for every pair of indices 
%	$i$ and $j$ removed are convex domains in $B$ of 
%	dimension $\dim B$. 
%	Any diagonal subspace not in $\mathcal{I}_B$ is called 
%	{\em extra} diagonal subspaces. 

\purple{Hueristically speaking, 
under $\iota_Z$, $B$ is mapped piecewise linearly in $\mathfrak{a}_n$ and when it 
meets $\Delta_{i, j}$ for some $i, j$, it continues in the subspace obtained by 
applying a reflection about $\Delta_{i, j}$.}  

\purple{Let $\mathcal{I}'_B$ the collection of indices $(k, l)$, $k< l$, of 
	$\Delta_{k, l}$ so that $B$ meets $\Delta_{k, l}$ transversally}. 
	We call the closures of the components of 
	\[B - \bigcup_{(k, l)\in \mathcal{I}'_B} \iota_Z^{-1}(\Delta_{k, l}) \] %- \bigcup_{i=1}^n (x_i\circ \iota)^{-1}(0) \] 
	the {\em generic flat \purple{cones} of $B$}.  
	There are finitely many of these 
	to be denoted $B_1, \dots, B_m$ which are convex 
	\purple{cones} in $B$.  
\purple{These are mapped linearly inside the positive Weyl chamber of 
$\mathfrak{a}_n$.}
	

We make notes of two ``discrete continuity" properties due to Danciger and Stecker: 
\begin{itemize} 
\item	In the interior of each $B_i$, $\mathcal{I}$ is constant since the 
$\mathcal{I}$-value is never empty and the change  in $\mathcal{I}$-values 
will happen only in the inverse images of some of $\Delta_{k,l}$. 
	
%	Also, $I_1$ is constant on each generic open domains of $B$. 
%	$\iota_+(B)$ since $\iota_+$ is a polyhedral map.
%	Let us denote by $B_1, \dots, B_m$ these polyhedral images. 
	
\item	For adjacent regions, their images under $\mathcal{I}$ differ by
	adding or removing some top or bottom consecutive sets. 
\end{itemize} 
Heuristically speaking, worms leave ``traces".
	
	We now aim to prove that the following never occurs: 
	\begin{quotation} 
	There exists a pair of elements $b_1, b_2 \in B$
	with $b_1 \in B_p^o$ and $b_2 \in B_q^o$ for some $p, q$
	where $\mathcal{I}(b_1)$ has an element $\leq k$ and 
	$\mathcal{I}(b_2)$ has an element $> k$---(*).
	\end{quotation} 
(The idea is that $k, k+1$ both cannot be in a zero set since they differ by a constant due to Anosov properties.)

We use the discrete version of continuity: 
	Suppose that this (*) is true. 
	Since $B$ is convex with $\dim B \geq 2$, 
	there is a chain of adjacent polyhedral images 
	$B_{l_1}, \dots, B_{l_m}$ where $b_1 \in B_{l_1}^o$
	and $b_2 \in B_{l_m}^o$
	where $B_{l_j}\cap B_{l_{j+1}}$ contains a nonzero point.  
	For adjacent $B_{l_p}$ and $B_{l_{p+1}}$, 
	we have by \eqref{eqn:limI} 
%	the $I$-value of the points of the intersection of the respective closures
\begin{equation} \label{eqn:intI} 
 \mathcal{I}(B_{l_p}^o)\cup \mathcal{I}(B_{l_{p+1}}^o) \subset \mathcal{I}(x)  \text{ for }
x\in \clo(B_{l_p})\cap \clo(B_{l_{p+1}}) \setminus \{O\}. 
\end{equation} 
%	must contain the both sets $\mathcal{I}(B_{l_p})$ and $\mathcal{I}(B_{l_{p+1}})$. 
	Considering $\mathcal{I}(B_{l_j}^o)$, $j=1, \dots, m$, and 
	the $\mathcal{I}$-values of the intersections of adjacent generic flat \purple{cones}, we have 
$\mathcal{I}(b') \ni k, k+1$ for some nonzero element $b'$ of $B$ by 
the connectedness of the $\mathcal{I}$-values as we change
generic flat \purple{cones} to adjacent ones according to 
\eqref{eqn:intI}.
Hence, we obtain
\begin{equation}\label{eqn:kk+1}  
\log \lambda_k(b') = 0 
\hbox{ and }\log \lambda_{k+1}(b')=0.
\end{equation}  





%	by the density of images of $\lambda_Z(\rho(\Gamma_M))$ in $B$
%	in directions by the main theorem of Benoist 
%	\cite{Benoist98} and \eqref{eqn:lklkp}. 
	% where $S_n$ is the unit sphere in
	%the Cartan algebra $A_n$.
	Since we can assume $\iota_Z(b') \in S_n$ for $b'$ in \eqref{eqn:kk+1}
	by normalizing, we obtain a contradiction to \eqref{eqn:lklkpC}. 
	

	
	Hence, we proved (*), and  
only one of the following holds:
	\begin{itemize} 
	\item for all $b\in B$ in the interiors of generic flat \purple{cones},  every element of $\mathcal{I}(b)$ is $\leq k$ 
	or 
	\item for all $b \in B$ in the interiors of generic flat \purple{cones}, every element of 
	$\mathcal{I}(b)$ is $> k$.
	\end{itemize} 

Note that the opposition involution $\iota_n$ sends $\iota_Z(B)$ to itself
since the inverse map $g\ra g^{-1}$ preserves $Z$ 
and the set of eigenvalues of $g$ is sent to the set of 
the eigenvalues of $g^{-1}$.
	In the first case, by the opposition involution $\iota_n$ preserving $G$, 
	$\mathcal{I}(\iota_n(b))$ for $b \in B$ contains an element $\geq n-k+1$.
	This again contradicts the above. 
	Hence, we conclude $\mathcal{I}(b) \subset [k+1, \dots, n]$ for all $b \in B$
	in the interiors of generic flat \purple{cones}. 
Acting by the opposition involution and \green{using similar} arguments, we obtain 
$\mathcal{I}(b)\subset [k+1, \dots, n-k]$ for all $b \in B$ in the interiors of generic flat 
\purple{cones}. 

Furthermore, we can show that $\mathcal{I}(b) \subset [k+1, n-k]$ for all $b \in B \setminus \{O\}$ since otherwise, 
we have $\mathcal{I}(b') \ni k, k+1$ for some $b'\ne O$ by a similar reason to the above applied to lower-dimensional strata. 
This is a contradiction as before. 
In addition, the argument shows $[k+1, n-k] \ne \emp$ and $k < n/2$. 


	 Now we consider the remaining case where $Z$ has rank $1$. 
Then the maximal abelian group $A_Z$ is $1$-dimensional. 
For each $g \in Z$, there is an 
index $i(g)$ for which $\log\lambda_{i(g)}(g) =0$. 
Thus, $\mathfrak{a}_Z$ is in the null space of $\log \lambda_i$ for some index $i$. 
Since $g \mapsto g^{-1}$ preserves $\mathfrak{a}_Z$, 
we have $\log \lambda_{n-i+1} = 0$ on $\mathfrak{a}_Z$. 
Since $\clo(\mathfrak{a}_Z^+) = \mathfrak{a}_Z \cap \clo(\mathfrak{a}^+)$, we also have 
$\log \lambda_{n-i+1}=0$ on $\mathfrak{a}_Z^+$. 
Hence, 
\[\mathcal{I}(B\setminus \{O\}) = [i, \dots, n-i+1] \hbox{ for some index } i, 0< i \leq n/2.\] 
Now by \eqref{eqn:lklkpC}, we must have $k < i$. This 
implies $\mathcal{I}(B \setminus \{O\}) \subset [k+1, n-k]$. 
%Also, $k< n/2$ since $[i, \dots, n-i+1] = \emptyset$ otherwise. 







	
	(III) To complete the proof, we show the 
	partial hyperbolicity property: 
%We write $\vec{a}$ 
%the Cartan projection of $\SL_\pm(n, \bR) \ra A_n^+$, 
We let $\log \vec{a}_Z|Z: Z \ra \mathfrak{a}_Z^+$ denote the Cartan projection of $Z$. 
Recall 
\[\iota_Z\circ \log \vec{a}_Z| \rho(G) 
= \log \vec{a}|\rho(G).\] 
%by choosing compact subgroups carefully. 
Since $\rho(G)$ has a connected reductive Zariski closure, 
the sequence 
\[\left\{\frac{\log \vec{a}(\rho(g_l))}{w(g_l)}\right\}\]
\purple{has} limit points only in $\iota_Z(B)$
by Theorem 1.3 of \cite{BS2021}. 
%and that of 
%\[(\log \lambda_1(\rho(g_l))), \cdots, \log \lambda_n(\rho(g_l))/w(g_l)\]
%have the same limit points in $B$.  
By  the conclusion of Step (II), % and 
%Theorem 1.3 of \cite{BS2018p}, 
for any sequence $\{g_l\}$ in  $G$, 
\[
\frac{\min_{i=k+1,\dots,n-k} |\log a_i(\rho(g_l))|}{w(g_l)} \ra 0,
\]
uniformly in terms of $w(g_l)$. 
%Here, $i$ may depend on $\{g_l\}$. 
Hence, this implies that for any small constant $C> 0$, there is $N> 0$ such that 
\[a_i(\rho(g)) \geq \exp(-C w(g)), g \in G\] for at least one $i \in [k+1, \dots, n-k]$ 
provided $w(g) > N$.
This means that 
\[a_i(\rho(g)) \geq A \exp(-C w(g))\] for at least one $i \in [k+1, \dots, n-k]$ for all $g\in G$ 
for a constant $A > 0$ depending on $C$. 
By the $P_\theta$-Anosov property of $\rho$ and Lemma \ref{lem:Equiv}, we have 
\[\frac{a_k(\rho(g))}{a_{k+1}(\rho(g))}\geq A'\exp(C'w(g)), g \in \Gamma\] 
for some constants $A', C'> 0$. 
Since 
\[a_{k+1}(\rho(g)) \geq a_i(\rho(g)), i=k+1, \cdots, n-k, 
\hbox{ for every } g \in G,\] 
it follows that 
\[a_k(\rho(g)) \geq A'' \exp(C'' w(g))\] for some constants $C'', A''> 0$. 
Taking inverses, we also showed that $\rho'$ is partially hyperbolic 
in the singular-value sense with the index $k$. 

Finally, since $[k+1, n-k]$ is not empty, $k < n/2$.  
	\end{proof} 


For a general representation $\phi: \Gamma \ra G$, we define the {\em semisimplification} 
$\phi^{ss}: \Gamma \ra G$ by post-composing $\phi$ with the projection to 
the Levi factor of the Zariski closure of $\phi(\Gamma)$.  
(See Section 2.5.4 of \cite{GGKW17ii} for details.)
Also, the Zariski closure of a finite-index group is reductive if so is the Zariski closure of the original discrete group.
(See Remark 2.38 of \cite{GGKW17ii}.) 


	
%	\marginpar{Now theorem has a converse part... do that}
	\begin{proof}[Proof of Theorem  \ref{thm:main2}] 
		\purple{Suppose that $\rho$ is $P_\theta$-Anosov in the sense of \cite{GGKW17ii}
where $k \in \theta$ and $\theta^\ast =\theta$. 
		Hence, $\Gamma:= \rho(\pi_1(N))$ is word-hyperbolic by 
the definition in \cite{GGKW17ii}.}  
	%	Proposition 5.61 of 
	%	Appendix 5.11 of \cite{KLP18} shows that 
	%	$\rho$ is $P^I$-Anosov for $I$ missing an index $k, 1\leq k \leq n/2$
	%	in the singular value sense. 
	
		First, suppose that the Zariski closure of the image of the linear holonomy group is 
	reductive. 
	By Lemma \ref{lem:Equiv}, $\rho$ is a $k$-dominated
	representation of the index $k$. 
%		 Since we can prove the corollary for any finite index subgroup of  $\Gamma$, 
We take a finite-index normal subgroup $\Gamma'$ of $\Gamma$ 
\purple{so that} the Zariski closure $Z$ of $\rho(\Gamma')$ is a connected 
	reductive Lie group. 	 
%	By Theorem \ref{thm:HKS} and 
	By Proposition \ref{prop:PAnosovStrict}, $\rho|\Gamma'$ is 
	partially hyperbolic in the singular-value sense with the index $k$, $k < n/2$.
   So is $\rho$ since $\Gamma$ is a finite-index extension of $\Gamma'$. 
	By Theorem \ref{thm:bundle}, $\rho$ is a partially hyperbolic 
	representation in the bundle sense with the index $k$.

%	For each $g \in \Gamma_M -\Gamma'$, 
%	$\mathcal{L}(g)(\bV_+) \subset \bV_+$ and 
%	$\mathcal{L}(g)(\bV_-) \subset \bV_-$ 
	
%	By Theorem \ref{thm:main}, this is a contradiction. 
	%\qed
	
	Now, we drop the reductive condition of the Zariski closure. 
Mainly, we have to prove the expanding and contracting properties in the subbundles. 
\purple{From now on, $P:= P_\theta$. The semisimplification $\rho^{ss}$ is $P$-Anosov by Proposition 1.8 of \cite{GGKW17ii},
and the Zariski closure of the image of $\rho$ is reductive. } 
%	Since conjugations do not change eigenvalues,
% by the proof of Lemma 2.40 of \cite{GGKW17ii}, 
\purple{We can conjugate $\rho$ arbitrarily close to $\rho^{ss}$ in 
$\Hom(\Gamma, \GL(n, \bR))$
since the orbit of $\rho^{ss}$ is the closed orbit in the closure of the orbit of $\rho$ under the conjugation 
action. (See Section 2.5.4 of \cite{GGKW17ii}. )}
\purple{We can assume that $\rho$ has the property (i) by Lemma 2.11 of \cite{BS2021}. 
Since eigenvalues do not change under conjugation, 
$\rho^{ss}(g)$	for each $g \in \Gamma$ has an eigenvalue equal to $1$. 
Hence (i) holds for $\rho^{ss}$. }
%\purple{Suppose that $\rho$ has the property (ii).  
%Each point of the joint spectrum $J(S)$ of $\rho$ has a coordinate equal to zero
%by Theorem 1.1 of \cite{BS2021}. 
%By Theorem 1.2 of \cite{BS2021}, each $\rho(g)$ has an eigenvalue with norm $1$,
%and hence (i) holds for $\rho^{ss}$ as well.}
% by	Theorem \ref{thm:HKS}. 
(See also the proof of Lemma 2.40 of \cite{GGKW17ii}.)
	By Proposition \ref{prop:PAnosovStrict} and Theorem \ref{thm:bundle}, 
	$\rho^{ss}$ is a partially hyperbolic representation in the bundle sense with the index $k$, $k < n/2$. 
%	By Theorem 5.13 of \cite{GW12}, 
	By Proposition 1.3 of \cite{GGKW17ii}, $\rho$ is $P$-Anosov with respect to 
	a continuous Anosov map $\zeta: \partial_\infty \Gamma \ra {\mathcal{F}}$ for the flag variety 
	$\mathcal{F} := \GL(n, \bR)/P$.

Let $\zeta_i: \partial_\infty \Gamma \ra {\mathcal{F}}$ denote the Anosov map for $\rho_i$ 
conjugate to $\rho$ converging to $\rho^{ss}$. 
Theorem 5.13 of \cite{GW12} generalizes to $\GL(n, \bR)$ 
since we can multiply a representation 
by a homomorphism $\Gamma \ra \bR^+$ not changing the $P$-Anosov property
to $\SL_\pm(n, \bR)$-representations. 
(See Section 2.5.3 of \cite{GGKW17ii}). 
 %$\rho_i$ is $P$-Anosov for sufficiently large $i$, and 
Hence, Theorem 5.13 of \cite{GW12} implies that
$\zeta_i$ converges to $\zeta^{ss}: \partial_\infty \Gamma \ra {\mathcal{F}}$
the map for $\rho^{ss}$ in the $C^0$-topology. 

Let $\bV^{ss+}, \bV^{ss0}$, and $\bV^{ss-}$ denote the bundles of the partially hyperbolic decomposition of $\bR^n_{\rho^{ss}}$ 
as in Definition \ref{defn:phyp}. 
Let $\bV^{+}_i, \bV^{0}_i$, and $\bV^{-}_i$ denote the bundles of the decomposition of $\bR^n_{\rho_i}$ 
of respective dimensions $k$, $n-2k$, and $k$ obtainable from $\zeta_i$
since $\rho_i$ is $k$-dominated.  
The  single product bundle $\Uc \hat M \times \bR^n$ covers the direct sums of these bundles.
\begin{itemize} 
\item Let $\hat \bV^{ss+}, \hat \bV^{ss0}$, and $\hat \bV^{ss-}$ denote the cover of 
$\bV^{ss+}, \bV^{ss0}$, and $\bV^{ss-}$ \purple{respectively}. 
\item Let $\hat \bV^{+}_i, \hat \bV^{0}_i$, and $\hat \bV^{-}_i$ denote the cover of 
$\bV^{+}_i, \bV^{0}_i$, and $\bV^{-}_i$ \purple{respectively}.  
\end{itemize} 
Then the above convergence $\zeta_i \ra \zeta^{ss}$ means that 
\begin{equation} \label{eqn:convB}
\hat \bV^{+}_i(x) \ra \hat \bV^{ss+}(x), \hat \bV^{0}_i(x) \ra \hat \bV^{ss0}(x), 
\text{ and } \hat \bV^{-}_i(x) \ra \hat \bV^{ss-}(x) \text{ as } i \ra \infty
\end{equation}
in the respective Grassmannian spaces pointwise for each $x \in \Uc \hat{M}$. 
 
The convergence is also uniform over every compact subset of $\Uc\hat{M}$
since these vector spaces depend only on the endpoints of quasi-geodesics through $x$
and $\hat M$ is quasi-isometric to $\partial^{(2)}_\infty \Gamma_M \times \bR$. 

%Over $\Uc \hat M$, 
%The fiberwise normss $\llrrV{\cdot}_{\bR^n_{\rho_i}}$ on  $\bR^n_{\rho_i}$ and 
%$\llrrV{\cdot}_{\bR^n_{\rho^{ss}}}$ on $\bR^n_{\rho^{ss}}$ respectively  lift to 
%the fiberwise norms on the common cover $\Uc \hat M \times \bR^n$, which we will denote 
%by the same notations. 
We may construct a sequence of the fiberwise metrics $\llrrV{\cdot}_{\bR^n_{\rho_i}}$  so that 
the sequence of the lifted norms of $\llrrV{\cdot}_{\bR^n_{\rho_i}}$  
uniformly converging to that of  $\llrrV{\cdot}_{\bR^n_{\rho^{ss}}}$ on $\Uc \hat M \times \bR^n$
over each compact subset of $\Uc \hat M$ 
using a fixed set of a partition of unity subordinate to a covering with trivializations.  

The flow $\phi^{ss}$ on $\bR_{\rho^{ss}}$ satisfies the expansion and contraction properties
in Definition \ref{defn:phyp}. 
Note that the lift $\Phi^{ss}$ of $\phi^{ss}$ and that $\Phi_i$ of the flow $\phi_i$ for $\bR^n_{\rho_i}$ in the product bundle 
$\Uc \hat M \times \bR^n$ are the same by construction
induced from trivial action on the second factor. 
Since $M$ is compact, there exists $t_0$ such that for $t> t_0$, 
\begin{itemize} 
\item $\Phi^{ss}_{t}| \bV^{ss+} $  is expanding by a factor $c^+> 1$ 
\item $\Phi^{ss}_{t}| \bV^{ss-}$ is contracting by a factor $c^- < 1$  with respect to 
the metric $\llrrV{\cdot}_{\bR^n_{\rho^{ss}}}$. 
\end{itemize} 
Hence, for sufficiently large $i$, 
the flow $\Phi_i$ on $\bR^n_{\rho_i}$  satisfies the corresponding properties
as well as the domination properties for $\llrrV{\cdot}_{\bR^n_{\rho_i}}$ 
by \eqref{eqn:convB}. 
%$\rho_i$ is 
%$P_i$-Anosov with index $k$ for $P_i$ conjugate to $P$. 
Hence, $\Phi_i$ is partially hyperbolic with the index $k$.  
%		Again, by Theorem \ref{thm:main}, this is a contradiction. %\qed
		
		
\purple{Conversely}, suppose that $\rho$ is partially hyperbolic with the index $k$. 
Proposition \purple{4.6} of \cite{BPS} shows that $\rho$ is $k$-dominated. 
%\purple{(See also the second part of Lemma 4.4 of \cite{BPS}.)}
Lemma \ref{lem:Equiv} \purple{(or Proposition 4.9 of \cite{BPS})} shows that $\rho$ is $P_\theta$-Anosov for $k \in \theta$. 
Also, $k < n/2$ since the neutral bundle exists. 
\end{proof} 

%\begin{proof}[Proof of Corollary \ref{cor:cpt}]
%	By Theorem \ref{thm:main2}, $\rho$ is partially hyperbolic.  
%	There is a constant $C$ of parallelism obtained in Theorem \ref{thm:neutral}. 
%	Independent of $z \in \partial_\infty \pi_1(N)$, $\mathcal{R}_z$ embeds quasi-isometrically into 
%	a generalized stable affine subspace $A_z$ by lifting points to $\Uc\tilde M$ and then applying $\tilde s_\infty$
%	moving only a distance bounded above by $C$ by Proposition \ref{prop:quasiiso}.
%	Let $K$ be a compact fundamental domain of $\Uc \tilde M$. 
%	Since $\mathcal{R}_z$ is $C'$-dense for some $C'> 0$, 
%	a point of $\mathcal{R}_z$ is in a $C'$-neighborhood $K'$ of $K$. 
%	Since $A_z$ contains $\tilde s_{\infty}(K' \cap \mathcal{R}_z)$,
%	the collection is bounded in $\mathcal{AG}_{n-k}(\bR^n)$. 
%	Taking the closure, we obtain a compact subset. 
%	Since all $A_z$ for $z\in \partial_\infty \pi_1(M)$ are considered, 
%	the invariance under the affine group follows. 
%\end{proof} 

%\begin{remark}
%	Provided $M$ is closed, 
%we may have assumed in the above proof that  
%the linear part homomorphism $\rho:\Gamma \ra \GL(n, \bR)$ 
%is injective by Corollary 1.1 of Bucher-Connel-Lafont \cite{BCL}
%since $\Gamma$ is hyperbolic and hence the simplicial volume is nonzero
%by Gromov \cite{Gromov82}. 
%\end{remark}

%\begin{remark}
%Finally, the so-called Markus conjecture, probably due to Goldman and Thurston, 
%is that a closed affine $n$-manifold is complete if and only if 
%its linear part is in $\SL_\pm(n, \bR)$. If the forward direction
%is known to be true, then we could just work with 
%$\SL_\pm(n, \bR)$ here. 
%\end{remark} 

%\begin{lemma}\label{lem:change}  
%	Suppose that we change a Euclidean metric on $\bR^n$
%	to a new one $\llrrV{\cdot}_2$. 
%	Then the singular value $a_i'(g), g \in \SL_\pm(n, \bR)$ 
%	satisfies 
%	\[  \frac{1}{C} a_i(g) \leq a_i'(g) \leq C a_i(g)  \]
%	for a constant $C> 1$ depending only on $\llrrV{\cdot}$ and 
%	$\llrrV{\cdot}$. 
%	\end{lemma} 
%\begin{proof} 
%We can use the min-max theorem for
%singular values restricted to a fixed sphere $\SI^{n-1}$ where
%\[ a_k(g) = \min_{S, \dim S = k} \max_{v\in S}  \frac{\llrrV{gv}}{\llrrV{v}}.\]
%We now change the metric to $\llrrV{\cdot}_2$ 
%and every ratio in the set we are considering are 
%comparable by a constant $C> 0$ independent of $g\in \SL_\pm(n, \bR)$
%since $\SI^{n-1}$ and the Grassmannians are compact. 
%\end{proof}

%\begin{proof}[Proof of Corollary \ref{cor:main2}] 
%
%	
%	Suppose that the linear part $\rho$ of
%	the affine holonomy homomorphism $\rho'$ is $P$-Anosov
%	for a parabolic subgroup $P$ of $\SL_\pm(n, \bR)$.
%%	where $1 \leq k < n/2$. 
%	If $k=0$ or $k=n$, then our affine manifold is Euclidean 
%and the fundamental group of a closed Euclidean manifold is of polynomial growth by the Bieberbach theorem
%and cannot be hyperbolic. 
%Hence, we may assume $1 \leq k \leq n/2$. 
%	
%	Assume $k < n/2$. 
%	From Theorems 5 and 9 of \cite{Sawyer}, 
%	we obtain that singular values occur in pairs %there is an integer $p$ such that 
%	\begin{multline} \label{eqn:ai} 
%		a_i(g), a_{n-i+1}(g) = \frac{1}{a_i(g)} \hbox{ for }i=1, \dots, k,
%		\hbox{ and } \\
%		a_j(g) = 1 \hbox{ for every } j \in [k+1, n-k], g \in \SO(k, n-k)
%	\end{multline} 
%	under an orthogonal change of the coordinate system. (See also Section 7.2 of 
%	Borel \cite{BorelN}.)
%	
%	Since the $P$-Anosov property holds for 
%	$P = P_{\tau_{mod}}$ for some model simplex $\tau_{mod}$ 
%	contained in $x_j = x_l=0$ for every $j, l\in [k+1, n-k]$ 
%	in the sense of Kapovich-Leeb-Porti \cite{KLP17}, 
%	$\rho$ is $p$-Anosov for some $p$ for $p \leq k$. 
%	Hence, $\rho$ is $p$-dominated
%	by Proposition 4.5 of \cite{BPS}.
%	Since $k < n/2$, the above property tells us that 
%	for $i\in [k+1, n-k]$ so that $a_{i}(\rho(g)) = 1$. 
%	%For other coordinates, $a_{i(g)}$ and $a_{j(g)}$ are bounded 
%	%above and below by constants by Lemma \ref{lem:change}. 
%	%Also, $a_{k+1}(\rho(g)) \geq 1$ and $a_{n-k}(\rho(g)) \leq 1$. 
%	%for $i \in [p+1, n-p]$. 
%	%Hence, the neutrality conditions are satisfied. 
%	%Corollary \ref{cor:main} implies the result. 
%	%As in the proof of 
%	Since $a_{p+1}(\rho(g))\geq 1$, we obtain 
%	$a_p(\rho(g)) \geq C \exp(w(g))$ by the p-domination property. 
%	Hence, Lemma \ref{lem:Equiv}, Theorem \ref{thm:bundle}, and Theorem \ref{thm:main} imply the %result.
% 
%	 
%	Finally, we consider $\SO(k, k)$ for $n= 2k$. This is the last remaining case. 
%	As in the above paragraph, $\rho$ is $p$-dominated for $p \leq n/2$. 
%	Suppose that $p < n/2$. Then $a_{p+1}(\rho(g)) \geq 1$ and $a_{n-p-1}(\rho(g)) \leq 1$ 
%	for every $g \in h(\Gamma_M)$
%	by \eqref{eqn:ai}. Then the result follows as above. 
%%	Lemma \ref{lem:Equiv} and Theorem \ref{thm:bundle} prove the conclusion
%	for these two cases. 
%	
%	%The neutrality conditions are satisfied for index $p$. 
%	%Corollary \ref{cor:main} implies the result. 
%	
%	Suppose $\rho$ is $k$-dominated for $k= n/2$. 
%	Then 
%	\begin{multline}  
%	\rho(g) = k_1(\rho(g)) t(\rho(g)) k_2(\rho(g)) \hbox{ for } g \in\Gamma_M\\
%\hbox{ where } 
%	k_i(\rho(g)) \in \SI(\Ort(k, \bR)\times \Ort(k, \bR))
%	\end{multline} 
%	holds,  and 
%	$t(\rho(g))$ a positive diagonal matrix with diagonal entries 
%	\begin{equation} \label{eqn:ak} 
%	a_1(\rho(g)), \dots, a_k(\rho(g)), \frac{1}{a_k(\rho(g))}, \dots, \frac{1}{a_1(\rho(g))}
%	\end{equation} 
%	under the coordinate system of Section 7.2 of Borel \cite{BorelN}.
%	We have 
%	\begin{equation} \label{eqn:ak2} 
%	a_k(\rho(g))^2 \geq C\exp(a w(g)) \hbox{ for }g\in \Gamma_M
%	\hbox{ for fixed } a> 0, C> 0.
%\end{equation}
%%	This implies that for the subspace 
%%	corresponding to $1, \dots, k$, 
%%	$\rho(g)$ is strictly expanding and for the subspace 
%%	corresponding to $k+1, \dots, n$, 
%%	$\rho(g)$ is strictly contracting. 
%	By \eqref{eqn:ak2} and 
%%	Theorem 1.2 of \cite{BS2018p} applied to a single element set
%%	$\{g\}$, 
%%the last corollary in Section 2 of Benoist \cite{Benoist97}
%and Lemma 2.11 of \cite{BS2018p}, 
%\[\log \vec{\lambda}(\rho(g) = \lim_{m \ra \infty} \frac{1}{m} \log(\vec{a}(\rho(g)^m)),\] 
%	we obtain that no eigenvalue equals $1$
%%	Hence, this means that $\rho(g)$ fixes a point in $\bR^n$
%%	for sufficiently large $w(g)$ by 
%	contradicting Theorem \ref{thm:HKS}. 
%	%for $g = \rho(q)$ for $q$ in $\Gamma_M$ with sufficiently large $|q|$. 
%%	This contradicts the free action property of $\rho$. 
%	
%	For symplectic group $\mathsf{SP}(2n)$, there is a Bloch-Messiah decomposition 
%	$\mathsf{SP}(2n) = U(2n)DU(2n)$ where
%	$U(2n) = \Ort(2n, \bR)\cap \mathsf{SP}(2n)$, 
%	$D$ is the group of diagonal matrix 
%	with entries $\lambda_1, \dots, \lambda_n, \lambda_n^{-1}, \dots, \lambda_1^{-1}$ in a decreasing order. Hence, the similar argument as 
%	$\SO(k, k)$ will apply. 
%	(See Tables 8 and 9 of Reference Chapter of \cite{OV90}.)
%	%\qed
%%	
%	%	 \marginpar{ still mess}
%%	
%	%	This would be satisfied if the linear part of $\rho$ is conjugate 
%	%	into $\SO(k, n-k)$: From equations (2.7) and (2, 9)
%	%	of Cornwell \cite{Cornwell75ii}, we see the 
%	%	basis of the maximal algebras are diagonal with entries $\pm 1$ 
%	%	whose traces are $0$. Hence, the number of $\log a_i$ greater than $0$ 
%	%	equal the number of $\log a_i$ less than $0$. 
%	%	Hence, we can let $k$ to be the maximum of the number of 
%	%	$a_i(\rho(g))$s with $\log a_i(\rho(g)) > 0$ for all $g$ in $\rho(\Gamma_M)$. 
%	%	Then $n-k+1$ is the minimum
%	%	Hence, neutral index is always in $[k+1, n-k]$. 
%%	
%	%\marginpar{Still unclear what to do}
%\end{proof} 


%\bibliography{AAII}{}
%
%\bibliographystyle{plain}
%
% BibTeX users please use one of
%\bibliographystyle{spbasic}      % basic style, author-year citations
%\bibliographystyle{spmpsci}      % mathematics and physical sciences
%\bibliographystyle{spphys}       % APS-like style for physics
%\bibliography{AAIII}   % name your BibTeX data base


\section{Complete affine manifolds} \label{sec:affine} 



%\subsection{$(G,X)$-structures and affine manifolds}
Let $G$ be a Lie group acting transitively and faithfully on 
a space $X$. 
A {\em $(G, X)$-structure} on a manifold $N$ is a maximal atlas of charts to $X$ so that 
the transition maps are in $G$. 
This is equivalent to
$N$ having a pair $(\dev, h)$ where 
$\hat N$  is a regular cover of $N$ with a deck transformation group $\Gamma_N$
such that
\begin{itemize} 
	\item $h:\Gamma_N \ra G$ is a homomorphism, called a {\em holonomy homomorphism}, and 
	\item $\dev: \hat N \ra X$ is an immersion,  
	called a {\em developing map}, satisfying 
	\[\dev \circ \gamma = h(\gamma)\circ \dev \]
 for each deck transformation  $\gamma \in \Gamma_N$.
\end{itemize}

%We recall the way to see this from a bundle sense: 
%Construct $X_h = \hat M \times X/\Gamma_M$ 
%where the action is given by 
%\[g(x, y) = (g(x), h(g)(y)), x\in \hat M, y\in X, g\in \Gamma_M.\]  
%This is a fiber bundle over $M$ with fibers $X$. 
%$X_h$ is a bundle over $M$
%with a flat connection induced from 
%the product structure.
%There is a {\em developing} section 
%\[s: M \ra X_h \hbox{ given 
%by } \hat M \ni x  \mapsto (x, \dev(x)) \in 
%\hat M \times X.\] 
%The section is transverse to the leaves of flat connection.
%Conversely, a transverse section 
%$s: M \ra X_h$ gives us a $(G,X)$-structure. 
%(See \cite{G88}.)

We say that an $m$-manifold $N$ 
with a $(G, X)$-structure  is {\em complete} if $\dev: \hat N \ra X$ is a diffeomorphism. 
If $N$ is complete, 
we have a diffeomorphism 
$N \ra X/h(\Gamma_N)$, and $h(\Gamma_N)$ acts properly discontinuously and freely on $X$. Complete $(G, X)$-structures on $N$ are classified 
by conjugacy classes of representations $\Gamma_N \ra G$ with properly discontinuous and free actions on $X$. 
(See Section 3.4 of \cite{Thurston97}.)

Note that the completeness and compactness of $N$ have 
no relation for some $(G,X)$-structures 
unless $G$ is in the isometry group of $X$ with a Riemannian metric.
(The Hopf-Rinow lemma may fail.)

Now, we move on to the affine structures. 
Let $\mathds{A}^n$ be a complete affine space. 
	Let $\Aff(\mathds{A}^n)$ denote the group of affine transformations of $\mathds{A}^n$
	whose elements are of the form: 
	\[  x \mapsto A x + \bv   \]
	for a vector $\bv \in \bR^{n}$ and $A \in \GL(n, \bR)$. 
	Let $\mathcal{L}: \Aff(\mathds{A}^n) \ra \GL(n, \bR)$ denote the map that sends each element of $\Aff(\mathds{A}^n)$ to its linear part in $\GL(n, \bR)$.
	
	
	
	
		An affine $n$-manifold is an $n$-manifold with  an $(\Aff(\mathds{A}^n), \mathds{A}^n)$-structure. 
%An affine $n$-manifold is {\em special} if $\mathcal{L}(\Gamma)\subset \SL_{\pm}(n, \bR)$. 
		A {\em complete affine $n$-manifold}
			is an $n$-manifold $N$ of the form $\mathds{A}^n/\Gamma$ for $\Gamma \subset \Aff(\mathds{A}^n)$
			acting properly discontinuously and freely on $\mathds{A}^n$.  
			
			
	\begin{theorem}[Hirsch-Kostant-Sullivan \cite{KS75}]\label{thm:HKS}
	Let $N$ be a complete affine $n$-manifold.
	Let $\rho:\pi_1(N) \ra \GL(n, \bR)$ be the linear part of the affine holonomy representation $\rho'$.
	Then $g$ has an eigenvalue equal to $1$ for each $g$ in the Zariski closure $Z(\rho(\pi_1(N)))$ of $\rho(\pi_1(N))$. 
\end{theorem}
	\begin{proof} 
		We define the function $f: Z(\rho(\pi_1(N))) \ra \bR$ 
		by $f(x) = \det(x-\Idd)$ for each element $x$. 
	%	First, we show this for $\rho(g) \in \rho(\Gamma_M) -\{\Idd\}$. 
		Suppose that $f(\rho(g)) \ne 0$, $g\in \Gamma_M \setminus \{\Idd\}$. Then $\rho(g)(x) + b = x$ has 
		a solution for each $b$ by Cramer's rule. 
		Hence, $\rho'(g)$ has a fixed point, which is a contradiction.
		We find that $f$ is zero in the Zariski closure, 
        which implies that each element has at least one eigenvalue equal to $1$.  
        %\qed
		\end{proof} 



\begin{proof}[Proof of Corollary \ref{cor:affine}] 
By Theorem \ref{thm:HKS}, the main premise of Theorem \ref{thm:main2} is satisfied. 
Hence, the result follows. 
\end{proof} 
	
	


\bibliographystyle{plain}
\bibliography{AA-PI}   

\end{document}

